\documentclass[11pt]{article}

% ── Packages ────────────────────────────────────────────────────────────────
\usepackage[margin=1in]{geometry}
\usepackage{amsmath,amssymb}
\usepackage{booktabs}
\usepackage{enumitem}
\usepackage{titlesec}
\usepackage{fancyhdr}
\usepackage{xcolor}
\usepackage{hyperref}
\usepackage{listings}
\usepackage[protrusion=true,expansion=false]{microtype}
\usepackage{url}
\urlstyle{tt}
\usepackage{seqsplit}
\usepackage{parskip}
\usepackage{array}
\usepackage{longtable}
\usepackage{mdframed}
\usepackage{tabularx}
\usepackage{multirow}

% ── Colors ───────────────────────────────────────────────────────────────────
\definecolor{sectionblue}{RGB}{30,80,140}
\definecolor{darkblue}{RGB}{20,60,110}
\definecolor{lightgray}{RGB}{245,245,245}
\definecolor{taskgray}{RGB}{230,235,242}
\definecolor{codegray}{RGB}{100,100,100}
\definecolor{done}{RGB}{220,240,220}
\definecolor{warn}{RGB}{255,243,220}

% ── Hyperref ─────────────────────────────────────────────────────────────────
\hypersetup{colorlinks=true, linkcolor=sectionblue, urlcolor=sectionblue}

% ── Section formatting ───────────────────────────────────────────────────────
\titleformat{\section}
  {\large\bfseries\color{sectionblue}}
  {\thesection}{1em}{}[\vspace{-4pt}\color{sectionblue}\rule{\linewidth}{0.6pt}]

\titleformat{\subsection}
  {\normalsize\bfseries\color{darkblue}}
  {\thesubsection}{1em}{}

\titleformat{\subsubsection}
  {\normalsize\bfseries}
  {\thesubsubsection}{1em}{}

% ── Code style ───────────────────────────────────────────────────────────────
\lstdefinestyle{pycode}{
  language=Python,
  basicstyle=\ttfamily\footnotesize,
  keywordstyle=\bfseries\color{sectionblue},
  commentstyle=\color{codegray},
  stringstyle=\color{darkblue},
  backgroundcolor=\color{lightgray},
  frame=single, framesep=4pt,
  breaklines=true,
  showstringspaces=false,
  tabsize=4
}
\lstset{style=pycode}

% ── Task box ─────────────────────────────────────────────────────────────────
\newmdenv[
  backgroundcolor=taskgray,
  linecolor=sectionblue,
  linewidth=0.8pt,
  innerleftmargin=10pt,
  innerrightmargin=10pt,
  innertopmargin=8pt,
  innerbottommargin=8pt,
  skipabove=6pt,
  skipbelow=6pt
]{taskbox}

% ── Global line-breaking tolerance ───────────────────────────────────────────
\tolerance=2000
\emergencystretch=4em
\hbadness=10000
\hfuzz=22pt

% ── Header / Footer ──────────────────────────────────────────────────────────
\pagestyle{fancy}
\fancyhf{}
\lhead{\small\textit{Electoral Equilibrium --- Task Instructions}}
\rhead{\small\textit{Juhi Damley, May 2026}}
\cfoot{\small\thepage}
\renewcommand{\headrulewidth}{0.4pt}

% ── Macros ───────────────────────────────────────────────────────────────────
\newcommand{\hrs}[1]{\textbf{[#1 hrs]}}
\newcommand{\veq}{V_{\mathrm{eq}}}
\newcommand{\boldmu}{\boldsymbol{\mu}}
\newcommand{\boldSigma}{\boldsymbol{\Sigma}}
\newcommand{\dayhead}[2]{%
  \subsection*{Day #1 \normalfont\small\color{codegray}(#2)}}
\newcommand{\task}[2]{%
  \noindent\textbf{Task:} #1 \hfill \hrs{#2}\\}
\newcommand{\done}{\textcolor{sectionblue}{$\square$}\ }

% ── Document ─────────────────────────────────────────────────────────────────
\begin{document}

% ── Title ────────────────────────────────────────────────────────────────────
\begin{center}
  {\LARGE\bfseries Electoral Equilibrium}\\[4pt]
  {\large Task Instructions \& Hour Ledger}\\[6pt]
  {\normalsize 320 Hours $\cdot$ 8 Weeks $\cdot$ 40 hrs/week}\\[8pt]
  \rule{0.6\linewidth}{0.8pt}\\[6pt]
  {\normalsize
    \textbf{Research Assistant:} Juhi Damley\quad
    \textbf{Supervising Mentor:} Prof.\ Gaston Espinosa\\[2pt]
    Claremont McKenna College\quad\textit{Summer 2026}
  }
\end{center}

\vspace{6pt}

% ── Executive Summary ────────────────────────────────────────────────────────
\begin{mdframed}[backgroundcolor=taskgray, linecolor=sectionblue,
  linewidth=1pt, innerleftmargin=14pt, innerrightmargin=14pt,
  innertopmargin=10pt, innerbottommargin=10pt]
\textbf{\large The Big Idea}

\medskip
Electoral forecasting is static. Polling aggregators tell you where the race stands
today but cannot answer: \textit{``If \textbf{X} happens tomorrow, which voter groups
must shift for the Democrat to still win?''}

\medskip
\textbf{Electoral Equilibrium} answers in real time. Select a party (Democrat or
Republican), type any hypothetical event --- a scandal, an assassination, an October
Surprise --- and the app immediately shows how that party's coalition of ten key voter
groups (African Americans, Latinos, Asians, Seculars, Jews, Muslims, Women, Catholics,
Protestants, Evangelicals) must rebalance, and the probability of achieving it.

\medskip
Built with a fine-tuned Mistral 7B LLM, a mean-variance stochastic optimizer, and a
two-stage NLP pipeline drawing on social media bio inference --- grounded in two
decades of Prof.\ Espinosa's research on religion, race, and presidential voting.
\end{mdframed}

\vspace{8pt}

% ── How to use ───────────────────────────────────────────────────────────────
\section*{How to Use This Document}

Each task entry shows: a concrete deliverable, the files it touches, and its hour
budget. Hours are set conservatively --- if a task runs short, bank the time and
use it as flex on the next heavy day. Every week closes with a \textbf{pull request}
to \texttt{main} and a \textbf{brief weekly supervision check-in} with Prof.\ Espinosa.

\begin{center}
\begin{tabularx}{\linewidth}{p{0.16\linewidth}p{0.40\linewidth}X}
\toprule
\textbf{Symbol} & \textbf{Meaning} & \textbf{Note} \\
\midrule
\hrs{N}   & Hour budget for this task  & Running total per day = 8 hrs \\
\done     & Checkbox to mark complete  & Print and tick off \\
HPC       & Run on USC Laguna cluster (compute); CMC HPC for hosting only & Submit batch job, don't block locally \\
M5        & Run locally on MacBook Pro & Fast iteration / experimentation \\
PR        & Pull request to main       & Required before week closes \\
\bottomrule
\end{tabularx}
\end{center}

% ── Hour summary ─────────────────────────────────────────────────────────────
\section*{Hour Summary by Week}

\begin{center}
\begin{tabular}{llr}
\toprule
\textbf{Week} & \textbf{Focus} & \textbf{Hours} \\
\midrule
1 & Infrastructure \& Artifact Contracts   & 40 \\
2 & Historical Data Pipeline               & 40 \\
3 & Baseline Voter Portfolio               & 40 \\
4 & RoBERTa Pipeline \& Fine-Tune Dataset  & 40 \\
5 & LLM Fine-Tuning \& Inference Endpoint  & 40 \\
6 & Stochastic Optimization \& Monte Carlo & 40 \\
7 & Web Application                        & 40 \\
8 & Hardening, Validation \& Deliverables  & 40 \\
\midrule
  & \textbf{Total}                         & \textbf{320} \\
\bottomrule
\end{tabular}
\end{center}

\newpage

% ════════════════════════════════════════════════════════════════════════════
\section*{AI Research Assistant Setup (Pre-Week --- 4 hrs)}
% ════════════════════════════════════════════════════════════════════════════

\textbf{Do this before the first day of work.} The goal is a 24/7 AI research
assistant running on the Raspberry Pi 5 that you can message from Telegram or
WhatsApp at any time. Routine queries (paper summaries, quick code questions,
daily briefings) route to \textbf{Gemini Pro} for free. Deep architectural work
routes to \textbf{Claude} on demand.

\begin{center}
\begin{tabularx}{\linewidth}{p{0.34\linewidth}Xp{0.22\linewidth}}
\toprule
\textbf{Component} & \textbf{Role} & \textbf{Cost} \\
\midrule
OpenClaw on Raspberry Pi 5 & Gateway, memory, scheduling & Free (open-source) \\
Google Gemini Pro (AI Studio API) & Lightweight daily queries & Free tier / Pro sub \\
Anthropic Claude API & Deep reasoning and code & Pay per token \\
Telegram or WhatsApp & Chat interface & Free \\
\bottomrule
\end{tabularx}
\end{center}

% ── Day 0 ────────────────────────────────────────────────────────────────────
\dayhead{0}{Sun --- Environment Setup: OpenClaw, Syncthing \& Pi NPU (pre-work, 6 hrs)}

\begin{taskbox}
\done \task{Apply for student programs and free tiers \textbf{before anything else}.
  (i) \textbf{GitHub Student Developer Pack} at \url{education.github.com/pack}
  --- unlocks JetBrains IDEs, Namecheap domain (.me free for 1 year),
  Canva Pro, and more. Takes minutes.
  (ii) \textbf{Social media access --- entirely free, no subscription needed.}
  Sign up for \textbf{Bluesky} at \url{bsky.app} (primary live collection,
  AT Protocol firehose, no cap) and \textbf{Apify} at \url{apify.com}
  (secondary live collection, free tier 500 results/run, covers the X
  demographic profile without any API cost). The bulk of historical Twitter
  data comes from free public archives downloaded in Week~4: 56M Kavanaugh
  tweets, 2020 election datasets, MeToo tweets, and more.
  \textbf{Do not sign up for the X Basic API (\$100/month) --- the archives
  are superior for historical shocks and Bluesky/Apify covers live
  collection at zero cost.}
  (iii) \textbf{Meta Content Library} at
  \url{developers.facebook.com/docs/content-library} --- free academic access;
  takes 2--4 weeks. Apply today.
  (iv) \textbf{Modal account} at \url{modal.com} --- \$30/month free credit;
  at Mistral 7B on an A100 (approximately \$0.0003 per call), covers around
  100{,}000 inference calls for free. Sign up immediately.}{1}

\done \task{SSH into the Raspberry Pi 5. Install OpenClaw:
  \texttt{curl -fsSL https://openclaw.ai/install.sh | bash}. Run
  \texttt{openclaw onboard} and follow the wizard. Select Telegram or
  WhatsApp as the primary channel. Configure Gemini as default model and
  Claude Sonnet as the \texttt{@claude} route. Store both API keys in
  \texttt{/etc/environment}. Docker-containerise OpenClaw for security.}{1}

\done \task{Install \textbf{Syncthing} on the four local machines only
  (M5, Intel Mac, Pi). \textbf{Do not install Syncthing on USC Laguna or CMC HPC}
  --- HPC nodes kill background daemons and a Syncthing daemon will be
  terminated silently, starving any SLURM jobs that depend on it. Create a
  shared \texttt{electoral-sync/} folder synced across the four local devices.
  Subdirectories: \texttt{rawdata/} (collected posts and articles),
  \texttt{artifacts/} (pipeline outputs), \texttt{data/finetune/} (training
  data). Test by writing a file on the Intel Mac and confirming it appears
  on the M5 within 30 seconds. HPC data transfer uses explicit
  \texttt{rsync}/\texttt{scp} before each job submission.}{1}

\done \task{\textbf{Compile the Hailo HEF now in Week 0 --- not Week 4.}
  NPU compilation to custom silicon routinely causes driver mismatches,
  ONNX graph translation failures, and quantization errors. This is a
  multi-day risk if hit in Week 4. Eliminate it now.
  (i) Install the Hailo Dataflow Compiler and SDK:
  \texttt{pip install hailo-sdk-client hailo-model-zoo onnx sentence-transformers}.
  Verify: \texttt{hailortcli fw-control identify} should return device info.
  (ii) Export \texttt{all-MiniLM-L6-v2} to ONNX.
  (iii) Run \texttt{hailo optimize minilm.onnx --hw-arch hailo8l}.
  (iv) Run \texttt{hailo compile --hw-arch hailo8l -o models/minilm.hef}.
  (v) Verify: \texttt{from hailo\_platform import HEF; HEF("models/minilm.hef")}.
  \textbf{CPU fallback (mandatory):} if compilation fails for any reason,
  set \texttt{"pi\_npu\_enabled": false} in \texttt{configs/base.json}.
  The bio classifier runs on Pi CPU at $\sim$20ms/bio instead of sub-ms ---
  slower but correct. The rest of the pipeline is unchanged.
  Week 4 must not be blocked on NPU compilation under any circumstances.}{3}

\done \task{Set up the \textbf{Intel Mac} as the social media collection machine.
  Configure both free collectors:
  \begin{itemize}[leftmargin=1.6em, itemsep=1pt]
    \item \textbf{Bluesky:} install \texttt{atproto} (\texttt{pip install atproto}).
          Test \texttt{app.bsky.feed.searchPosts} with the Ayatollah keywords.
          Confirm author \texttt{description} and \texttt{langs} fields are returned.
    \item \textbf{Apify:} install \texttt{apify-client} (\texttt{pip install
          apify-client}). Set the API token in \texttt{.env}. Test the X Scraper
          actor with the Ayatollah keywords. Confirm output matches the same
          \texttt{posts.jsonl} schema as \texttt{BlueSkyCollector}.
  \end{itemize}
  Set up a \textbf{launchd plist} so both collectors restart on reboot.
  \texttt{BlueSkyCollector} runs continuously (no cap).
  \texttt{ApifyCollector} is triggered per-shock (500 results/run free tier).
  Both write to \texttt{rawdata/social/\{platform\}/\{shock\_id\}/posts.jsonl}.
  \textbf{Do not set up the X Basic API --- it is not used in this project.}}{2}

\done \task{Set up the \textbf{Intel Mac} for news scraping (previously Windows role):
  Python 3.11, \texttt{jupyter}, \texttt{pandas}, \texttt{matplotlib},
  install Python 3.11, project dependencies, and Syncthing. Configure Syncthing to sync
  the shared folder. Confirm Jupyter can read files from
  \texttt{electoral-sync/rawdata/}.}{1}
\end{taskbox}

\begin{mdframed}[backgroundcolor=warn, linecolor=darkblue, linewidth=0.8pt,
  innerleftmargin=10pt, innerrightmargin=10pt, innertopmargin=6pt,
  innerbottommargin=6pt]
\textbf{Parallel machine strategy for Weeks 4--5:} Intel Mac runs Bluesky collector
continuously and Apify per-shock. Pi NPU classifies bios in real time as posts arrive
via Syncthing. Intel Mac runs the news scraper. HPC scores all collected text via
SLURM array jobs. M5 does interactive development. No machine is idle.
\end{mdframed}

\begin{mdframed}[backgroundcolor=warn, linecolor=darkblue, linewidth=0.8pt,
  innerleftmargin=10pt, innerrightmargin=10pt, innertopmargin=6pt,
  innerbottommargin=6pt]
\textbf{Routing strategy for the 8 weeks:} Use Gemini for paper summaries,
quick questions, drafts, and anything you would normally Google. Switch to
\texttt{@claude} for pipeline architecture, LaTeX writing, debugging complex
logic, and anything requiring the full project context. This keeps your Claude
Max usage reserved for real development sessions.
\end{mdframed}

% ════════════════════════════════════════════════════════════════════════════
\section{Week 0 --- Infrastructure \& Artifact Contracts (40 hrs)}

\textbf{Goal.} By Friday, the pipeline skeleton exists, every stage has a typed
artifact contract, and round-trip serialization tests pass. No economic logic is
implemented this week.

\textbf{End condition.}
\begin{itemize}[leftmargin=1.6em, itemsep=1pt]
  \item \done \texttt{pytest} passes with zero failures.
  \item \done All payload classes serialize and deserialize without data loss.
  \item \done Smoke pipeline runs and writes one artifact file per stage.
  \item \done PR merged to \texttt{main}.
\end{itemize}

% ── Day 1 ────────────────────────────────────────────────────────────────────
\dayhead{1}{Mon --- Repo Setup \& Project Scaffolding}

\begin{taskbox}
\done \task{Create the GitHub repository \texttt{electoral-equilibrium}. Initialize
  with \texttt{README.md}, \texttt{.gitignore} (Python), and \texttt{pyproject.toml}.
  Add project description, license (MIT), and a brief abstract in the README.}{1}

\done \task{\textbf{Apply for Meta Content Library access now} at
  \texttt{developers.facebook.com/docs/content-library}. Approval takes 2--4 weeks
  and is required for Week 4 Facebook data collection. Fill in the academic research
  use-case form: describe the electoral shock sentiment analysis project, cite
  Prof.\ Espinosa as faculty mentor, and list CMC as the institution.}{1}

\done \task{Set up the Python environment. Create a virtual environment, install
  core dependencies: \texttt{numpy}, \texttt{pandas}, \texttt{scipy}, \texttt{scikit-learn},
  \texttt{cvxpy}, \texttt{pytest}, \texttt{black}, \texttt{ruff}. Pin all versions
  in \texttt{pyproject.toml}.}{1}

\done \task{Create the full directory skeleton from the appendix of the development
  plan. Every folder and empty \texttt{\_\_init\_.py} file should exist. Commit as
  \texttt{chore: scaffold repository structure}.}{1}

\done \task{Write \texttt{AGENTS.md}: contributor conventions, branch naming rules
  (\texttt{feature/weekN-description}), PR format, and the pull-request checklist
  from the development plan.}{1}

\done \task{Write \texttt{DOMAIN.md}: the specification for instantiating this
  pipeline in a new domain. This is the generalization contract --- the document
  that makes the codebase a platform rather than a one-off tool. Required sections:
  (i) \textit{What is a domain} --- define stakeholder groups, loyalty metric,
  win condition, and equilibrium threshold; (ii) \textit{Config fields to set} ---
  list every field in \texttt{base.json} that is domain-specific;
  (iii) \textit{Data to provide} --- historical event-response dataset format,
  prediction market contract mapping file; (iv) \textit{Fine-tuning schema} ---
  confirm the JSONL format is domain-agnostic; (v) \textit{Example: electoral
  vertical} --- the current implementation as a worked example.
  The electoral build is instance 1 of a general framework; \texttt{DOMAIN.md}
  makes that explicit from day one.}{1}

\done \task{Create \texttt{electoral/core/types.py} with all type aliases:
  \texttt{Cycle}, \texttt{BlocId}, \texttt{Weight}, \texttt{ElasticityScore},
  \texttt{ShockDelta}, \texttt{LoraRank}. Write a short docstring for each
  explaining the convention (format, valid range, rationale).}{1}

\done \task{Create \texttt{configs/smoke.json} and \texttt{configs/base.json}.
  Smoke config uses two blocs, two cycles, and a tiny dataset path so the
  full pipeline can run end-to-end in under 10 seconds for testing.}{1}

\done \task{Create \texttt{electoral/core/rng.py}: implement \texttt{derive\_seed(tokens)}
  that hashes a list of string tokens into a reproducible integer seed, and
  \texttt{make\_rng(seed)} that returns a seeded \texttt{np.random.Generator}.
  Write three unit tests in \texttt{tests/test\_rng.py} confirming
  (i) identical inputs produce identical seeds, (ii) order matters, (iii)
  the generator produces identical draws across calls with the same seed.}{1}
\end{taskbox}

% ── Day 2 ────────────────────────────────────────────────────────────────────
\dayhead{2}{Tue --- Schema Helpers \& Artifact Envelope}

\begin{taskbox}
\done \task{Implement \texttt{electoral/core/schema.py} with all five validation
  helpers: \texttt{assert\_required\_keys}, \texttt{assert\_unique},
  \texttt{assert\_sorted\_unique}, \texttt{assert\_valid\_share},
  \texttt{assert\_shares\_sum\_to\_one}. Every helper must raise \texttt{ValueError}
  with a message naming the class and field. No external dependencies.}{3}

\done \task{Write \texttt{tests/test\_schema.py}. For each helper, write at least
  two positive tests (valid inputs pass silently) and two negative tests
  (invalid inputs raise \texttt{ValueError} with an informative message).
  Run \texttt{pytest -v} and confirm all pass.}{2}

\done \task{Create \texttt{electoral/core/io.py} implementing the Parquet + JSON
  artifact I/O layer (Option~A). For the envelope (\texttt{stage}, \texttt{run\_key},
  \texttt{metadata}): \texttt{write\_json(path, payload)} and \texttt{read\_json(path)}
  using \texttt{indent=2}, \texttt{sort\_keys=True}. For tabular \texttt{data}
  payloads (voter panel, scored posts, fine-tuning dataset): \texttt{write\_parquet(path,
  df)} using \texttt{engine="pyarrow", compression="snappy"} and
  \texttt{read\_parquet(path)} returning a \texttt{pd.DataFrame}. Non-tabular payloads
  (optimizer output, simulation results) remain JSON. \texttt{write\_artifact(path,
  envelope, df=None)} dispatches automatically. Write four tests: (i) JSON round-trip,
  (ii) Parquet round-trip with a 100-row fixture DataFrame, (iii) Parquet output is
  smaller than equivalent JSON (assert \texttt{parquet\_size < json\_size}), (iv)
  reading a non-existent path raises \texttt{FileNotFoundError}.}{3}

\done \task{Define the artifact envelope in \texttt{electoral/artifacts.py}.
  Every artifact must contain \texttt{stage}, \texttt{run\_key},
  \texttt{metadata}, and \texttt{data} fields. Write \texttt{StageArtifact}
  as a frozen dataclass with \texttt{to\_dict} and \texttt{from\_dict}.}{1}
\end{taskbox}

% ── Day 3 ────────────────────────────────────────────────────────────────────
\dayhead{3}{Wed --- Payload Classes (Part 1)}

\begin{taskbox}
\done \task{Implement \texttt{VoterPanelData} in \texttt{electoral/artifacts.py}.
  Fields: \texttt{cycles}, \texttt{blocs}, \texttt{columns}, \texttt{n\_rows},
  \texttt{source}. Implement \texttt{to\_dict}, \texttt{from\_dict},
  \texttt{validate} using schema helpers. Validation must check: cycles sorted
  unique, blocs unique, \texttt{n\_rows} non-negative.}{2}

\done \task{Implement \texttt{BaselinePortfolioData}. Fields: \texttt{method},
  \texttt{party}, \texttt{weights}, \texttt{mu}, \texttt{target}.
  \texttt{party} is \texttt{"democrat"} or \texttt{"republican"} and must
  be validated as such. \texttt{weights[bloc\_id]} is the optimized coalition
  weight in $[0,1]$; must sum to 1.0. \texttt{mu[bloc\_id]} is the historical
  mean within-bloc vote share \emph{for the selected party} in $[0,1]$.
  Validation: weights sum to 1.0 within \texttt{tol=1e-6}, all weight and mu
  values in $[0,1]$, both dicts have identical key sets matching the 10
  canonical bloc IDs, party is one of the two allowed values.}{2}

\done \task{Implement \texttt{SentimentData}. Fields: \texttt{model},
  \texttt{shocks}, \texttt{scores}. Validation must check: shocks unique,
  all score values in $[-1, 1]$, every bloc key in \texttt{scores} maps to
  a dict keyed by all shocks.}{2}

\done \task{Implement \texttt{LLMFineTuneData}. Fields: \texttt{base\_model},
  \texttt{lora\_rank}, \texttt{n\_examples}, \texttt{cycles\_used},
  \texttt{adapter\_path}. Validation: \texttt{lora\_rank} positive,
  \texttt{n\_examples} positive, \texttt{cycles\_used} sorted unique.}{2}
\end{taskbox}

% ── Day 4 ────────────────────────────────────────────────────────────────────
\dayhead{4}{Thu --- Payload Classes (Part 2) \& Round-Trip Tests}

\begin{taskbox}
\done \task{Implement \texttt{ShockResponseData}. Fields: \texttt{shock},
  \texttt{cycle}, \texttt{deltas}, \texttt{covariance}, \texttt{source}.
  \texttt{deltas[bloc\_id]} is the LLM-estimated change in within-bloc
  Democratic vote share ($\Delta\mu_i$). \texttt{source} must be one of
  \texttt{"llm\_unified"}, \texttt{"roberta\_news\_only"},
  \texttt{"roberta\_social\_only"} --- required for the three-way baseline
  comparison. Validation: all delta values finite and in $[-0.15, 0.15]$,
  covariance is square with side matching number of bloc keys, source
  non-empty string.}{2}

\done \task{Implement \texttt{EquilibriumData}. Fields: \texttt{method},
  \texttt{party}, \texttt{shock}, \texttt{weights}, \texttt{mu\_shifted},
  \texttt{feasible}, \texttt{target\_met}, \texttt{target}.
  \texttt{weights[bloc\_id]} is the rebalanced coalition weight $\tilde{w}_i$;
  \texttt{mu\_shifted[bloc\_id]} is the post-shock within-bloc vote share
  $\tilde{\mu}_i^{(P)}$ for the selected party. Both are required by the
  frontend to display coalition weight shifts and vote-share loyalty shifts
  for all 10 blocs. Validation: weights sum to 1.0, all weight and
  mu\_shifted values in $[0,1]$, party is valid, both dicts have identical
  key sets.}{2}

\done \task{Implement \texttt{SimulationData}. Fields: \texttt{n\_simulations},
  \texttt{seed}, \texttt{win\_probability}, \texttt{percentiles}. Validation:
  \texttt{win\_probability} in $[0,1]$, all percentile lists have exactly
  5 values (p5, p25, p50, p75, p95) each in $[0,1]$.}{1}

\done \task{Implement \texttt{MetricsTablesData}. Fields: \texttt{tables}.
  Validation: tables keys non-empty strings, values JSON-serializable
  (call \texttt{json.dumps} inside validate).}{1}

\done \task{Write \texttt{tests/test\_artifact\_roundtrip.py}. Implement the
  \texttt{assert\_roundtrip(obj)} helper. Write one round-trip test per
  payload class using the minimal instances from the development plan.
  Add at least two negative tests: duplicate cycles in \texttt{VoterPanelData},
  weights not summing to 1 in \texttt{BaselinePortfolioData}.}{3}
\end{taskbox}

% ── Day 5 ────────────────────────────────────────────────────────────────────
\dayhead{5}{Fri --- Pipeline Skeleton, Smoke Run \& PR}

\begin{taskbox}
\done \task{Implement \texttt{electoral/config.py}: \texttt{PipelineConfig}
  dataclass that loads \texttt{smoke.json} or \texttt{base.json}. Required
  fields: \texttt{run\_key}, \texttt{seed}, \texttt{data\_path}, \texttt{blocs},
  \texttt{party} (\texttt{"democrat"} or \texttt{"republican"}),
  \texttt{target} (equilibrium threshold: Democrat $\approx 0.535$, Republican
  $\approx 0.520$), \texttt{output\_dir}. Validate that \texttt{party} is one
  of the two allowed values and that \texttt{target} is in $(0.5, 0.7)$.
  Add \texttt{"party": "democrat", "target": 0.535} to \texttt{smoke.json}
  and \texttt{base.json}. Implement \texttt{derive\_seed} using \texttt{rng.py}.}{2}

\done \task{Create \texttt{electoral/stages.py}: one stub function per stage
  (\texttt{build\_voter\_panel}, \texttt{build\_baseline\_portfolio},
  \texttt{build\_sentiment\_data}, \texttt{build\_llm\_finetune},
  \texttt{build\_shock\_response}, \texttt{build\_optimization},
  \texttt{run\_simulations}). Each stub accepts \texttt{PipelineConfig},
  constructs a minimal placeholder payload using the correct dataclass,
  wraps it in a \texttt{StageArtifact}, and writes it to disk via
  \texttt{io.write\_artifact} (not \texttt{write\_json} --- tabular payloads
  go to Parquet). Every week replaces its stub with real logic.}{2}

\done \task{Write a \texttt{Justfile} at the repo root (install \texttt{just}
  via \texttt{brew install just}). Required targets: \texttt{smoke} (runs
  pipeline on \texttt{configs/smoke.json}), \texttt{test} (runs
  \texttt{pytest tests/ -v}), \texttt{score} (submits RoBERTa SLURM array),
  \texttt{train} (submits HPC fine-tuning job), \texttt{deploy} (runs
  \texttt{modal deploy deploy/modal\_app.py}),
  \texttt{sample} (submits \texttt{scripts/sample\_archives.py} as a
  SLURM array to draw stratified samples from raw archives),
  \texttt{clean} (runs \texttt{scripts/clean\_with\_llm.py} locally on
  the sampled dataset using Gemini Pro),
  and \texttt{continuous} (runs \texttt{prefect run pipeline.py --config
  configs/continuous.json} for nightly incremental updates post-SRP).
  This saves repeated long commands and eliminates flag typos
  across 8 weeks of daily use.}{1}

\done \task{Implement \texttt{electoral/core/io.py}: \texttt{write\_artifact(path,
  envelope)} writes the JSON envelope; for tabular \texttt{data} payloads
  (panel, scored posts, fine-tuning dataset) writes the data as Parquet via
  \texttt{pandas.DataFrame.to\_parquet()}. \texttt{read\_artifact(path)} reads
  the envelope and auto-detects Parquet vs JSON for the data payload.
  Parquet is 5--10$\times$ faster than JSON for large tabular data and
  smaller on disk (Syncthing benefit). Confirm round-trip: write a 1,000-row
  toy DataFrame, read it back, assert equality.}{2}

\done \task{Replace \texttt{electoral/pipeline.py} with a \textbf{Prefect flow}.
  Install \texttt{prefect}. Convert each stage function into a
  \texttt{@prefect.task} with \texttt{retries=2} and \texttt{retry\_delay\_seconds=30}.
  Wrap the full pipeline in a \texttt{@prefect.flow}. Benefits: automatic retry
  on failure, parallel task execution where tasks are independent, a visual
  dashboard (\texttt{prefect server start}) showing stage completion status,
  and persistent state so a crash restarts from the last successful stage.
  Run \texttt{prefect server start} locally and confirm the smoke pipeline
  appears in the dashboard.}{2}

\done \task{Weekly PR: open \texttt{feature/week1-artifact-contracts} against
  \texttt{main}. PR description must include: summary of what was implemented,
  the pre-submission checklist from \texttt{AGENTS.md}, and any design choices
  made. Run all tests locally and merge after they pass.}{1}

\done \task{Brief supervision check-in with Prof.\ Espinosa: confirm the canonical
  bloc names and equilibrium threshold $\veq = 0.535$. This is a domain
  confirmation only --- Juhi drives the agenda. Document decisions in
  \texttt{DECISIONS.md}.}{1}
\end{taskbox}

\newpage

% ════════════════════════════════════════════════════════════════════════════
\section{Week 1 --- Historical Data Pipeline (40 hrs)}
% ════════════════════════════════════════════════════════════════════════════

\textbf{Goal.} By Friday, the data kernel ingests raw survey exports from all
five sources and produces a validated \texttt{VoterPanelData} artifact.

\textbf{End condition.}
\begin{itemize}[leftmargin=1.6em, itemsep=1pt]
  \item \done Panel passes all invariants: unique tuples, sorted, finite values in $[0,1]$.
  \item \done Toy fixture (\texttt{toy\_panel.csv}) runs through the full cleaning pipeline.
  \item \done All data tests pass.
  \item \done PR merged.
\end{itemize}

% ── Day 1 ────────────────────────────────────────────────────────────────────
\dayhead{1}{Mon --- Data Source Audit \& Fixture Creation}

\begin{taskbox}
\done \task{Download and inspect raw data from all five sources: ARDA, GSS, Gallup,
  NEP, and Pew. For each source, document: available years, column names,
  demographic breakdowns, and any known data quality issues. Write findings in
  \texttt{rawdata/SOURCE\_NOTES.md}.}{3}

\done \task{Create \texttt{tests/fixtures/toy\_panel.csv}: a minimal 20-row CSV
  with columns \texttt{cycle}, \texttt{bloc}, \texttt{vote\_share}, \texttt{turnout},
  \texttt{source}. Cover all 10 blocs across two cycles (2016, 2020). Include one
  intentional dirty row (missing value, wrong type) to test cleaning logic.}{2}

\done \task{Identify and document the column name mapping for each real data source
  to the canonical panel schema. Store as \texttt{rawdata/column\_maps.json}:
  a dict mapping source name to a dict of \texttt{\{raw\_col: canonical\_col\}}.}{3}
\end{taskbox}

% ── Day 2 ────────────────────────────────────────────────────────────────────
\dayhead{2}{Tue --- Loaders \& Panel Utilities}

\begin{taskbox}
\done \task{Implement \texttt{electoral/data/loaders.py}: \texttt{load\_csv\_panel(path)}
  that reads a CSV with explicit dtypes (\texttt{cycle} as int, \texttt{bloc} as str,
  \texttt{vote\_share} and \texttt{turnout} as float64). No date inference.
  No random sampling. Returns a \texttt{pd.DataFrame}.}{2}

\done \task{Implement \texttt{electoral/data/panel.py}: \texttt{validate\_panel(df,
  required\_cols, context)} that enforces all five invariants from the development
  plan. Each violated invariant must raise a \texttt{ValueError} naming the
  class and field. Write five unit tests --- one per invariant.}{3}

\done \task{Implement source-specific loader wrappers in \texttt{electoral/data/loaders.py}:
  \texttt{load\_arda}, \texttt{load\_gss}, \texttt{load\_gallup}, \texttt{load\_nep},
  \texttt{load\_pew}. Each applies the correct column map from
  \texttt{column\_maps.json} and returns a raw (uncleaned) DataFrame.}{3}
\end{taskbox}

% ── Day 3 ────────────────────────────────────────────────────────────────────
\dayhead{3}{Wed --- Cleaning Logic}

\begin{taskbox}
\done \task{Implement \texttt{electoral/data/cleaning.py}: \texttt{clean\_raw\_panel(df)}
  that applies all required cleaning steps: coerce \texttt{cycle} to int YYYY,
  \texttt{bloc} to lowercase snake-case str, \texttt{vote\_share} and \texttt{turnout}
  to float64, drop rows with missing required values (log how many dropped),
  sort by (\texttt{cycle}, \texttt{bloc}), raise on duplicate
  (\texttt{cycle}, \texttt{bloc}, \texttt{source}) tuples rather than silently
  aggregating.}{3}

\done \task{Implement bloc name normalization in \texttt{cleaning.py}:
  \texttt{normalize\_bloc(raw\_name)} that maps source-specific demographic
  labels to the canonical bloc identifiers confirmed by Prof.\ Espinosa
  (e.g., \texttt{"White Evangelical"} $\to$ \texttt{"evangelical"}). Build the
  map from \texttt{DECISIONS.md}. Raise on unrecognized bloc names.}{3}

\done \task{Write \texttt{tests/test\_data\_cleaning.py}. Test: (i) clean path
  with toy fixture produces correct dtypes and sorting, (ii) dirty row is
  dropped and logged, (iii) unrecognized bloc name raises informatively,
  (iv) duplicate tuple raises rather than silently aggregating.}{2}
\end{taskbox}

% ── Day 4 ────────────────────────────────────────────────────────────────────
\dayhead{4}{Thu --- Data Kernel \& Multi-Source Merge}

\begin{taskbox}
\done \task{Implement \texttt{electoral/kernels/data.py}: \texttt{build\_voter\_panel(config)}
  that loads all five sources, applies source-specific loaders, concatenates,
  cleans via \texttt{clean\_raw\_panel}, validates via \texttt{validate\_panel},
  constructs a \texttt{VoterPanelData} payload, and returns both the payload and
  the cleaned DataFrame. Wire into \texttt{stages.py} replacing the Week~1 stub.}{3}

\done \task{Implement multi-source conflict resolution: when the same
  (\texttt{cycle}, \texttt{bloc}) pair appears in multiple sources with different
  \texttt{vote\_share} values, compute the weighted average (weight = inverse
  standard error if available, else equal weight). Document the rule in
  \texttt{DECISIONS.md}.}{3}

\done \task{Run the full data kernel on real data for all cycles 2000--2024.
  Inspect the output DataFrame. Log counts per source per cycle. Flag any
  bloc-cycle pairs with missing data across all sources. Document gaps in
  \texttt{rawdata/SOURCE\_NOTES.md}.}{2}
\end{taskbox}

% ── Day 5 ────────────────────────────────────────────────────────────────────
\dayhead{5}{Fri --- Integration, Visualization \& PR}

\begin{taskbox}
\done \task{Write a data audit script \texttt{scripts/audit\_panel.py} that loads
  the cleaned panel and prints: (i) coverage matrix (bloc $\times$ cycle showing
  which cells have data), (ii) vote-share summary statistics per bloc, (iii) any
  obvious outliers (vote share $< 0.05$ or $> 0.95$). Save output as
  \texttt{rawdata/audit\_report.txt}.}{2}

\done \task{Manually cross-check 10 spot values from the panel against the
  original source documents. Record the spot-check results in
  \texttt{rawdata/spot\_check.md}. Flag any large discrepancies for the
  supervision check-in.}{2}

\done \task{Weekly PR: \texttt{feature/week2-data-pipeline}. Ensure all tests
  pass, the audit report is committed, and the PR description summarizes any
  data quality decisions made. Merge after tests pass.}{1}

\done \task{Brief supervision check-in with Prof.\ Espinosa: share the coverage
  matrix and get domain input on any missing bloc-cycle cells. Juhi decides
  and implements the imputation strategy; document in \texttt{DECISIONS.md}.}{2}

\done \task{Implement the agreed imputation strategy (likely: linear interpolation
  for consecutive missing cycles, or carry-forward for isolated gaps) in
  \texttt{cleaning.py}. Add unit tests.}{1}
\end{taskbox}

\newpage

% ════════════════════════════════════════════════════════════════════════════
\section{Week 2 --- Baseline Voter Portfolio (40 hrs)}
% ════════════════════════════════════════════════════════════════════════════

\textbf{Goal.} By Friday, the ML pipeline produces a validated
\texttt{BaselinePortfolioData} artifact representing the optimal steady-state
voter equilibrium free of shocks.

\textbf{End condition.}
\begin{itemize}[leftmargin=1.6em, itemsep=1pt]
  \item \done Baseline weights sum to 1.0 and every weight is in $[0, 1]$.
  \item \done The portfolio achieves $\mathbf{w}^\top\boldmu^{(P)} \geq \veq^{(P)}$
        on historical winning cycles for the configured party.
  \item \done Leave-one-cycle-out validation documented for both parties.
  \item \done PR merged.
\end{itemize}

% ── Day 1 ────────────────────────────────────────────────────────────────────
\dayhead{1}{Mon --- Moment Estimation}

\begin{taskbox}
\done \task{Implement \texttt{electoral/models/ml\_baseline.py}:
  \texttt{estimate\_moments(df, party)} that computes $\boldmu^{(P)}$ (mean
  within-bloc vote share for party $P$ across winning cycles) and $\boldSigma$
  (empirical covariance across cycles). For \texttt{party="democrat"} use the
  \texttt{vote\_share} column directly; for \texttt{party="republican"} use
  $1 - \texttt{vote\_share}$. Define ``winning cycle'' as cycles where the
  selected party's candidate exceeded 50\% of the two-party vote share
  nationally --- this is party-configurable. Apply a PSD safeguard to
  $\boldSigma$: add $\epsilon I$ if the minimum eigenvalue is negative.}{3}

\done \task{Write helper \texttt{psd\_repair(cov, eps=1e-6)} in
  \texttt{ml\_baseline.py} that adds $\epsilon I$ to a matrix until it is
  positive semi-definite. Write a unit test confirming the output is PSD
  for a known near-singular matrix.}{2}

\done \task{Log and inspect the estimated $\boldmu$ and $\boldSigma$. Print
  bloc names alongside their mean vote shares. Cross-reference against
  Prof.\ Espinosa's published research (e.g., Evangelicals near 24\%).
  Flag any large discrepancies for the supervision check-in.}{3}
\end{taskbox}

% ── Day 2 ────────────────────────────────────────────────────────────────────
\dayhead{2}{Tue --- CVXPY Baseline Optimization}

\begin{taskbox}
\done \task{Implement \texttt{electoral/portfolios/cvx.py}: \texttt{solve\_baseline(mu,
  cov, target)} that formulates and solves the constrained quadratic program:
  $\min_{\mathbf{w}} \mathbf{w}^\top\boldSigma\mathbf{w}$ subject to
  $\sum w_i = 1$, $w_i \geq 0$, $\mathbf{w}^\top\boldmu \geq \veq$.
  Use \texttt{cvxpy.Problem}. Return weights as a dict keyed by bloc name.
  Raise informatively if problem is infeasible.}{3}

\done \task{Implement solver fallback in \texttt{cvx.py}: if CVXPY returns
  infeasible, relax the equilibrium constraint by 1\% increments and retry
  up to 5 times. Log each relaxation. If still infeasible after 5 retries,
  raise \texttt{ValueError} with the final constraint value attempted.}{2}

\done \task{Write \texttt{electoral/portfolios/weights.py}: \texttt{equal\_weight\_baseline(blocs)}
  and \texttt{value\_weight\_baseline(df, blocs)} as simple benchmark alternatives
  to the CVXPY optimizer. These serve as sanity checks and paper baselines.}{3}
\end{taskbox}

% ── Day 3 ────────────────────────────────────────────────────────────────────
\dayhead{3}{Wed --- Gaussian Process Classifier}

\begin{taskbox}
\done \task{Implement the \textbf{Gaussian Process} win-probability classifier
  in \texttt{ml\_baseline.py} (replacing XGBoost; Option~E).
  Use \texttt{sklearn.gaussian\_process.GaussianProcessClassifier} with a
  composite kernel: \texttt{RBF(length\_scale=1.0) + Matern(length\_scale=1.0,
  nu=1.5)}. Feature matrix $X$: per-bloc vote shares, turnout rates, cycle
  index (numeric, so the GP sees temporal ordering). Target $y$: binary win.
  The GP provides calibrated posterior predictive probabilities and explicitly
  models temporal correlations --- the 2020 cycle is weighted more heavily than
  2000 when predicting 2024. Extract the posterior standard deviation per
  held-out cycle and store in \texttt{artifacts/baseline\_loco.json} alongside
  accuracy. These uncertainty estimates feed the paper's epistemic humility
  framing.}{4}

\done \task{Write \texttt{tests/test\_baseline.py}: (i) \texttt{estimate\_moments}
  returns correct shapes, (ii) \texttt{psd\_repair} output is PSD, (iii)
  \texttt{solve\_baseline} weights sum to 1.0, (iv) GP classifier outputs
  calibrated probabilities in $[0,1]$, (v) infeasible QP raises correctly.}{4}
\end{taskbox}

% ── Day 4 ────────────────────────────────────────────────────────────────────
\dayhead{4}{Thu --- Kernel Integration \& Validation}

\begin{taskbox}
\done \task{Implement \text{electoral/kernels/baseline.py}:
  \texttt{build\_baseline\_portfolio(config, panel\_df)} that orchestrates
  moment estimation, CVXPY optimization, and payload construction.
  Wire into \texttt{stages.py} replacing the Week~1 stub.}{3}

\done \task{Run leave-one-cycle-out (LOCO) validation: for each cycle $c$,
  train on all other cycles, predict win probability on $c$, compare to
  ground truth. Record accuracy in \texttt{artifacts/baseline\_loco.json}.
  \textbf{Also} compare the GP's predicted win probability against the
  IEM/PredictIt market-implied price at the time of that election (for
  cycles 2004+). If the model diverges from the market by more than 10
  percentage points on a held-out cycle, investigate the moment estimates
  and constraint bounds before proceeding to Week~4 --- a large divergence
  signals a calibration problem, not just model uncertainty.}{4}

\done \task{Inspect the optimized weights. Cross-reference high-weight blocs
  against Prof.\ Espinosa's published demographic research. Write a 1-page
  interpretation in \texttt{docs/baseline\_interpretation.md}.}{1}
\end{taskbox}

% ── Day 5 ────────────────────────────────────────────────────────────────────
\dayhead{5}{Fri --- Constraints, Edge Cases \& PR}

\begin{taskbox}
\done \task{Implement \texttt{electoral/portfolios/constraints.py}:
  \texttt{ConstraintSpec} dataclass specifying simplex constraints plus
  any user-defined per-bloc bounds (e.g., Evangelicals must be at least 20\%).
  Integrate into \texttt{solve\_baseline}.}{2}

\done \task{Add edge case handling: (i) single-bloc input, (ii) all blocs
  equal weight produces feasible solution, (iii) target $> 1.0$ raises
  immediately without calling solver. Add tests for each.}{2}

\done \task{Weekly PR: \texttt{feature/week3-baseline-portfolio}. Include the
  LOCO validation results and interpretation doc in the PR. Merge after
  tests pass.}{1}

\done \task{Brief supervision check-in with Prof.\ Espinosa: share the baseline
  weights and LOCO results. Get domain feedback on whether the demographic
  ordering aligns with his published research. Juhi decides on any
  constraint adjustments and documents in \texttt{DECISIONS.md}.}{3}

\done \task{Implement \texttt{electoral/kernels/raking.py}: marginal calibration
  (iterative proportional fitting) as an optional post-processing step on the
  stratum weights from \texttt{build\_constraint\_spec}. The algorithm:
  (1) rescale race weights to sum to 1 matching voter panel race marginals;
  (2) rescale religion weights matching religion marginals;
  (3) rescale gender weights matching gender marginals;
  repeat until convergence (tolerance $< 10^{-6}$, typically $< 10$ iterations).
  Output the raked weights under key \texttt{"raked"} in
  \texttt{configs/layer\_weights.json} alongside the additive \texttt{"additive"}
  weights. Log the per-iteration change and number of iterations to converge.
  This is the lightweight middle ground between the additive model and full MRP;
  it enforces joint consistency across all three marginals without requiring
  cross-tabulation data. The paper must compare additive vs.\ raked outputs.}{4}
\end{taskbox}

\newpage

% ════════════════════════════════════════════════════════════════════════════
\section{Week 3 --- RoBERTa Pipeline, Social Media \& Fine-Tuning Dataset (40 hrs)}
% ════════════════════════════════════════════════════════════════════════════

\textbf{Goal.} By Friday, two parallel sentiment pipelines are running --- a news
corpus scraper and a four-platform social media collector. Both are scored with
RoBERTa, social media sentiment is correlated with 14-day lagged favorability
polls, and a unified fine-tuning dataset is assembled.

\textbf{End condition.}
\begin{itemize}[leftmargin=1.6em, itemsep=1pt]
  \item \done \texttt{SentimentData} and \texttt{SocialMediaSentimentData} artifacts
        validate and round-trip.
  \item \done At least 400 fine-tuning examples in \texttt{train.jsonl}, fusing
        both news and social media signals.
  \item \done Baseline-adjusted social sentiment computed for all platforms.
  \item \done \texttt{eval.jsonl} contains 2020 cycle examples only.
  \item \done PR merged.
\end{itemize}

\begin{mdframed}[backgroundcolor=warn, linecolor=darkblue, linewidth=0.8pt,
  innerleftmargin=10pt, innerrightmargin=10pt, innertopmargin=6pt, innerbottommargin=6pt]
\textbf{Hard fallbacks --- read before starting Week 4.}

\textbf{Meta Content Library:} If not approved by Week 4 Day 1, do not wait.
Activate the fallback immediately: scrape public Facebook page reaction counts
(like, angry, sad, love) from major news pages as a valence-only signal. The
\texttt{FacebookCollector} stub must work in both modes --- full Content Library
and reaction-count fallback --- before Day 2 ends.

\textbf{Social media collectors:} Bluesky and Apify should both be running from
Week~0. If Apify's free tier runs out for a given shock, reduce the per-run result
count or skip that shock for Apify and rely on Bluesky alone. The historical archives
(MeToo hydrated dataset, COVID-19, BLM/George Floyd 2020, 2020 election) are the primary source of Twitter-equivalent
signal --- live collection is supplementary.

\textbf{Hailo NPU:} If Week 0 compilation failed, confirm \texttt{pi\_npu\_enabled=false}
in \texttt{configs/base.json}. The CPU fallback path must be running and tested
before Day 3's SetFit training begins.
\end{mdframed}

\begin{taskbox}
\textbf{Hard fallback implementation checklist} --- all three must be complete before
the deadlines shown. These are not optional; the pipeline is blocked on them.

\begin{itemize}[leftmargin=1.6em, itemsep=4pt]
  \item \done \textbf{[Day 1 end]} \texttt{scripts/compile\_hailo.py}
        implemented. The script exports \texttt{all-MiniLM-L6-v2} to ONNX,
        runs \texttt{hailo optimize} and \texttt{hailo compile}, verifies the
        HEF loads, and writes \texttt{pi\_npu\_enabled=true} to
        \texttt{configs/base.json} on success. On \emph{any} failure it writes
        \texttt{pi\_npu\_enabled=false} and exits 1. Week 3 is never blocked.
        \textbf{File:} \texttt{scripts/compile\_hailo.py}

  \item \done \textbf{[Day 2 end]} \texttt{FacebookCollector} implemented in
        dual-mode --- \emph{not stubs}. \texttt{mode="content\_library"} uses
        the Meta Content Library API (requires \texttt{FACEBOOK\_CL\_TOKEN}).
        \texttt{mode="reaction\_fallback"} uses the Graph API to scrape public
        page reaction counts (LIKE, LOVE, WOW, HAHA, ANGRY, SAD) from major
        news pages and computes a valence score per post.
        \texttt{mode="auto"} (default) auto-selects based on whether
        \texttt{FACEBOOK\_CL\_TOKEN} is set in the environment.
        Both modes write the same \texttt{posts.jsonl} schema;
        reaction-fallback records carry a \texttt{reaction\_valence} field
        consumed by \texttt{elasticity.py}.
        \textbf{File:} \texttt{electoral/nlp/social.py}

  \item \done \textbf{[Day 2 end]} \texttt{ApifyCollector} quota fallback
        implemented. On actor-run failure, \texttt{\_call\_actor()} returns
        \texttt{None} (not raise). \texttt{run()} detects the \texttt{None}
        result, retries once with \texttt{max\_items=100}
        (\texttt{MAX\_ITEMS\_QUOTA\_RETRY}), and if that also fails, logs a
        warning and returns~0. The pipeline falls back to Bluesky alone for
        that shock and continues without crashing.
        \textbf{File:} \texttt{electoral/nlp/collectors/apify\_x\_scraper.py}

  \item \done \textbf{[Day 3 start — before SetFit training]} Hailo CPU
        fallback verified. Run:
        \texttt{python scripts/pi\_bio\_server.py --config configs/base.json
        --cpu-only} and confirm \texttt{curl http://pi.local:9000/health}
        returns \texttt{\{"status":"ok","mode":"cpu",...\}}. The
        \texttt{/health} endpoint returns both \texttt{"mode"} (canonical,
        \texttt{"cpu"}~or~\texttt{"npu"}) and \texttt{"inference\_backend"}
        (legacy alias). Do \textbf{not} start Day 3 SetFit training until
        this check passes.
        \textbf{File:} \texttt{scripts/pi\_bio\_server.py}
\end{itemize}
\end{taskbox}

% ── Day 1 ────────────────────────────────────────────────────────────────────
\dayhead{1}{Mon --- Shock Taxonomy, API Setup \& News Scraper}

\begin{taskbox}
\done \task{Build the shock event registry: \texttt{configs/shocks.json} listing
  all historical shock events 2000--2024. Each entry: \texttt{id} (slug),
  \texttt{date}, \texttt{category} (Security, Geopolitical, Moral/Scandal,
  Electoral Surprise), \texttt{description} (one sentence),
  \texttt{target\_blocs}, \texttt{keywords}, and \texttt{archive\_ids}
  (list of pre-collected dataset slugs that cover this event, or empty list
  if live collection only). Aim for at least 25 events. Mark the following
  shocks as archive-covered in \texttt{archive\_ids}:

  \smallskip
  \begin{itemize}[leftmargin=1.6em, itemsep=1pt]
    \item \texttt{metoo\_2017}: \texttt{["metoo\_liwc", "reddit\_pushshift"]} ---
          Foster \& Rathlin 2022 LIWC dataset (3,683 tweets, CC~BY-NC-SA~4.0,
          DOI:~10.5683/SP3/1YWCW1) plus Reddit Pushshift subreddits
          (\texttt{TwoXChromosomes}, \texttt{Feminism}, \texttt{AskFeminists},
          \texttt{AskWomen}) filtered to Oct~15--Dec~31, 2017 at sampling time.
          Shock category: Moral/Scandal. 15~October 2017.
    \item \texttt{kavanaugh\_2018}: \texttt{["metoo\_liwc", "reddit\_pushshift"]} ---
          Foster \& Rathlin dataset covers \texttt{kavanaugh\_ford}
          (27~Sept 2018) and \texttt{kavanaugh\_confirmed} (6~Oct 2018) phases.
          Reddit Pushshift filtered to Sept--Oct 2018 at sampling time.
          Shock category: Moral/Scandal (Electoral). Oct 2018.
    \item \texttt{covid\_pandemic\_2020}: \texttt{["covid\_pandemic\_2020"]} ---
          $\sim$179k tweets, gpreda Kaggle dataset.
          Shock category: Geopolitical/Health. 11~March 2020, WHO pandemic
          declaration.
    \item \texttt{covid\_vaccine\_2020}: \texttt{["covid\_vaccine\_2020"]} ---
          $\sim$200k tweets, kaushiksuresh147 Kaggle dataset.
          \texttt{user\_description} confirmed present --- use for bio
          classification. Shock category: Geopolitical/Health.
          December 2020, first vaccine approvals.
    \item Any 2020 cycle events: \texttt{["election\_2020\_ieee",
          "election\_2020\_kaggle"]} as applicable
    \item \texttt{blm\_george\_floyd\_2020}: \texttt{["reddit\_pushshift"]} ---
          \textbf{Twitter BLM coverage dropped (2026-06-03).}
          Reddit Pushshift \texttt{r/BlackLivesMatter} and
          \texttt{r/BlackPeopleTwitter} filtered to May~25--August~31, 2020
          is the sole training source.
          Shock category: Moral/Scandal (Social Movement).
          Primary signal for African American, secular, and diverse bloc delta bins.
    \item \texttt{daca\_rescission\_2017}: \texttt{[]} --- live collection only;
          no pre-collected archive. Shock category: Geopolitical (Immigration).
          Trump rescinds DACA 5~Sept 2017. Target blocs: Latino, Asian,
          Muslim, Secular.
    \item \texttt{daca\_scotus\_2020}: \texttt{["daca\_scotus\_2020"]} ---
          Schema: \texttt{status}$\to$\texttt{text},
          \texttt{timestamp}$\to$\texttt{created\_at};
          no \texttt{lang} field (hardcode \texttt{"en"});
          no \texttt{user.description} (apply Latino/diverse platform proxy,
          \texttt{inference\_method="platform\_proxy"}).
          Shock category: Geopolitical (Immigration). 18~June 2020,
          SCOTUS blocks Trump DACA rescission 5--4 (\emph{DHS v.\ Regents of
          the University of California}). Target blocs: Latino, Asian, Secular.
    \item \texttt{election\_2016}: \texttt{["election\_2016"]} ---
          kinguistics Kaggle dataset, Election Day tweets only.
          Shock category: Electoral. 8~November 2016, Trump defeats Clinton.
          Major shock for Latino, Muslim, Secular, and Jewish blocs.
          Manually downloaded to
          \texttt{/Volumes/JUHIDRIVE/electoralData/archives/twitter/}\allowbreak
          \texttt{election\_2016/}.
          \textbf{Schema verification required on first load:} confirm
          \texttt{user\_description} field is present before assigning final
          tier. If present: Tier~1 (bio classification). If absent: Tier~2
          (platform proxy).
    \item \texttt{election\_2012}: \texttt{["election\_2012"]} ---
          jgoodman8 Kaggle dataset, $\sim$1M tweets, JSON format.
          Shock category: Electoral. 6~November 2012, Obama defeats Romney.
          Baseline election; lower shock value than 2016 but useful for
          pre-Trump sentiment calibration.
          Manually downloaded to
          \texttt{/Volumes/JUHIDRIVE/electoralData/archives/twitter/}\allowbreak
          \texttt{election\_2012/}.
          \textbf{Schema verification required on first load:} confirm
          \texttt{user\_description} field is present before assigning final
          tier. If present: Tier~1 (bio classification). If absent: Tier~2
          (platform proxy).
    \item \texttt{election\_2024}: \texttt{["truth\_social\_2024"]} ---
          USC Humans Lab / kashish-s Kaggle dataset (1.5M posts, Feb 2022--Oct 2024)
          as primary Truth Social source; plus \texttt{"x\_2024"} if sinking8
          schema check confirms full objects (not IDs only --- still pending).
          Shock category: Electoral. Covers three sub-events:
          Trump assassination attempt (13~July 2024),
          Biden dropout from race (21~July 2024),
          and Election Day Trump victory (5~November 2024).
          Target blocs: all five race blocs, Evangelical, Catholic, Secular,
          women (gender gap). Primary Truth Social source is \texttt{truth\_social\_2024}.
  \end{itemize}
  \smallskip
  Use the Ayatollah assassination as the pilot/test case throughout this week.}{2}

\done \task{Verify all Week~0 social media setup is fully operational before
  writing any Week~4 code: (i) \texttt{BlueSkyCollector} returns posts with
  author \texttt{description} and \texttt{langs} fields for the Ayatollah
  keywords; (ii) \texttt{ApifyCollector} returns posts in the correct schema;
  (iii) the Pi bio server responds at \texttt{http://pi.local:9000/health}
  in either NPU or CPU fallback mode; (iv) Syncthing is syncing between
  Intel Mac and M5. Log the status of both collectors in \texttt{DECISIONS.md}
  with the date.

  \smallskip
  \textbf{Resolution (2026-06-03):}
  \begin{itemize}[leftmargin=1.6em, itemsep=1pt]
    \item[(i)] \textbf{BlueSkyCollector} --- intentionally deferred.
          Archive-first strategy adopted; live collectors not needed for
          historical shock training data. To be configured only if live
          shock collection is needed post-SRP. Logged in \texttt{DECISIONS.md}.
    \item[(ii)] \textbf{ApifyCollector} --- intentionally deferred.
          Same reason as BlueSky. Logged in \texttt{DECISIONS.md}.
    \item[(iii)] \textbf{Pi bio server} --- \textbf{COMPLETE.}
          Confirmed running at \texttt{http://100.125.90.19:9000/health}
          in CPU fallback mode, returning
          \texttt{\{"status":"ok","mode":"cpu","inference\_backend":"cpu",}%
          \texttt{"model":"all-MiniLM-L6-v2"\}}.
          Hailo NPU compilation skipped.
          Running from \texttt{feature/week4-roberta-social-pipeline} branch.
    \item[(iv)] \textbf{Syncthing} --- \textbf{COMPLETE.}
          Operational and syncing across local machines.
  \end{itemize}
  \smallskip
  Ayatollah assassination pilot test not run (live collectors deferred).
  Pilot shock will instead use archive data from the 2020 election dataset
  at sample time.}{1}

\done \task{Compile \texttt{all-MiniLM-L6-v2} for the Pi AI HAT+ NPU.
  On the Pi, run \texttt{scripts/compile\_hailo.py}: (i) load the model
  via \texttt{sentence-transformers}, (ii) export to ONNX with a dummy
  input of shape $(1, 128)$, (iii) run \texttt{hailo optimize minilm.onnx
  --hw-arch hailo8l}, (iv) run \texttt{hailo compile minilm.onnx
  --hw-arch hailo8l -o models/minilm.hef}. Verify the HEF loads:
  \texttt{from hailo\_platform import HEF; HEF("models/minilm.hef")}.
  This step can run while the Intel Mac setup runs in parallel.

  Then start the bio server: \texttt{uvicorn pi\_bio\_server:app
  --host 0.0.0.0 --port 9000}. \textbf{Hostname:} \texttt{pi.local}
  requires mDNS/Bonjour. If it does not resolve from the Intel Mac or M5,
  assign a static IP to the Pi in your router, add it to \texttt{/etc/hosts}
  on all machines, and update \texttt{"pi\_bio\_server"} in
  \texttt{configs/base.json}. Confirm reachability:
  \texttt{curl http://pi.local:9000/health}.

  \smallskip
  \textbf{Resolution (2026-06-03):} Hailo NPU compilation \textbf{deferred
  indefinitely.} The Hailo developer SDK is gated behind vendor registration,
  and the Hailo compiler must run on the host Mac (not the Pi itself) ---
  adding a toolchain dependency that is not justified for the SRP timeline.
  NPU compilation will only be revisited if CPU inference becomes a bottleneck
  after SetFit training is complete.

  \smallskip
  \textbf{CPU fallback mode is active and sufficient.} Bio server confirmed
  running at \texttt{http://100.125.90.19:9000/health} (static IP; mDNS
  blocked on CMC WiFi). Response:
  \texttt{\{"status":"ok","mode":"cpu","inference\_backend":"cpu",}%
  \texttt{"model":"all-MiniLM-L6-v2"\}}.
  CPU inference ($\sim$20\,ms/bio on Pi 5) is acceptable for the archive
  scoring workload.}{2}

\begin{mdframed}[backgroundcolor=warn, linecolor=darkblue, linewidth=1.5pt,
  innerleftmargin=12pt, innerrightmargin=12pt, innertopmargin=10pt, innerbottommargin=10pt]
\textbf{\large \textrightarrow\ Mount JUHIDRIVE before running any pipeline script.}

\smallskip
\textbf{Drive:} \texttt{/Volumes/JUHIDRIVE} --- 5\,TB exFAT, shared between Mac and
Pi over the local network.
All raw archives, sampled JSONL, and training data live here;
\textbf{nothing in \texttt{archives/} is committed to git.}

\medskip
\begin{tabular}{@{}ll@{}}
  \textbf{Mac path:} & \texttt{/Volumes/JUHIDRIVE/electoralData/archives/} \\
  \textbf{Pi path:}  & set \texttt{"pi\_data\_root"} in \texttt{configs/base.json} \\
\end{tabular}

\medskip
\textbf{Archive layout on JUHIDRIVE:}
\begin{itemize}[leftmargin=1.6em, itemsep=1pt]
  \item \texttt{archives/twitter/}    --- Twitter/X datasets (by shock slug)
  \item \texttt{archives/reddit/}     --- Pushshift subreddit dumps (.zst)
  \item \texttt{archives/telegram/}   --- Telegram message exports
  \item \texttt{archives/discord/}    --- Discord server exports
  \item \texttt{archives/snap/}       --- Snapchat Snap Map data (future)
  \item \texttt{archives/validation/} --- bot blocklists, held-out sets
  \item \texttt{archives/news/}       --- 3DLNews2 URL metadata + Webhose dumps
\end{itemize}

\medskip
\textbf{Pi network access.}
The Pi reads and writes archives over the local network via a Samba or NFS mount of
JUHIDRIVE. Set \texttt{"pi\_data\_root"} in \texttt{configs/base.json} to the
Pi-side mount path (e.g.\ \texttt{/mnt/juhidrive/electoralData/archives/})
\textbf{before} running any pipeline script on the Pi.
This field is separate from \texttt{"data\_root"}, which is the Mac-side path.
\end{mdframed}

\begin{mdframed}[backgroundcolor=warn, linecolor=darkblue, linewidth=0.8pt,
  innerleftmargin=10pt, innerrightmargin=10pt, innertopmargin=6pt, innerbottommargin=6pt]
\textbf{Storage strategy --- read before downloading anything.}

All raw archives write directly to JUHIDRIVE under
\texttt{/Volumes/JUHIDRIVE/electoralData/archives/\{platform\}/\{slug\}/}.
Use the platform subdirectory convention: \texttt{twitter/}, \texttt{reddit/},
\texttt{telegram/}, \texttt{discord/}, \texttt{snap/}, \texttt{validation/},
\texttt{news/}.

\begin{itemize}[leftmargin=1.6em, itemsep=1pt]
  \item \textbf{JUHIDRIVE} is the primary archive storage.
        The 5\,TB capacity means raw archives do not need to be deleted after
        sampling --- retain them for schema re-checks and re-sampling runs.
  \item \textbf{3DLNews2:} download only the URL metadata file to
        JUHIDRIVE \texttt{archives/news/3dlnews/} (small). Filter by date window
        and state \emph{before} fetching any HTML.
        Run the HTML fetch-and-parse as a SLURM array job on HPC scratch so raw
        HTML is never written to local disk. \texttt{rsync} only the extracted
        \texttt{articles.jsonl} back to JUHIDRIVE.
  \item \textbf{Syncthing} syncs processed outputs only (sampled JSONL, cleaned
        JSONL, scored embeddings). Never configure Syncthing to sync raw archives.
  \item \textbf{M5} holds no raw archives --- processed outputs only.
\end{itemize}
\end{mdframed}

\done \task{Download all pre-collected archives and inspect their schemas.
  Save each to \texttt{/Volumes/JUHIDRIVE/electoralData/archives/\{platform\}/\{slug\}/}
  and document schema findings in
  \texttt{/Volumes/JUHIDRIVE/electoralData/archives/README.md} before Day 2.

  \smallskip
  \textbf{Tier 1 --- full objects, directly usable:}
  \begin{itemize}[leftmargin=1.6em, itemsep=1pt]
    \item \textbf{MeToo / Kavanaugh LIWC training set (Foster \& Rathlin 2022):}
          DOI:~\texttt{10.5683/SP3/1YWCW1}. Local path:
          \texttt{validation/metoo\_liwc/fulldataCHBR.sav} (2.81~MB,
          3,683 tweets). Schema: \texttt{Tweet}$\to$\texttt{text},
          \texttt{Event}$\to$\texttt{shock\_slug}. No \texttt{user\_description}
          --- apply secular/diverse platform proxy to all posts.
          Four shock phases: \texttt{metoo\_2017} (Oct 15, 2017),
          \texttt{kavanaugh\_ford} (Sept 27, 2018),
          \texttt{kavanaugh\_confirmed} (Oct 6, 2018),
          \texttt{weinstein\_convicted} (Feb 24, 2020).
          \textbf{License: CC~BY-NC-SA~4.0 --- flag for commercial review.}
          Use as training signal despite small size; supplement with Reddit
          Pushshift at sampling time.
          \textbf{Note:} The Ruest 25.2M Harvard Dataverse corpus
          (DOI~10.7910/DVN/2SRSKJ) is \textbf{permanently dropped} ---
          do not reference or attempt to download.
    \item \textbf{2020 Election --- IEEE (3.5M, July--Aug 2020):}
          \url{ieee-dataport.org/open-access/us-november-2020-election-tweets-dataset}.
    \item \textbf{2020 Election --- Kaggle (Oct 15--Nov 4 2020):}
          \url{kaggle.com/manchunhui/us-election-2020-tweets}.
    \item \textbf{[DROPPED 2026-06-03] BLM / George Floyd 2020 Twitter (alenjose,
          TweetBLM.csv, BLM2020\_Dataset.csv) ---} Extensive search found no suitable
          Twitter/X BLM 2020 dataset with \texttt{user\_description}: all
          candidates are IDs-only, wrong language, wrong time period, or too
          small. The alenjose 100k dataset is entirely 2023 data.
          TweetBLM.csv and BLM2020\_Dataset.csv also rejected for the same reasons.
          Training signal for \texttt{blm\_george\_floyd\_2020} comes entirely
          from Reddit Pushshift \texttt{r/BlackLivesMatter} and
          \texttt{r/BlackPeopleTwitter} filtered to
          May~25--August~31, 2020. Delete \texttt{twitter/blm\_2020/} from
          the archive directory.
    \item \textbf{Congress tweets (ongoing JSON):}
          \url{alexlitel.github.io/congresstweets/}.
    \item \textbf{Polarization corpus (Kiran et al.):}
          \url{gvrkiran.github.io/polarizationTwitter/}.
          \textbf{High priority.} Add bio text + ideological label to
          \texttt{data/bio\_labels/labeled\_bios.jsonl} before SetFit
          training (Week 4 Day 3).
    \item \textbf{USC Telegram 2024 election (leonardo-blas):}
          \url{github.com/leonardo-blas/usc-tg-24-us-election} and
          \url{huggingface.co/datasets/leonardoblas/}\allowbreak
          \texttt{us\_election\_2024\_telegram\_distilled}.
          No user bio field --- apply Evangelical/conservative Protestant
          platform-level proxy to all posts. Use \texttt{TelegramLoader}.
    \item \textbf{Truth Social --- notmooodoo9 (PRIMARY, 31.8M comments + 18k posts):}
          HuggingFace: \texttt{notmooodoo9/TrumpsTruthSocialPosts}.
          Downloaded to \texttt{/Volumes/JUHIDRIVE/electoralData/archives/}\allowbreak
          \texttt{truthsocial/notmooodoo9/}. Slug: \texttt{truth\_social\_2025}.
          License: CC-BY 4.0, commercial use permitted.

          \smallskip
          Three files: \texttt{truthsocial.comments[Trump-FROM-10-8-25].csv}
          (31.8M rows, $\sim$5.8\,GB),
          \texttt{truthsocial.posts[Trump-FROM-10-8-25].csv}
          (18,476 Trump posts, Feb 2022--Oct 2025),
          \texttt{truthsocial.users[Trump-FROM-10-8-25].csv}
          (1.5M rows --- \textbf{crawler state table, not profile data};
          no bio, description, or display\_name field).

          \smallskip
          \textbf{Schema (comments):} \texttt{\_id} (Mastodon snowflake ---
          decode date via \texttt{datetime.utcfromtimestamp((id >> 16) / 1000)}),
          \texttt{owner} (user ID, no bio available),
          \texttt{reply\_to},
          \texttt{text}.
          No \texttt{created\_at} column --- date must be derived from
          \texttt{\_id}.

          \smallskip
          \textbf{Tier 2} --- no bio field anywhere in dataset.
          Apply Evangelical/MAGA/conservative Protestant platform proxy to all
          comments (\texttt{inference\_method="platform\_proxy"}).
          Exclude from $\Sigma_\Delta$ covariance estimation.

          \smallskip
          \textbf{Date range:} comments start 8~October 2025 (post-election
          Trump activity). Trump's posts go back to Feb 2022; these cover
          shocks such as \texttt{trump\_indictment\_2023},
          \texttt{israel\_hamas\_war\_2023}, \texttt{trump\_conviction\_2024},
          and \texttt{trump\_assassination\_attempt\_2024}, but as
          retrospective Truth Social posts rather than reaction-window
          social media --- use for baseline sentiment only, not shock window
          scoring.

    \item \textbf{Truth Social --- Zenodo Notre Dame (\texttt{truth\_social\_2022},
          Feb--Sept 2022, 823K posts):}
          Downloaded to \texttt{/Volumes/JUHIDRIVE/electoralData/archives/}\allowbreak
          \texttt{truthsocial/zenodo\_notredame/}. Slug: \texttt{truth\_social\_2022}.
          Network-structure scrape (12 TSV files): users (454k), truths (854k raw lines,
          $\sim$824k well-formed), follows, replies, hashtags, media, external URLs.
          \textbf{No bio field} in users.tsv --- \texttt{profile\_url} is a link,
          not text. \textbf{Tier 2}: apply Evangelical/conservative Protestant
          platform proxy. Filter \texttt{is\_retruth=f}; decode dates from
          Mastodon snowflake \texttt{\_id}. Use \texttt{csv.QUOTE\_NONE} for
          truths.tsv (embedded newlines in $\sim$30,500 rows).

    \item \textbf{Truth Social --- USC Humans Lab / kashish-s Kaggle
          (PRIMARY 2024 election source, \texttt{truth\_social\_2024}):}
          \url{kaggle.com/datasets/kashishashah/truthsocial-2024-election-integrity-initiative}.
          Led by Prof.\ Emilio Ferrara, USC Humans Lab.
          1.5~million posts, February 2022 through October 2024.
          Save to \texttt{/Volumes/JUHIDRIVE/electoralData/archives/}\allowbreak
          \texttt{truthsocial/kashish\_usc/}. Slug: \texttt{truth\_social\_2024}.

          \smallskip
          \textbf{Shock coverage:} covers the full 2024 election window including
          Trump assassination attempt (13~July 2024), Biden dropout (21~July 2024),
          and Election Day (5~November 2024). This is the primary Truth Social
          archive for \texttt{election\_2024} shock analysis.

          \smallskip
          \textbf{Tier 2} --- no user bio field in Truth Social API data at
          scrape time. Apply Evangelical/MAGA/conservative Protestant platform proxy
          (\texttt{inference\_method="platform\_proxy"}).
          Exclude from $\Sigma_\Delta$ covariance estimation.

          \smallskip
          Note: \texttt{github.com/kashish-s/TruthSocial\_2024ElectionInitiative}
          is the scraping scripts repository only --- it contains no data.
          The actual dataset is the Kaggle link above.
    \item \textbf{[DROPPED 2026-06-03] BLM / George Floyd 2020 Twitter ---}
          see note above. \texttt{twitter/blm\_2020/} removed from archive structure.
    \item \textbf{COVID-19 tweets (gpreda):}
          \url{kaggle.com/datasets/gpreda/covid19-tweets}.
          $\sim$179k tweets with \#covid19, March--April 2020. Twitter API
          scrape --- confirm \texttt{user\_description} field is present on
          download. Field mapping: \texttt{text}$\to$\texttt{text},
          \texttt{date}$\to$\texttt{created\_at},
          \texttt{user\_description}$\to$\texttt{user.description}.
          Maps to \texttt{covid\_pandemic\_2020} shock slug; covers the WHO
          pandemic declaration shock window (March~11, 2020).
          Save to \texttt{/Volumes/JUHIDRIVE/electoralData/archives/twitter/}\allowbreak
          \texttt{covid\_pandemic\_2020/}.
    \item \textbf{COVID Vaccine tweets (kaushiksuresh147):}
          \url{kaggle.com/datasets/kaushiksuresh147/covidvaccine-tweets}.
          $\sim$200k tweets with \#CovidVaccine, Jan 2020--March 2021.
          Schema confirmed: \texttt{user\_name}, \texttt{user\_location},
          \texttt{user\_description}, \texttt{user\_created},
          \texttt{user\_followers}, \texttt{user\_verified}, \texttt{date},
          \texttt{text}, \texttt{hashtags}, \texttt{source},
          \texttt{is\_retweet}. Tier~1 --- \texttt{user\_description}
          confirmed present; use directly for bio classification.
          Maps to \texttt{covid\_vaccine\_2020} shock slug; covers vaccine
          rollout and approval shock window (Dec~2020--March~2021).
          Save to \texttt{/Volumes/JUHIDRIVE/electoralData/archives/twitter/}\allowbreak
          \texttt{covid\_vaccine\_2020/}.
  \end{itemize}

  \smallskip
  \textbf{Tier 2 --- scorer / demographic training:}
  \begin{itemize}[leftmargin=1.6em, itemsep=1pt]
    \item \textbf{Political tweets (kaushiksuresh147):}
          \url{kaggle.com/datasets/kaushiksuresh147/political-tweets}.
    \item \textbf{Political partisanship (lingshuhu):}
          \url{kaggle.com/datasets/lingshuhu/political-partisanship-tweets}.
    \item \textbf{Political tweets + gender (lingshuhu):}
          \url{kaggle.com/datasets/lingshuhu/political-tweets-with-gender-info}.
    \item \textbf{PoliticalTweets (HuggingFace):}
          \url{huggingface.co/datasets/Jacobvs/PoliticalTweets}.
    \item \textbf{US Politicians tweets (mrmorj):}
          \url{kaggle.com/datasets/mrmorj/us-politicians-twitter-dataset}.
    \item \textbf{TikTok 2024 election (gabbypinto):}
          \url{github.com/gabbypinto/US2024PresElectionTikToks}.
          Load caption text as \texttt{text} field; apply younger/secular/diverse
          proxy. Use \texttt{TikTokLoader}. No bio field available.
    \item \textbf{DACA SCOTUS 2020 tweets (schairezv):}
          \url{github.com/schairezv/daca-tweets-analysis}.
          Several thousand tweets, June 2020. Covers shock slug
          \texttt{daca\_scotus\_2020}. Schema:
          \texttt{status}$\to$\texttt{text},
          \texttt{timestamp}$\to$\texttt{created\_at};
          no \texttt{lang} field --- hardcode \texttt{"en"}.
          No \texttt{user.description} --- apply Latino/diverse platform proxy
          to all posts (\texttt{inference\_method="platform\_proxy"}).
          No \texttt{user\_id} --- bot blocklist cannot be applied; skip
          deduplication step that depends on author ID.
          Pre-computed \texttt{polarity} and \texttt{sentiment} fields are
          sanity-check reference only --- \textbf{do not} use as training
          inputs or fold into $\Delta\mu$ estimation.
  \end{itemize}

  \smallskip
  \textbf{ID-only / rehydration backlog --- not usable at current API tier:}
  \begin{itemize}[leftmargin=1.6em, itemsep=1pt]
    \item \textbf{SMAPP NYU elections (SMAPPNYU):}
          \url{github.com/SMAPPNYU/twitter_elections_public_interest}.
          \textbf{Rejected: tweet IDs only.} Rehydration requires X Basic API
          (\$100/month); not in budget. Same decision as echen102, Harvard
          2016, and VoterFraud 2020. Revisit if X API pricing changes or a
          pre-hydrated mirror becomes available.
    \item \textbf{UW ResearchWorks (b88ba40a):}
          \url{digital.lib.washington.edu/researchworks/items/}\allowbreak
          \texttt{b88ba40a-4001-4fb4-b960-2ea2e49d58a0}.
          \textbf{Rejected: tweet IDs only.} Rehydration requires X Basic API
          (\$100/month); not in budget. Same decision as SMAPP NYU, echen102,
          Harvard 2016, and VoterFraud 2020. Revisit if X API pricing changes
          or a pre-hydrated mirror becomes available.
    \item \textbf{X 2024 election (sinking8):}
          \url{github.com/sinking8/x-24-us-election}.
          Visit repo and check: if full tweet objects, save to
          \texttt{/Volumes/JUHIDRIVE/electoralData/archives/twitter/x\_2024/}
          and add to Tier~1.
          If IDs only, skip (not usable without paid API).
    \item \textbf{US Pres Elections 2020 (echen102):}
          \url{github.com/echen102/us-pres-elections-2020}.
          Same check. Emily Chen's datasets are typically IDs --- but
          check for any full-object sample files in the repo.
  \end{itemize}

  \smallskip
  \textbf{Bot filter --- blocklist only:}
  \begin{itemize}[leftmargin=1.6em, itemsep=1pt]
    \item \textbf{Political Twitter bots (khanradcoder):}
          Kaggle: \textit{khanradcoder/political-tweets-from-twitter-bots}.
          Extract bot user IDs into
          \texttt{/Volumes/JUHIDRIVE/electoralData/archives/validation/}\allowbreak
          \texttt{bot\_blocklist/ids.txt}.
          \textbf{Do not use as training data.}
  \end{itemize}

  \smallskip
  \textbf{News archives --- historical article text:}
  \begin{itemize}[leftmargin=1.6em, itemsep=1pt]
    \item \textbf{3DLNews2 (W\&M NewsLab) --- Tier 1, primary:}
          \url{github.com/wm-newslab/3DLNews}.
          8M+ URLs from 14,000 local outlets, 1995--2024.
          \textbf{Do not fetch HTML yet.} Download the URL metadata file
          only and save to \texttt{/Volumes/JUHIDRIVE/electoralData/archives/}\allowbreak
          \texttt{news/3dlnews/}. Filter
          by date window and state before fetching --- this reduces the
          fetch set from millions to tens of thousands of URLs. Run the
          HTML fetch-and-parse as a SLURM array job on HPC scratch (see
          the storage strategy box above). Tag with \texttt{source="3dlnews"}.
    \item \textbf{Webhose free news dataset --- Tier 2, secondary:}
          \url{github.com/Webhose/free-news-datasets}.
          Full article text pre-extracted, multi-outlet. Save to
          \texttt{/Volumes/JUHIDRIVE/electoralData/archives/news/webhose/}.
          Tag with \texttt{source="webhose"}.
  \end{itemize}

  \smallskip
  \textbf{Reddit archives --- Watchful1 Pushshift dump (AcademicTorrents):}
  \begin{itemize}[leftmargin=1.6em, itemsep=1pt]
    \item \textbf{Source:} Watchful1 Pushshift dump on AcademicTorrents
          (hash \texttt{3e3f64dee22dc304cdd2546254ca1f8e8ae542b4}).
          Contains the top 40,000 subreddits from 2005-06 to 2025-12 (inclusive) as
          zstandard-compressed ndjson files. Total torrent is 3.97TB but
          supports \textbf{selective download by subreddit} --- only
          download the identity subreddits needed for bloc proxies.
    \item \textbf{Download only these files:}
          \texttt{Catholicism\_comments.zst},
          \texttt{Catholicism\_submissions.zst},
          \texttt{BlackPeopleTwitter\_comments.zst},
          \texttt{BlackPeopleTwitter\_submissions.zst},
          \texttt{LatinoPeopleTwitter\_comments.zst},
          \texttt{LatinoPeopleTwitter\_submissions.zst},
          \texttt{Conservative\_comments.zst},
          \texttt{Conservative\_submissions.zst},
          \texttt{Republican\_comments.zst},
          \texttt{Republican\_submissions.zst},
          \texttt{politics\_comments.zst},
          \texttt{politics\_submissions.zst},
          \texttt{progressive\_comments.zst},
          \texttt{progressive\_submissions.zst},
          \texttt{Jewish\_comments.zst},
          \texttt{Jewish\_submissions.zst},
          \texttt{Judaism\_comments.zst},
          \texttt{Judaism\_submissions.zst},
          \texttt{islam\_comments.zst},
          \texttt{islam\_submissions.zst},
          \texttt{Muslim\_comments.zst},
          \texttt{Muslim\_submissions.zst},
          \texttt{exchristian\_comments.zst},
          \texttt{exchristian\_submissions.zst}.
    \item \textbf{Storage:} save to
          \texttt{/Volumes/JUHIDRIVE/electoralData/archives/reddit/\{subreddit\_name\}/}.
    \item \textbf{Parsing:} use the \texttt{zstandard} library
          (\texttt{pip install zstandard}). Example parsing scripts at
          \url{github.com/Watchful1/PushshiftDumps}.
    \item \textbf{No user bio field on Reddit} --- apply subreddit-level
          demographic proxy per the devplan mapping table.
          Set \texttt{inference\_method="subreddit\_proxy"} on all posts
          and exclude from $\Sigma_\Delta$ covariance estimation.
    \item \textbf{Single-torrent coverage:} this dump covers all shocks in
          the registry from 2005 to 2025, so no per-shock Reddit downloads
          are needed.
    \item \textbf{Additional subreddits --- African American bloc:}
          \texttt{BlackLivesMatter} (direct BLM shock signal),
          \texttt{askblackpeople},
          \texttt{Blackpeople},
          \texttt{BlackPeopleofReddit},
          \texttt{blackmen},
          \texttt{BlackMentalHealth}.
          \textbf{NSFW check required before processing:}
          \texttt{BlackIsBetter} --- inspect content before including;
          skip if majority off-topic.
    \item \textbf{Additional subreddits --- Mormon bloc (new):}
          \texttt{mormon},
          \texttt{mormonpolitics}.
          Apply Mormon demographic proxy. \textbf{Note:} Mormons are a
          distinct bloc from Evangelical Protestant --- they trend
          Republican but broke partially from Trump in 2016 and 2020.
          Save to \texttt{/Volumes/JUHIDRIVE/electoralData/archives/reddit/mormon/}.
    \item \textbf{Additional subreddits --- Secular/progressive bloc:}
          \texttt{Liberal},
          \texttt{AskALiberal},
          \texttt{LibertarianLeft},
          \texttt{ProgressiveHQ},
          \texttt{atheism}.
    \item \textbf{Additional subreddits --- Deconversion signal:}
          \texttt{exlldm} (ex La Luz del Mundo --- Latino religious
          deconversion signal),
          \texttt{exmormon} (now complete with Mormon additions above).
    \item \textbf{Do not process --- accidental downloads (irrelevant):}
          \texttt{askajudge},
          \texttt{askanelectrician},
          \texttt{aternos},
          \texttt{atheismbot},
          \texttt{atheismindia},
          \texttt{ateismo\_br},
          \texttt{morganbaiiley},
          \texttt{morningsomewhere},
          \texttt{moroccancuties1},
          \texttt{morrissey},
          \texttt{mortarsportfeed},
          \texttt{mossberg},
          \texttt{mostre\_os\_peitos},
          \texttt{mother4},
          \texttt{motherbussnark},
          \texttt{RefugeOfSilksong},
          \texttt{RegalUnlimited},
          \texttt{Referrals},
          \texttt{atc2}.
          Delete these files after confirming they are not needed.
    \item \textbf{Incomplete downloads --- re-download from torrent:}
          \texttt{politics\_comments.zst},
          \texttt{Conservative\_submissions.zst},
          \texttt{SoCalLatina\_comments.zst},
          \texttt{DebateACatholic\_submissions.zst},
          \texttt{excatholic\_comments.zst}.
          Verify file size against the torrent manifest before processing.
  \end{itemize}

  {\sloppy For each dataset, confirm presence of \texttt{text} (or
  \texttt{full\_text}), \texttt{created\_at}, and \texttt{lang}.
  Record all findings in the archives README before Day 2.\par}

  \smallskip
  \textbf{Resolution (2026-06-03):} Schema verified for all downloaded archives.
  Full findings in
  \texttt{/Volumes/JUHIDRIVE/electoralData/archives/README.md}.
  Summary:

  \begin{itemize}[leftmargin=1.6em, itemsep=2pt]
    \item \textbf{\texttt{election\_2012} (\texttt{tweets.json}):} Tier~1.
          Fields: \texttt{text}, \texttt{created\_at} (Twitter date string),
          \texttt{user.description} confirmed present.
          \textbf{No \texttt{lang} field} --- hardcode \texttt{"en"}.

    \item \textbf{\texttt{election\_2016} (\texttt{election\_day\_tweets.csv}):} Tier~1.
          Fields: \texttt{text}, \texttt{created\_at}, \texttt{lang},
          \texttt{user.description} (dot-notation CSV column) --- all confirmed.

    \item \textbf{\texttt{election\_2020} (manchunhui Kaggle):} Tier~1.
          Two files (\texttt{hashtag\_donaldtrump.csv}, \texttt{hashtag\_joebiden.csv}).
          Fields: \texttt{tweet}$\to$text, \texttt{created\_at}, \texttt{user\_description}
          confirmed present.
          \textbf{No \texttt{lang} field} --- infer with langdetect or hardcode \texttt{"en"}.
          Deduplicate by tweet\_id across both files before sampling.

    \item \textbf{\texttt{x\_2024} (sinking8 / USC):} \textbf{Tier~1 confirmed ---
          full objects, NOT IDs-only.} CSV.gz, snscrape format, $\sim$47~parts
          $\times$ $\sim$1M tweets each.
          Fields: \texttt{rawContent}$\to$text, \texttt{date}$\to$created\_at,
          \texttt{lang}, \texttt{user.rawDescription}$\to$bio (confirmed populated).
          \textbf{Note:} \texttt{user} column is a Python dict repr --- parse with
          \texttt{ast.literal\_eval()} to extract \texttt{rawDescription}.
          Filter by \texttt{date} window per shock slug.

    \item \textbf{\texttt{kashish\_usc} (\texttt{truth\_social\_2024}):} Tier~2.
          1,585,187 rows confirmed. Fields: \texttt{status}$\to$text,
          \texttt{timestamp}$\to$created\_at, \texttt{author\_username}.
          No bio, no lang column. Platform proxy: Evangelical/MAGA/conservative Protestant.

    \item \textbf{\texttt{telegram\_2024}:} Tier~2.
          Fields: \texttt{content}$\to$text, \texttt{date}$\to$created\_at,
          \texttt{language}$\to$lang, \texttt{from\_id}$\to$author.
          Platform proxy: Evangelical/conservative Protestant.

    \item \textbf{TikTok (\texttt{tiktok/}):} Tier~2 with Whisper transcripts.
          \texttt{TikTok\_IDs\_and\_Transcripts\_Published\_2025\_Feb.csv} has
          single column \texttt{whisper\_voice\_to\_text} ($\sim$53\% coverage).
          No \texttt{created\_at}, no \texttt{lang}, no bio.
          Use non-empty transcript rows only; hardcode \texttt{"en"}.

    \item \textbf{\texttt{blm\_2020} (alenjose): \textcolor{red}{DROPPED.}}
          All 100,000 tweets are from 2023-04-03, not May--June~2020.
          TweetBLM.csv and BLM2020\_Dataset.csv also rejected --- same issues
          (IDs-only, wrong language, or wrong time period).
          \textbf{Decision (2026-06-03):} BLM 2020 training signal comes entirely
          from Reddit Pushshift \texttt{r/BlackLivesMatter} and
          \texttt{r/BlackPeopleTwitter} filtered to May~25--August~31, 2020.
          \texttt{twitter/blm\_2020/} removed from active archive structure.

    \item \textbf{Reddit Pushshift:} all targeted subreddits downloaded.
          \texttt{DebateAChristian\_comments.zst} and
          \texttt{DebateCommunism\_comments.zst} are \texttt{.part} files ---
          re-download needed.
          \texttt{politics\_comments.zst} (26~GB) and
          \texttt{Conservative\_submissions.zst} (410~MB) are full-size and
          assumed complete.

    \item \textbf{3DLNews2:} \textbf{still downloading} as of 2026-06-03.
          Do not attempt HTML fetch until download completes.
  \end{itemize}}{5}

\done \task{Write \texttt{scripts/sample\_archives.py} and submit it as a
  \textbf{HPC SLURM array job} to draw stratified samples from the raw archives.
  This runs \emph{before} any RoBERTa scoring and reduces the corpus from
  hundreds of millions of posts to a tractable, representative sample.

  \textbf{Three strata per archive per shock event:}
  \begin{enumerate}[leftmargin=1.6em, itemsep=1pt]
    \item \textbf{Temporal:} sample proportionally by day across the shock
          window (e.g.\ 3 days before, day-of, 3 days after, 7 days after)
          so the full arc of public reaction is captured rather than just
          the peak day.
    \item \textbf{Demographic:} after a fast keyword-based bio pass, sample
          proportionally by inferred bloc so all 10 groups are represented.
          Without this step, high-volume blocs (secular, evangelical) crowd out
          low-volume blocs (Jewish, Muslim).
    \item \textbf{Sentiment extremity:} run a lightweight keyword sentiment
          score (VADER or similar) and oversample posts at the extremes
          (top and bottom 20\% of the distribution). Extreme examples are
          more informative for fine-tuning than neutral ones.
  \end{enumerate}

  \textbf{Target output:} $\sim$200,000 posts per major archive dataset,
  $\sim$5,000 per shock event per platform. Submit as a SLURM array job
  (\texttt{\#SBATCH --array=0-N} where N is the number of archive/shock
  combinations). Add \texttt{just sample} to the Justfile.

  \textbf{Status (2026-06-09):} Sampling complete on USC Laguna HPC ---
  \textbf{1,553,246 posts} across \textbf{77 shock/archive combinations},
  submitted as SLURM array job
  (\texttt{scripts/hpc/sample\_archives\_laguna.slurm}).
  All three strata implemented: temporal (proportional by day across shock
  window), demographic (keyword bio pass with religion$\to$race$\to$gender
  priority, CNN~2024 NEP weights), sentiment extremity (VADER~2$\times$
  weight on $|\text{compound}| \geq 0.5$).
  Target of 200,000 posts per major archive met for
  \texttt{election\_2012} (200,000), \texttt{election\_2020} (113,659),
  \texttt{blm\_george\_floyd\_2020} (74,938).
  Reddit Pushshift covers 2013--2025 shocks;
  \texttt{reddit\_monthly} covers pre-2013 shocks pending
  \texttt{filter\_reddit\_monthly} completion.
  3DLNews local news parsed for 16 states, 2015--2024.
  \texttt{just sample} and \texttt{just sample-laguna} added to Justfile.

  \textbf{Two pending items before fully complete:}
  \begin{enumerate}[leftmargin=1.6em, itemsep=1pt]
    \item \texttt{filter\_reddit\_monthly\_laguna.slurm} submitted but output
          not yet sampled --- will add $\sim$200k posts for pre-2013 shocks
          once filtering completes.
    \item Discord extraction jobs on Laguna not yet run.
  \end{enumerate}}{3}

\done \task{Write \texttt{scripts/clean\_with\_llm.py} and run it on the
  \textbf{M5} using a local open-weight model (Mistral-7B-Instruct)
  (free tier: 1,000 requests/day) during the initial build phase. In
  continuous mode post-SRP, this script runs on the \textbf{Intel Mac}
  as part of the nightly launchd pipeline. Four cleaning operations:

  \begin{enumerate}[leftmargin=1.6em, itemsep=1pt]
    \item \textbf{Off-topic detection.} For each post, prompt Gemini:
          \textit{``Is this post about [shock event]? Reply yes or no.''}
          Drop all ``no'' responses. This removes posts that mention the
          shock keyword coincidentally (e.g.\ MeToo off-topic tweets
          such as entertainment industry gossip unrelated to political activation).
    \item \textbf{Spam and low-information filtering.} Remove posts that
          survive the bot blocklist but add no signal: pure retweet text,
          copy-paste chain tweets, emoji-only posts, bare URL shares with
          no commentary.
    \item \textbf{Text normalisation.} Expand political abbreviations
          (SCOTUS, POTUS, MAGA, FLOTUS), normalise hashtag formatting,
          strip URLs and \texttt{@mentions} so RoBERTa receives cleaner
          input without domain-specific noise.
    \item \textbf{Cross-platform deduplication.} The same post sometimes
          appears in multiple archives. Deduplicate on
          \texttt{hash(user\_id + text[:50])} before scoring.
  \end{enumerate}

  \textbf{Throughput:} $\sim$50 posts per Gemini prompt $\times$ 1,000 requests/day
  = 50,000 posts/day. For 200k posts per archive this takes 4 days.
  \textbf{Start the most important archives (COVID pandemic tweets from Kaggle (gpreda), COVID vaccine tweets from Kaggle (kaushiksuresh147), Reddit BLM 2020, 2020 election) during
  Weeks 2--3} while other components are being built, so cleaned data is
  ready for the Week~4 RoBERTa scorer. Add \texttt{just clean} to the
  Justfile. Log the discard rate per cleaning step in
  \texttt{/Volumes/JUHIDRIVE/electoralData/archives/README.md}.

  \textbf{Status (2026-06-09):} Cleaning complete ---
  \textbf{1,553,246} sampled posts $\to$ \textbf{231,725} cleaned posts
  (\textbf{14.9\%} kept) across \textbf{77 shock/archive combinations}.
  Off-topic filter via Gemini~2.5 Flash Lite with async concurrency
  (50 concurrent batches, \texttt{asyncio.gather}). Discard breakdown:
  $\sim$85\% off-topic, ${<}1\%$ spam, ${<}1\%$ dedup. Cleaned data at
  \texttt{/scratch/JDamley28@cmc.edu/electoralData/cleaned/} on Laguna
  and synced to \texttt{/Volumes/JUHIDRIVE/electoralData/cleaned/}.}{3}

\done \task{Implement \texttt{electoral/nlp/scraper.py} on the
  \textbf{Intel Mac}: \texttt{scrape\_articles(shock\_id,
  sources, date\_range)} for news sources (Christianity Today, CBN, Univision,
  NYT, WaPo, Fox). Strip HTML, extract plain text, save to the Syncthing
  \texttt{rawdata/articles/} folder so the M5 can read results immediately.
  Enforce 100-word minimum, log discard rate, test on Ayatollah shock.

  \smallskip\textit{Status (2026-06-10):} Implemented as \texttt{electoral/nlp/scraper.py} using direct RSS feeds (BBC, Guardian, NPR, Christianity Today, CBN, Fox News). Google News RSS dropped due to URL encoding issues. Tested on \texttt{iran\_war\_2026} shock---4 articles written from BBC and Guardian. Scraper is live-collection only; RSS feeds do not archive past articles so historical shocks cannot be scraped. \texttt{just scrape} added to Justfile.}{1}

\done \task{Implement \texttt{electoral/nlp/news\_loader.py}:
  \texttt{NewsArchiveLoader} that normalises pre-collected news articles
  into the same \texttt{articles.jsonl} schema as the live scraper,
  so the RoBERTa scorer is blind to source.

  \textbf{Two sources, handled as subclasses or branches:}
  \begin{itemize}[leftmargin=1.6em, itemsep=1pt]
    \item \textbf{\texttt{3DLNewsLoader}:} reads URL entries from
          \texttt{/Volumes/JUHIDRIVE/electoralData/archives/news/3dlnews/};
          filters to the shock
          date window using \texttt{publication\_date}; fetches and
          HTML-strips each article (same logic as \texttt{scraper.py});
          optionally filters by \texttt{state} or outlet name for
          targeted demographic coverage; enforces 100-word minimum;
          tags \texttt{source="3dlnews"}. This is the primary historical
          news source --- 14,000 local outlets covering the communities
          that constitute each of the 10 blocs.
    \item \textbf{\texttt{WebhoseLoader}:} reads pre-extracted text
          from \texttt{/Volumes/JUHIDRIVE/electoralData/archives/news/webhose/};
          filters by
          \texttt{published} date and optionally by \texttt{url} domain;
          enforces 100-word minimum; tags \texttt{source="webhose"}.
          No HTML fetch needed.
  \end{itemize}

  Both write to \texttt{rawdata/articles/\{shock\_id\}/} with fields
  \texttt{\{id, text, published, url, source, archive\_id\}}.
  Test \texttt{3DLNewsLoader} on the \texttt{metoo\_2017\_2023} shock, filtering
  to Oct--Nov 2017 articles from Pennsylvania, Georgia, and Texas
  local outlets as a demographic cross-section.

  \smallskip\textit{Status (2026-06-10):} Implemented as \texttt{electoral/nlp/news\_loader.py}
  with \texttt{ScrapedNewsLoader}, \texttt{3DLNewsLoader}, and \texttt{WebhoseLoader}.
  Tested \texttt{3DLNewsLoader} on \texttt{metoo\_2017} shock filtered to PA, GA, TX
  states Oct--Nov 2017: loaded 16,095 articles (51,303 skipped by date, 55,853 by
  state, 4,794 under 100 words). Outlet-level demographic proxies embedded in
  \texttt{author\_description}. Module is a library---imported by
  \texttt{score\_roberta.py} and other pipeline scripts.}{2}
\end{taskbox}

% ── Day 2 ────────────────────────────────────────────────────────────────────
\dayhead{2}{Tue --- Social Media Collectors \& Bio Keyword Lexicon}

\begin{taskbox}
\done \task{Implement \texttt{electoral/nlp/social.py} with a
  \texttt{SocialCollector} abstract base class and the following concrete
  subclasses. All write identically-formatted \texttt{posts.jsonl};
  the scorer is blind to source.

  \begin{itemize}[leftmargin=1.6em, itemsep=1pt]
    \item \texttt{BlueSkyCollector}: primary free live path. AT Protocol
          \texttt{app.bsky.feed.searchPosts} via \texttt{atproto}. Returns
          author \texttt{description} and \texttt{language} fields for bio
          classification. No rate limit on public posts.
    \item \texttt{ApifyCollector}: secondary free live path.
          \texttt{apify-client} Python SDK running the X Scraper actor.
          Covers the Evangelical, Catholic, and urban demographic profile
          that Bluesky underrepresents. Free tier: 500 results/run.
          \textbf{Do not implement \texttt{XCollector} or install
          \texttt{tweepy} --- the X Basic API is not used.}
    \item \texttt{TruthSocialCollector}: Mastodon-compatible
          \texttt{/api/v1/timelines/public?q=} endpoint. No auth. Small
          user base; weight accordingly.
    \item \texttt{FacebookCollector}: Meta Content Library client if
          approved, else public page reaction-count fallback.
  \end{itemize}
  For the Ayatollah pilot: collect at least 1,000 posts from Bluesky
  and run Apify for the same event to confirm schema parity.

  \smallskip\textit{Status (2026-06-10):} Implemented as \texttt{electoral/nlp/social.py}
  with \texttt{SocialCollector} ABC and \texttt{BlueSkyCollector},
  \texttt{ApifyCollector}, \texttt{TruthSocialCollector}, \texttt{FacebookCollector}
  subclasses. \texttt{BlueSkyCollector} tested on \texttt{ayatollah\_assassination}
  and \texttt{iran\_war\_2026} shocks---317 posts in 60 seconds, schema confirmed
  correct. \texttt{author\_description} null until Pi bio server classifies bios in
  scoring stage. Apify schema parity test pending Apify token configuration.
  TruthSocial and Facebook collectors implemented but untested.}{3}

\done \task{Implement \texttt{electoral/nlp/archive.py}:
  \texttt{HistoricalArchiveLoader} class with
  \texttt{load(shock\_id, archive\_id, window\_hours=72,
  bot\_blocklist=None)}. Core behaviour:
  (i) reads raw archive files from
  \texttt{/Volumes/JUHIDRIVE/electoralData/archives/\{platform\}/\{archive\_id\}/};
  (ii) filters to the shock date window;
  (iii) drops posts whose \texttt{user.id} is in \texttt{bot\_blocklist}
  and logs the count;
  (iv) normalises to \texttt{\{id, text, created\_at, lang,
  bloc\_weights, source, archive\_id\}};
  (v) calls \texttt{/classify\_batch} for posts with \texttt{user.description};
  (vi) applies platform proxy for posts without a bio field;
  (vii) saves to \texttt{rawdata/social/archive/\{shock\_id\}/}.

  Implement two additional loaders as subclasses or standalone functions:

  \begin{itemize}[leftmargin=1.6em, itemsep=1pt]
    \item \textbf{\texttt{TelegramLoader}:} normalises Telegram message
          objects (field names differ from Twitter --- \texttt{message}
          instead of \texttt{text}, \texttt{date} instead of
          \texttt{created\_at}, no \texttt{user.description}).
          Applies Evangelical/conservative Protestant as the
          platform-level proxy for all posts. Tags with
          \texttt{platform="telegram"}.
    \item \textbf{\texttt{TikTokLoader}:} normalises TikTok video
          metadata (caption text as \texttt{text}, \texttt{createTime}
          as \texttt{created\_at}, no bio field). Applies
          younger/secular/diverse as the platform-level proxy.
          Tags with \texttt{platform="tiktok"}.
  \end{itemize}

  Handle Twitter schema differences: the MeToo dataset uses standard Twitter v2
  JSON \texttt{text} field; normalise all legacy archives to \texttt{text}.
  Test the main loader on \texttt{metoo\_2017\_2023} filtered to Oct 15, 2017
  $\pm$14 days (MeToo origin week).
  Test \texttt{TelegramLoader} on the USC 2024 dataset.
  Verify all outputs match the \texttt{posts.jsonl} schema exactly.

  \smallskip\textit{Status (2026-06-10):} Implemented as \texttt{electoral/nlp/archive.py}.
  \texttt{HistoricalArchiveLoader} tested on \texttt{metoo\_2017/metoo\_liwc}---SAV
  support added, loads correctly. \texttt{TelegramLoader} tested on
  \texttt{election\_2024/telegram\_2024}---3,122 posts loaded from CSV format.
  \texttt{TikTokLoader} implemented but untested (no TikTok data available).
  All outputs match \texttt{posts.jsonl} schema.}{4}

\done \task{Build the keyword lexicon in
  \texttt{electoral/nlp/bio\_classifier.py}: a \texttt{BIO\_LEXICON} dict
  mapping each canonical \texttt{BlocId} to a list of bio keyword patterns.
  Source the patterns from Prof.\ Espinosa's published research on linguistic
  and cultural markers per bloc. Required entries per bloc:

  \smallskip
  \begin{itemize}[leftmargin=1.6em, itemsep=1pt]
    \item \texttt{african\_american}: ``Black'', ``African American'', ``HBCU'',
          ``\#BlackTwitter'', ``NAACP''
    \item \texttt{latino}: ``Latino'', ``Latina'', ``Latinx'', ``Chicano'',
          ``Hispanic'', bio in Spanish
    \item \texttt{asian}: ``Asian American'', ``AAPI'', bio in Chinese/Korean/
          Vietnamese/Tagalog/Hindi
    \item \texttt{secular}: ``atheist'', ``agnostic'', ``humanist'',
          ``no religion'', ``nonreligious''
    \item \texttt{jewish}: ``Jewish'', ``\#JewishTwitter'', ``synagogue'',
          ``rabbi'', bio in Hebrew
    \item \texttt{muslim}: ``Muslim'', ``Islam'', ``Alhamdulillah'',
          ``Inshallah'', ``mosque'', bio in Arabic
    \item \texttt{catholic}: ``Catholic'', ``parish'', ``rosary'',
          ``Vatican'', ``\#CatholicTwitter''
    \item \texttt{protestant}: ``Protestant'', ``Lutheran'', ``Methodist'',
          ``Presbyterian'', ``Baptist'' (non-Evangelical)
    \item \texttt{evangelical}: ``Evangelical'', ``born again'',
          ``Southern Baptist'', ``\#ChristianTwitter'', ``Bible-believing''
  \end{itemize}
  \smallskip
  The bio classifier outputs \textbf{three separate probability distributions}
  per post (one per stratum):
  \texttt{race\_weights} over
  \texttt{\{african\_american, latino, asian, white, other\_race\}},
  \texttt{religion\_weights} over
  \texttt{\{evangelical, catholic, protestant, secular, jewish, muslim, other\_rel\}},
  and \texttt{gender\_weights} over \texttt{\{women, men, other\_gender\}}.
  These are three independent distributions; no cross-tabulation is produced.

  \textbf{Gender signal} (\texttt{gender\_signal: "F" | "M" | null}) is
  extracted as a third field from pronouns (``she/her'', ``he/him''),
  explicit words (``woman/women'', ``man/men'', ``mom'', ``feminist''), and
  name-based inference as a last resort.

  Store as: \texttt{configs/race\_lexicon.json},
  \texttt{configs/religion\_lexicon.json},
  \texttt{configs/gender\_lexicon.json}, all loaded at runtime.}{3}

\done \task{\textbf{Language fallback --- validation-only, exclude from all estimation.}
  When both keyword matching and SetFit return no signal (bio empty or below
  confidence threshold), fall back to language-based Pew priors (e.g.\
  \texttt{"es"} $\to$ \{latino/catholic: 0.38, latino/protestant: 0.22, \ldots\}).
  Tag every post assigned via this fallback with
  \texttt{inference\_method: "language\_prior"} in the post record.

  \textbf{Critical:} do \emph{not} include language-prior posts in mean $\mu$
  or covariance $\Sigma_\Delta$ estimation. Including them in mean estimation
  while excluding from covariance bifurcates the data generating process ---
  $\mu$ and $\Sigma$ would come from different corpora, violating the
  mean-variance optimizer's assumptions. Use language-prior posts only as a
  held-out validation set: after estimation, compare their sentiment signal
  against the bio-classified signal and report any divergence in the paper.
  Store language priors in \texttt{configs/language\_priors.json}.}{2}

\done \task{Implement \texttt{TruthSocialCollector} and \texttt{FacebookCollector}
  in \texttt{social.py} (\textbf{not stubs} --- see hard fallback checklist above),
  and \texttt{TelegramLoader} and \texttt{TikTokLoader} in \texttt{archive.py}.
  \texttt{TruthSocialCollector} uses the Mastodon-compatible
  \texttt{/api/v2/search} endpoint at \texttt{truthsocial.com} (no auth).
  \texttt{FacebookCollector} must support \texttt{mode="content\_library"}
  and \texttt{mode="reaction\_fallback"} before this task is marked done.
  Truth Social, Facebook, Telegram, and TikTok posts do not provide structured
  author bio fields --- for all of these apply the platform-level demographic
  proxy directly (\texttt{inference\_method="platform\_proxy"}).

  \smallskip\textit{Status (2026-06-10):} All four classes implemented.
  \texttt{TruthSocialCollector} tested---returns 403 Forbidden on
  unauthenticated requests (Truth Social has tightened API access since devplan
  was written). Fails gracefully with 0 posts returned, no crash.
  \texttt{FacebookCollector} implemented with both \texttt{content\_library} and
  \texttt{reaction\_fallback} modes. \texttt{TelegramLoader} tested and working
  (3,122 posts for \texttt{election\_2024}). \texttt{TikTokLoader} implemented
  but no TikTok data available to test.
  \texttt{inference\_method=platform\_proxy} applied correctly to all platforms
  without bio fields.}{3}
\end{taskbox}

% ── Day 3 ────────────────────────────────────────────────────────────────────
\dayhead{3}{Wed --- Bio Classifier: Embedding Model, Language Signal \& Scorer}

\begin{taskbox}
\done \task{Build the manual bio labelling dataset: create
  \texttt{data/bio\_labels/labeled\_bios.jsonl} with 100--200 real user
  bios per bloc (aim for 1,000+ total). Source bios from the polarization
  corpus (Kiran et al., already downloaded in Day~1 --- ideologically-sorted
  user bios with labels), the Kaggle political tweets datasets, and the
  senator-tweets corpus (m-newhauser/senator-tweets, already downloaded). For any blocs underrepresented in the archives,
  manually write example bios from public Bluesky profiles. Bios matching the
  keyword lexicon from Day~2 can be auto-labelled; focus manual effort on
  ambiguous cases.

  \smallskip\textit{Status (2026-06-10):} \texttt{data/bio\_labels/labeled\_bios.jsonl}
  created with 2,191 labeled bios. Distribution: 150 per race bloc (white, latino,
  african\_american, asian, other\_race), 150 per religion bloc (catholic, protestant,
  secular, muslim, other\_rel, jewish, evangelical), 150 per gender signal (M, F).
  All labeled via \texttt{keyword\_auto} from \texttt{election\_2020} and other corpus
  bios. No manually labeled ambiguous cases yet---acceptable for SetFit pre-training
  since keyword-matched bios are high confidence. Manual labeling of ambiguous cases
  deferred to post-SRP phase.}{2}

\done \task{Train the \textbf{SetFit} bio classifier (Option~B): install
  \texttt{setfit}. Load \texttt{all-MiniLM-L6-v2} as the embedding backbone.
  Run \texttt{SetFitTrainer(model, train\_dataset, eval\_dataset)} with the
  labelled bios; target $> 0.75$ macro-F1 across all 10 blocs. Export the
  trained head weights.

  \textbf{Update \texttt{scripts/pi\_bio\_server.py} to handle both
  execution modes explicitly.} On startup, read
  \texttt{configs/base.json["pi\_npu\_enabled"]}:
  \begin{itemize}[leftmargin=1.6em, itemsep=1pt]
    \item \textbf{If \texttt{true}:} load the embedding backbone from the
          compiled HEF via the Hailo SDK. SetFit classification head runs
          on CPU as a second step.
    \item \textbf{If \texttt{false}:} load \texttt{all-MiniLM-L6-v2} via
          HuggingFace \texttt{SentenceTransformer} into
          CPU memory. Do not attempt to initialise the Hailo SDK or load any
          HEF. The SetFit classification head runs on CPU as normal.
  \end{itemize}
  Never attempt to pass data to the Hailo interface when NPU is disabled ---
  this causes an immediate crash. The \texttt{/health} endpoint must return a
  JSON object with \texttt{"status":"ok"} and \texttt{"mode"} set to either
  \texttt{"npu"} or \texttt{"cpu"}, so callers can confirm which path is active.
  Test both modes explicitly before Week 4 Day 3 begins.}{3}

  \textit{Status (2026-06-07): server infrastructure and NPU/CPU path separation
  complete (\texttt{scripts/pi\_bio\_server.py} rewritten). Training script written
  (\texttt{scripts/train\_setfit.py}). 28 bio-classifier tests pass. Prerequisite
  blocking actual model training: \texttt{data/bio\_labels/labeled\_bios.jsonl}
  does not yet exist --- run \texttt{build\_bio\_labels.py} on M5 with JUHIDRIVE
  mounted to generate it, then \texttt{train\_setfit.py}. macro-F1 target of 0.75
  not yet measured (no trained weights exist).}

\done \task{Implement \texttt{electoral/nlp/scorer.py} with \textbf{embedding
  caching} (Option~H): \texttt{score\_text(text)} first checks
  \texttt{data/embeddings/\{post\_id\}.parquet} for a cached embedding vector.
  If found, skip the model forward pass and use the cached vector directly.
  If not found, run the model, save the embedding to Parquet, then score.
  Implement \texttt{score\_posts\_weighted(posts)} using the SetFit bio weights.
  The first run of the scorer is slow (embeds everything); all subsequent runs
  with different aggregation logic are near-instant. Log cache hit rate.}{3}
\end{taskbox}

% ── Day 4 ────────────────────────────────────────────────────────────────────
\dayhead{4}{Thu --- Lagged Favorability Correlation \& Elasticity Regression}

\begin{taskbox}
\done \task{Generate \textbf{synthetic training data} using Gemini 2.5 Pro
  (Option~D). Run \texttt{scripts/generate\_synthetic.py}: (i) construct
  a single prompt containing the full voter panel (as CSV), all 25+ real shock
  records, the canonical stratum IDs, and instructions to generate 500--1{,}000
  plausible hypothetical shock scenarios with demographically grounded delta
  estimates per stratum; (ii) parse the JSON response into JSONL records matching
  the fine-tuning schema; (iii) mark each record \texttt{"synthetic": true}
  and assign an \texttt{mmd\_weight} per the formula below; (iv) save to
  \texttt{data/finetune/synthetic.jsonl} and diagnostics to
  \texttt{data/finetune/synthetic\_diagnostics.json}.

  \textbf{Three required diagnostics} (compute for every batch; store in diagnostics file):
  \begin{itemize}[leftmargin=1.4em, itemsep=1pt]
    \item \textbf{MMD}: compute RBF-kernel MMD between synthetic and real joint
          $\Delta\mu$ distributions. Set $\lambda = 1/(2\sigma^2)$ where
          $\sigma^2$ = median pairwise distance over the real data (median heuristic).
          Assign weight $w = 0.5 \times \exp(-\lambda \cdot \text{MMD}^2)$.
    \item \textbf{PCA alignment}: fit PCA on real $\Delta\mu$ matrix; project
          synthetic batch onto the top-2 PCs; flag if $< 70\%$ of synthetic
          variance is explained.
    \item \textbf{PCD (Pairwise Correlation Difference)}: compute Frobenius norm
          of (real pairwise correlation matrix $-$ synthetic pairwise correlation
          matrix) normalized by number of off-diagonal elements. Log for paper.
  \end{itemize}
  Never use a hard KS-test rejection gate --- the soft MMD weight is sufficient.

  \smallskip\textit{Status (2026-06-08):} \texttt{scripts/generate\_synthetic.py}
  complete. Generated 557 synthetic shock scenarios via Gemini 2.5 Pro covering all
  15 blocs with delta\_bin estimates. Saved to \texttt{data/finetune/synthetic.jsonl}.
  Diagnostics (MMD, PCA alignment, PCD) deferred---require $\geq$3 real scored shocks
  in \texttt{data/finetune/real.jsonl}. Will recompute after RoBERTa scorer runs on
  cleaned corpus.}{3}

\done \task{Build \texttt{configs/market\_contracts.json}: a static mapping
  of every shock event to its exact contract identifier on each supported
  prediction market platform. Querying Polymarket, Manifold, or Metaculus
  with a generic slug like \texttt{metoo\_2017\_2023} returns empty or
  irrelevant results --- each platform uses its own internal IDs that must
  be looked up manually. Each entry maps a shock slug to a dict of
  \texttt{\{platform: contract\_id\_or\_null\}}, for example:
  \texttt{"metoo\_2017\_2023": \{"polymarket": null, "predictit": null,
  "metaculus": "3712", "manifold": null, "iem": null\}}.
  For each of the 25+ shock events, visit each platform's search interface,
  find the closest relevant contract, and record the exact ID. Mark
  \texttt{null} where no contract exists. This file is the source of truth
  for the market collector --- no dynamic slug resolution. Commit to the repo
  and add to the replication package.}{2}

\done \task{Implement \textbf{prediction market collection} in
  \texttt{electoral/markets/collector.py}. Load contract identifiers from
  \texttt{configs/market\_contracts.json} (built in the previous task) ---
  never resolve slugs dynamically; use only the hardcoded IDs. For each
  of the 25+ shock events, collect the market-implied win probability for
  both the Democrat and Republican candidate contracts at $t - 24$h,
  $t + 1$h, $t + 24$h, and $t + 72$h. API endpoints:
  \begin{itemize}[leftmargin=1.6em, itemsep=1pt]
    \item \textbf{Polymarket} (2020+): \texttt{clob-api.polymarket.com} REST API.
          Free, no auth required. Query by CLOB market address from
          \texttt{market\_contracts.json}.
    \item \textbf{PredictIt} (2014+): \texttt{www.predictit.org/api/marketdata/all/}
          JSON feed. Free, no auth. Query by contract ID from
          \texttt{market\_contracts.json}. Download full historical CSV from
          \url{predictit.org/research}.
    \item \textbf{Metaculus} (2015+): \texttt{www.metaculus.com/api2/questions/}
          Free REST API. Query by question ID from \texttt{market\_contracts.json}.
    \item \textbf{Manifold} (2022+): \texttt{manifold.markets/api/v0/markets}
          Free REST API. Query by market slug from \texttt{market\_contracts.json}.
    \item \textbf{Iowa Electronic Markets} (1988+): Download historical CSV from
          \url{tippie.uiowa.edu/iem}. Parse offline. Best coverage for 2000--2014.
  \end{itemize}
  Implement \texttt{electoral/markets/aggregator.py}:
  \texttt{volume\_weighted\_price(market\_prices, volumes)} that weights each
  market's price by its trading volume; assign equal weight when volume is
  unavailable (Metaculus, Manifold). Construct a
  \texttt{PredictionMarketData} artifact per shock event and write to
  \texttt{artifacts/markets/\{shock\_id\}.json}. For events before 2014,
  set \texttt{prediction\_market = null} in the fine-tuning record.}{3}

\done \task{Implement \texttt{electoral/nlp/elasticity.py}:
  \texttt{fit\_elasticity(news\_scores, social\_scores, polling\_deltas)} that
  regresses observed vote-share deltas on both news RoBERTa scores and
  baseline-adjusted social media scores, using platform as a demographic proxy
  weight. Use \texttt{sklearn.linear\_model.Ridge} with cross-validation on
  held-out shock events. Store coefficients per source (news vs.\ each platform
  vs.\ archive) per bloc. For shocks with archive coverage, the archive posts
  contribute to the same \texttt{social\_scores} dict as live posts --- they
  are already normalised to the same schema and the scorer has processed them
  identically.}{3}

\done \task{Compute and log the per-platform correlation between baseline-adjusted
  social sentiment and 14-day lagged favorability deltas. Write results to
  \texttt{docs/platform\_correlations.md}. If a platform shows near-zero
  correlation, down-weight or exclude it in the regression --- document the
  decision.}{2}
\end{taskbox}

% ── Day 5 ────────────────────────────────────────────────────────────────────
\dayhead{5}{Fri --- Merge Stage, Dataset Assembly, Full Corpus Run \& PR}

\begin{taskbox}
\done \task{Implement \texttt{merge\_posts(shock\_id)} in
  \texttt{electoral/kernels/sentiment.py} as a Prefect task. This stage runs
  \emph{before} the RoBERTa scorer and is required because multiple files may
  exist per shock (one per platform per collector run). The task:
  (i) finds all \texttt{rawdata/social/\{platform\}/\{shock\_id\}/\*.jsonl}
  files across all platforms and all archive sources;
  (ii) concatenates them into a single
  \texttt{rawdata/merged/\{shock\_id\}/posts.jsonl};
  (iii) logs the total post count broken down by platform so you can see the
  coverage split.
  The scorer always reads from \texttt{rawdata/merged/} and is completely
  blind to how many files or platforms contributed. Wire this task into
  \texttt{pipeline.py} as the stage immediately preceding \texttt{score\_posts}.
  Verify: running the pipeline twice produces identical merged files.}{2}

\done \task{Implement \texttt{build\_finetune\_dataset} in \texttt{elasticity.py}:
  create one JSONL record per shock event using the unified schema. The
  \texttt{social\_roberta\_scores} field is keyed by \texttt{bloc\_id}; include
  the \texttt{prediction\_market} block from the \texttt{PredictionMarketData}
  artifact (set to \texttt{null} for pre-2014 events). Apply differential
  weighting: market-backed examples get \texttt{weight=1.5}, poll-only examples
  get \texttt{weight=1.0}, synthetic examples get \texttt{weight=0.5}. Write to
  \texttt{data/finetune/train.jsonl} (all cycles except 2020) and
  \texttt{data/finetune/eval.jsonl} (2020 only). Validate: no duplicate shock
  IDs, all 10 bloc keys present, no train/eval overlap.}{3}

\done \task{Submit the RoBERTa scoring pipeline as a \textbf{USC Laguna SLURM array job}
  across all 25+ shock events simultaneously.

  \textbf{First: confirm sampled and cleaned data is in place.} The scorer
  reads from \texttt{rawdata/merged/\{shock\_id\}/posts.jsonl} which should
  contain the output of \texttt{sample\_archives.py} and
  \texttt{clean\_with\_llm.py}, not the raw full-corpus files.
  Confirm post counts are in the 5,000--200,000 range per shock, not millions.
  If raw data was accidentally staged, re-run sampling first.

  \textbf{Then: transfer and submit.}
  Do not rely on Syncthing on USC Laguna nodes --- scratch nodes kill background
  daemons. Run:
  \begin{itemize}[leftmargin=1.6em, itemsep=1pt]
    \item \texttt{rsync -avz rawdata/merged/ laguna:\$SCRATCH/electoral/merged/}
    \item \texttt{rsync -avz rawdata/articles/ laguna:\$SCRATCH/electoral/articles/}
  \end{itemize}
  Wait for rsync to complete before submitting the job.

  \textbf{Then: write and submit} \texttt{scripts/score\_array.sh}:
  \texttt{\#SBATCH --array=0-24}, one node per shock event, each node
  loads its assigned \texttt{posts.jsonl} from
  \texttt{\$SCRATCH/electoral/merged/} and runs the weighted scorer.
  Collecting 25 events serially on the M5 would take 2--3 days; the array
  job finishes in under 2 hours on USC Laguna.
  Monitor with \texttt{squeue -u \$USER}.

  \textbf{Scorer sanity check --- LIWC validation (run after scoring completes):}
  \texttt{validation/metoo\_liwc/fulldataCHBR.sav} (Foster \& Rathlin 2022,
  DOI:~10.5683/SP3/1YWCW1) contains 3,683 LIWC-pre-scored tweets across four
  shock phases. After the RoBERTa scorer has processed the MeToo/Kavanaugh
  archives, load this file and compute the Spearman correlation between
  RoBERTa emotion outputs and the LIWC \texttt{Tone}, \texttt{posemo},
  \texttt{negemo}, and \texttt{anger} columns on the same tweet IDs.
  \textbf{If Spearman $\rho < 0.4$ on \texttt{negemo} or \texttt{anger}, do
  not proceed to SetFit training.} Flag the result in \texttt{DECISIONS.md}
  (column name, $\rho$ value, and the decision) before moving to Day~3.}{3}

\done \task{Write \texttt{tests/test\_bio\_classifier.py}: (i) keyword match
  for an explicit Evangelical bio returns \texttt{evangelical} weight $> 0.8$,
  (ii) Spanish-language \texttt{lang} code boosts \texttt{latino} weight,
  (iii) \texttt{combine\_signals} output sums to 1.0 for all inputs,
  (iv) unclassifiable bio returns platform fallback. Write
  \texttt{tests/test\_social.py}: (i) \texttt{BlueSkyCollector} records include
  \texttt{author.description} and \texttt{lang} fields, (ii) \texttt{ApifyCollector}
  output schema matches \texttt{BlueSkyCollector} exactly, (iii) baseline
  adjustment produces per-bloc scores in $[-1, 1]$, (iv) weighted scorer
  produces 10 distinct elasticity values not a single platform score.
  Write \texttt{tests/test\_sentiment.py}: (i) scorer outputs sum to 1.0,
  (ii) dataset has no train/eval overlap.}{1}

\done \task{Weekly PR: \texttt{feature/week4-roberta-social-pipeline}. Include:
  dataset statistics (examples, shock categories, mean elasticity per bloc),
  platform correlation table, and the Ayatollah pilot findings.}{1}

\done \task{Brief supervision check-in with Prof.\ Espinosa: share platform
  correlation results and elasticity scores for 3--4 key shocks. Get domain
  input on whether Truth Social's Evangelical reactions align with his
  published patterns. Juhi decides on platform weighting adjustments and
  documents in \texttt{DECISIONS.md}.}{2}
\end{taskbox}

\newpage

% ════════════════════════════════════════════════════════════════════════════
\section{Week 4 --- LLM Fine-Tuning \& Inference Endpoint (40 hrs)}
% ════════════════════════════════════════════════════════════════════════════

\textbf{Goal.} By Friday, Mistral 7B is fine-tuned via QLoRA on the unified
news + social media dataset on USC Laguna, the FastAPI endpoint returns valid
\texttt{ShockResponseData} JSON with constrained decoding, and the three-way
baseline comparison (news-only, social-only, unified) is documented.

\textbf{End condition.}
\begin{itemize}[leftmargin=1.6em, itemsep=1pt]
  \item \done 100 consecutive inference calls all return schema-valid JSON.
  \item \done Endpoint responds in under 3 seconds on M5.
  \item \done Three-way MAE comparison (news-only, social-only, unified) committed.
  \item \done PR merged.
\end{itemize}

% ── Day 1 ────────────────────────────────────────────────────────────────────
\dayhead{1}{Mon --- Schema Migration, USC Laguna Setup \& Local MLX Baseline}

\begin{taskbox}
\done \task{\textbf{Migrate \texttt{ShockResponseData} to the per-stratum schema
  before writing any LLM inference code.}
  The current \texttt{artifacts.py} stores a flat \texttt{deltas: dict[str,
  float]} over all 15 blocs with an $N \times N$ covariance ($N = 15$).
  This contradicts two architectural decisions in \texttt{DECISIONS.md}:
  (i)~the 5$\times$5 race-only covariance rule (religion and gender do not
  enter $\boldSigma_\Delta$), and (ii)~the requirement to persist the raw
  LLM 9-token bin output alongside the converted floats.

  Replace the current fields with:
  \begin{itemize}[leftmargin=1.6em, itemsep=1pt]
    \item \texttt{delta\_bins\_race}, \texttt{delta\_bins\_religion},
          \texttt{delta\_bins\_gender}: \texttt{dict[str, str]} storing the
          raw constrained-decoding token (one of the nine bin labels) per bloc.
          Required for fine-tuning reproducibility.
    \item \texttt{deltas\_race}, \texttt{deltas\_religion},
          \texttt{deltas\_gender}: \texttt{dict[str, float]} storing the
          converted $\Delta\mu$ per bloc.
    \item \texttt{delta\_eff}: \texttt{float} --- scalar effective delta
          $\lambda_1\sum w_i\Delta\mu_i^{\text{race}} + \ldots$
    \item \texttt{covariance}: \texttt{list[list[float]]} --- 5$\times$5
          race-bloc only. Validation: assert side $= 5$, not $N$.
  \end{itemize}

  Update \texttt{validate()} to assert the 5$\times$5 covariance shape and
  validate all three \texttt{delta\_bins\_*} dicts against the nine canonical
  bin labels. Update \texttt{build\_shock\_response} in \texttt{stages.py}
  and its tests in \texttt{test\_stages.py} to match the new schema.
  Do \emph{not} start USC Laguna setup until this migration is committed and all
  tests pass --- the LLM endpoint return type depends on this schema.}{2}

\done \task{Set up USC Laguna access: log in via SSH, confirm SLURM quota and
  available GPU nodes (\texttt{sinfo}), check scratch storage allocation
  (\texttt{quota}). Create the project directory at
  \texttt{\$SCRATCH/electoral-equilibrium/}. Document node names, partition
  names, and max walltime in \texttt{scripts/LAGUNA\_NOTES.md}.
  \emph{Note: USC Laguna is used for all training and compute; CMC HPC is
  reserved for hosting (vLLM deployment) only.}}{1}

\done \task{Copy the fine-tuning dataset to USC Laguna scratch: \texttt{scp -r
  data/finetune/ laguna:\$SCRATCH/electoral-equilibrium/data/}. Verify file
  counts match locally. Set up a Python virtual environment on Laguna:
  \texttt{python -m venv \$SCRATCH/venv} and install
  \texttt{transformers}, \texttt{peft}, \texttt{bitsandbytes}, \texttt{torch},
  \texttt{accelerate}. Confirm \texttt{import torch; torch.cuda.is\_available()}
  returns \texttt{True} on a GPU node.}{2}

\done \task{Write \texttt{scripts/laguna\_submit.sh}: a SLURM batch script for
  USC Laguna that requests 1 GPU (A100 if available, else V100), 32GB RAM,
  4 CPU cores, and 4-hour walltime. The script activates the venv, sets
  \texttt{TRANSFORMERS\_CACHE} to scratch, and calls \texttt{python -m
  electoral.llm.trainer --backend hpc --config configs/train\_r16.json}.
  Test with \texttt{sbatch --test-only scripts/laguna\_submit.sh} to validate
  the script parses without errors.}{2}

\done \task{Install MLX and \texttt{mlx-lm} locally on the M5:
  \texttt{pip install mlx mlx-lm}. Download \texttt{mistralai/Mistral-7B-v0.3}
  weights in MLX format: \texttt{mlx\_lm.convert --hf-path
  mistralai/Mistral-7B-v0.3 --mlx-path models/mistral-7b-mlx}.
  Confirm the model loads and generates a test completion. This download
  is $\sim$14GB --- start it immediately and work on other tasks while
  it runs.}{1}

\done \task{Run a 10-step smoke fine-tune locally with MLX: \texttt{mlx\_lm.lora
  --model models/mistral-7b-mlx --train --data data/finetune/train.jsonl
  --batch-size 2 --lora-layers 4 --iters 10}. Confirm the run completes,
  loss decreases over 10 steps, and an adapter file is saved. Log the
  output to \texttt{docs/mlx\_smoke\_run.txt}.}{1}

\done \task{Implement \texttt{electoral/llm/trainer.py}: \texttt{TrainConfig}
  dataclass with fields: \texttt{base\_model}, \texttt{lora\_rank},
  \texttt{lora\_alpha}, \texttt{learning\_rate}, \texttt{epochs},
  \texttt{batch\_size}, \texttt{train\_path}, \texttt{eval\_path},
  \texttt{output\_dir}, \texttt{backend} (``mlx'' or ``hpc''). Write
  \texttt{load\_config(path)} and \texttt{save\_config(config, path)}.
  Create \texttt{configs/train\_r16.json} and \texttt{configs/train\_r32.json}
  with the hyperparameters from the development plan.}{1}
\end{taskbox}

% ── Day 2 ────────────────────────────────────────────────────────────────────
\dayhead{2}{Tue --- Prompt Engineering \& First Training Run}

\begin{taskbox}
\done \task{Design the prompt template. Write three candidate templates in
  \texttt{docs/prompt\_candidates.md} and test each on 5 training examples
  manually using the base (un-fine-tuned) Mistral 7B locally. Evaluate:
  does the model output valid JSON? Does the structure make the task obvious?
  Select the best template and document the reasoning.}{2}

\done \task{Implement the selected prompt template in \texttt{trainer.py}:
  \texttt{format\_prompt(example)} that takes a JSONL record and returns a
  formatted string. The prompt must include: (i) system context paragraph
  establishing the political science framing and the selected party; (ii) the
  event description; (iii) RoBERTa news scores as a labeled JSON block
  (per-bloc); (iv) X bio-weighted social scores as a labeled JSON block
  (per-bloc, sourced from \texttt{social\_roberta\_scores}); (v) the
  historical analog and shock category; (vi) a closing instruction:
  ``You are analyzing the impact on the \textbf{[party]} candidate's coalition.
  Output only a JSON object with three keys:
  \texttt{delta\_bins\_race} (keys: african\_american, latino, asian, white,
  other\_race), \texttt{delta\_bins\_religion} (keys: evangelical, catholic,
  protestant, secular, jewish, muslim, other\_rel), and
  \texttt{delta\_bins\_gender} (keys: women, men, other\_gender).
  Values are one of: strong\_neg, mod\_neg, mild\_neg, slight\_neg, neutral,
  slight\_pos, mild\_pos, mod\_pos, strong\_pos.
  Also output \texttt{delta\_eff} (scalar): the pipeline will verify it matches
  $\lambda_1\sum w_i\Delta\mu_i^{\text{race}} + \lambda_2\sum v_R\Delta\mu_R^{\text{rel}}
  + \lambda_3\sum g_G\Delta\mu_G^{\text{gen}}$.
  For non-gendered shocks, gender bins may be identical across women/men.''
  The party name must be substituted from \texttt{example["party"]}, never
  hardcoded.}{2}

\done \task{Implement the USC Laguna training path in \texttt{trainer.py}: when
  \texttt{backend=''hpc''} (parameter name unchanged; targets USC Laguna),
  use HuggingFace \texttt{transformers} +
  \texttt{peft.LoraConfig} + \texttt{BitsAndBytesConfig} (4-bit NF4).
  Implement \texttt{train\_hpc(config)} that loads the base model in 4-bit,
  applies LoRA adapters to Q/K/V/O projections, runs \texttt{Trainer.train()},
  and saves the adapter to \texttt{config.output\_dir}.}{2}

\done \task{Submit the first full USC Laguna training run: \texttt{sbatch
  scripts/laguna\_submit.sh} with \texttt{configs/train\_r16.json}.
  Monitor the job queue with \texttt{squeue -u \$USER}. Once the job starts,
  \texttt{tail -f} the SLURM output file to watch the loss curve.
  The job should complete overnight. If it times out, increase walltime and
  resubmit.}{1}

\done \task{While the USC Laguna job runs: implement \texttt{electoral/llm/eval.py}.
  \texttt{EvalResult} dataclass: \texttt{shock\_id}, \texttt{mae\_per\_bloc}
  (dict), \texttt{mean\_mae}, \texttt{output\_json} (raw model output).
  \texttt{evaluate\_model(adapter\_path, eval\_path)} loads the adapter,
  iterates through each eval example, runs \texttt{estimate\_shock}, computes
  MAE against ground-truth deltas, and returns a list of \texttt{EvalResult}.
  Write \texttt{print\_eval\_summary(results)} that logs a table of MAE per
  bloc sorted descending.}{1}
\end{taskbox}

% ── Day 3 ────────────────────────────────────────────────────────────────────
\dayhead{3}{Wed --- Constrained Decoding \& Inference Wrapper}

\begin{taskbox}
\done \task{Install \texttt{outlines}: \texttt{pip install outlines}. Confirm
  it imports correctly alongside the existing HuggingFace stack. Read the
  \texttt{outlines.generate.json} documentation. Write a standalone
  10-line proof-of-concept script \texttt{scripts/outlines\_poc.py} that
  loads a tiny HuggingFace model (\texttt{gpt2}), generates constrained JSON
  with a simple schema (\texttt{\{``x'': int\}}), and prints the result.
  Confirm the output is always valid JSON regardless of how many times
  you run it.}{1}

\done \task{Convert \texttt{ShockResponseData} to a \textbf{Pydantic} model
  in \texttt{electoral/artifacts.py} (or a parallel
  \texttt{ShockResponseSchema} Pydantic class if keeping the frozen
  dataclass for internal use). Define the 10 bloc fields as
  \texttt{Annotated[str, Literal["strong\_neg", "mod\_neg", "mild\_neg",
  "weak\_neg", "neutral", "weak\_pos", "mild\_pos", "mod\_pos",
  "strong\_pos"]]} --- one field per canonical bloc. Passing this Pydantic
  model directly to \texttt{outlines.generate.json(model,
  ShockResponseSchema)} has two critical advantages over a manual
  \texttt{json\_schema()} dict: (i) the LLM output schema and the
  FastAPI \texttt{response\_model} are the \emph{same Python object},
  so they can never drift out of sync; (ii) Pydantic's \texttt{Literal}
  type maps cleanly to the 9-token enum constraint that constrained
  decoding enforces. Write a unit test confirming the Pydantic model
  validates a correct output and rejects an unknown bin token.}{1}

\done \task{Implement \texttt{electoral/llm/inference.py}. Class
  \texttt{ShockEstimator} with \texttt{\_\_init\_\_(self, adapter\_path,
  base\_model)} that loads the model once and keeps it in memory.
  Method \texttt{estimate(self, event, intensity)} that: (i) formats
  the prompt via \texttt{format\_prompt}; (ii) passes the
  \texttt{ShockResponseSchema} Pydantic model directly to
  \texttt{outlines.generate.json(model, ShockResponseSchema)} ---
  no manual dict construction; (iii) runs generation;
  (iv) maps bin tokens to numeric deltas via \texttt{bin\_to\_delta};
  (v) scales delta values by \texttt{intensity}; (vi) constructs and
  validates \texttt{ShockResponseData}; (vii) returns it. The
  \texttt{ShockResponseSchema} Pydantic model is also used as the
  FastAPI \texttt{response\_model} in \texttt{shock\_endpoint.py},
  guaranteeing the two schemas are always identical.}{3}

\done \task{Write \texttt{tests/test\_constrained\_output.py}: instantiate
  \texttt{ShockEstimator} with the trained adapter. Run
  \texttt{estimator.estimate(event, intensity)} 100 times with 10 different
  event strings (10 calls each). Assert: (i) every output passes
  \texttt{ShockResponseData.validate()}, (ii) all delta values are finite and
  in $[-0.15, 0.15]$, (iii) all 10 bloc keys are present, (iv) no call
  raises an exception.}{2}

\done \task{Benchmark latency: time 20 sequential calls to
  \texttt{estimator.estimate()} with a standard Ayatollah event description.
  Print mean, median, and p95 latency. If p95 exceeds 3 seconds: (i) confirm
  the model is loaded once (not reloaded per call), (ii) try reducing
  \texttt{max\_tokens} in the generator, (iii) try 4-bit quantization if
  not already active. Document the final p95 in
  \texttt{docs/latency\_benchmark.md}.}{1}
\end{taskbox}

% ── Day 4 ────────────────────────────────────────────────────────────────────
\dayhead{4}{Thu --- FastAPI Endpoint \& Baseline Comparison}

\begin{taskbox}
\done \task{Implement \texttt{electoral/api/shock\_endpoint.py}. Use FastAPI
  lifespan context to load \texttt{ShockEstimator} once at startup and store
  in \texttt{app.state.estimator}. Define three routes: \texttt{POST /estimate}
  (body: \texttt{EstimateRequest(event: str, intensity: float = 1.0)}, returns
  \texttt{ShockResponseData}), \texttt{GET /health} (returns
  \texttt{\{"status": "ok", "model": adapter\_path\}}), \texttt{GET /blocs}
  (returns list of canonical bloc names from \texttt{types.py}). Add
  \texttt{CORSMiddleware} allowing \texttt{http://localhost:3000} for
  local development.}{2}

\done \task{Add request validation and error handling to the endpoint.
  If \texttt{event} is empty or under 10 characters, return HTTP 422
  with message ``Event description too short''. If \texttt{intensity} is
  outside $[0.1, 3.0]$, return HTTP 422 with message ``Intensity out of
  range''. Wrap the \texttt{estimator.estimate()} call in a try/except:
  on any unexpected error, return HTTP 500 with a sanitized message (no
  stack trace in the response body, but log the full trace server-side).}{2}

\done \task{Run the three-way baseline comparison on the 2020 eval set.
  For each method, call the appropriate function on every eval example
  and collect \texttt{EvalResult} lists: (i) news-only: use the
  \texttt{SentimentData} artifact's news RoBERTa scores mapped through
  the elasticity regression; (ii) social-only: same but using only
  platform scores; (iii) unified LLM: use \texttt{estimator.estimate()}.
  Compute mean MAE per bloc and overall. Save to
  \texttt{artifacts/baseline\_comparison.json} and
  \texttt{docs/baseline\_comparison.md} with a formatted table.}{2}

\done \task{Write \texttt{tests/test\_endpoint\_latency.py}: start the
  FastAPI app as a subprocess with \texttt{subprocess.Popen}, wait for
  the health check to return 200, then send 5 POST requests to
  \texttt{/estimate} with the Ayatollah event description. Assert each
  response: (i) HTTP 200, (ii) valid JSON, (iii) parses into
  \texttt{ShockResponseData} without error, (iv) response time under
  3 seconds. Teardown: terminate the subprocess.}{2}
\end{taskbox}

% ── Day 5 ────────────────────────────────────────────────────────────────────
\dayhead{5}{Fri --- Hyperparameter Decision, Kernel Wiring \& PR}

\begin{taskbox}
\done \task{Compare the $r=16$ and $r=32$ eval results. Decision rule: adopt
  $r=32$ if mean MAE improves by more than 5\% across all blocs. If neither
  run is satisfactory (mean MAE $> 0.05$ vote-share points), consider the
  Qwen2.5 14B fallback: re-run the smoke MLX test with
  \texttt{Qwen/Qwen2.5-14B-Instruct} and compare structured output reliability.
  Document the final model and rank decision in \texttt{DECISIONS.md} with
  the supporting MAE numbers.}{2}

\done \task{Retrieve the winning adapter weights from USC Laguna scratch using
  \textbf{\texttt{scp}}, not Syncthing. USC Laguna scratch nodes frequently kill
  background daemons including Syncthing; relying on it for model weight
  transfer risks silent failures or corrupted weights from race conditions.
  Run: \texttt{scp -r laguna:\$SCRATCH/electoral-equilibrium/adapters/
  mistral-7b-electoral/ adapters/}. Verify the adapter loads correctly on M5.
  Tag the version in \texttt{adapters/mistral-7b-electoral/VERSION} with the
  USC Laguna job ID, training config path, and eval MAE. After verifying the
  transfer, add the adapter directory to Syncthing for local machine availability.}{1}

\done \task{Wire \texttt{electoral/kernels/finetune.py}: implement
  \texttt{build\_llm\_finetune(config)} that reads \texttt{LLMFineTuneData}
  from the artifact store (or constructs it from the adapter directory if
  already trained). Wire into \texttt{stages.py}. The finetune stage is
  idempotent: if an adapter already exists at the configured path, skip
  training and log ``adapter found, skipping training''.}{2}

\done \task{Wire \texttt{electoral/kernels/shock.py}: implement
  \texttt{build\_shock\_response(config, event, intensity)} that calls
  \texttt{estimate\_shock}, \texttt{bootstrap\_cov}, and
  \texttt{solve\_rebalanced} in sequence, wraps results in
  \texttt{ShockResponseData} and \texttt{EquilibriumData} payloads, and
  writes both artifacts. Run the full pipeline smoke config end-to-end
  including the shock stage.}{2}

\done \task{Weekly PR: \texttt{feature/week5-llm-finetune}. PR description must
  include: (i) the three-way MAE comparison table, (ii) the model and rank
  decision with justification, (iii) USC Laguna job IDs and adapter path, (iv)
  a sample of 3 eval examples showing input event, predicted deltas, and
  ground-truth deltas side by side.}{1}

\done \task{Brief supervision check-in with Prof.\ Espinosa: share 3--4 eval
  examples where the model's predictions diverge most from ground truth.
  Get domain perspective on whether divergences reflect real political
  dynamics. Juhi decides on any model adjustments.}{2}
\end{taskbox}

\newpage

% ════════════════════════════════════════════════════════════════════════════
\section{Week 5 --- Stochastic Optimization \& Monte Carlo (40 hrs)}
% ════════════════════════════════════════════════════════════════════════════

\textbf{Goal.} By Friday, the CVXPY optimizer rebalances voter distributions
after a shock, $10{,}000+$ Monte Carlo simulations run on USC Laguna, and the
\texttt{SimulationData} artifact is produced and validated.

\textbf{End condition.}
\begin{itemize}[leftmargin=1.6em, itemsep=1pt]
  \item \done Optimizer produces feasible weights for all test shocks.
  \item \done Monte Carlo win probability converges within $\pm 0.005$
        between $N=5{,}000$ and $N=10{,}000$.
  \item \done All optimization and simulation tests pass.
  \item \done PR merged.
\end{itemize}

% ── Day 1 ────────────────────────────────────────────────────────────────────
\dayhead{1}{Mon --- Schema Migration, Post-Shock Optimization}

\begin{taskbox}
\done \task{\textbf{Add \texttt{win\_probability\_low} and
  \texttt{win\_probability\_high} to \texttt{SimulationData} before writing
  any Monte Carlo code.}
  The current \texttt{SimulationData} only stores \texttt{win\_probability}
  (a point estimate) and \texttt{percentiles} (per-bloc coalition-weight
  distributions). It is missing the 90\% confidence interval on the
  \emph{win-probability} itself, which is a named deliverable in
  \texttt{DECISIONS.md} and required by the WinGauge UI logic:

  \begin{itemize}[leftmargin=1.6em, itemsep=1pt]
    \item \texttt{win\_probability\_low: float} --- 5th percentile of
          $\Pr(\tilde{\mu}_{\mathrm{eff}}(\mathbf{w}^{(n)}) \geq \veq)$
          across the $N$ ILR draws.
    \item \texttt{win\_probability\_high: float} --- 95th percentile.
  \end{itemize}

  These come from propagating $\boldSigma_\Delta$ uncertainty through the
  ILR back-transform (not Bernoulli variance). The WinGauge fires
  \textit{``No feasible path''} when \texttt{win\_probability\_high}~$< \veq$
  and \textit{``Uncertain path''} when
  \texttt{win\_probability\_low}~$< \veq \leq$~\texttt{win\_probability\_high}.

  Add both fields to the frozen dataclass, update \texttt{validate()} to
  check $0 \leq \texttt{low} \leq \texttt{high} \leq 1$, update
  \texttt{from\_dict} / \texttt{to\_dict}, update the \texttt{run\_simulations}
  stub in \texttt{stages.py} to provide placeholder values of \texttt{0.0}
  and \texttt{1.0} respectively, and update \texttt{test\_artifact\_roundtrip.py}.
  Do \emph{not} start the Monte Carlo kernel until this schema change is
  committed and all tests pass.}{1}

\done \task{Implement \texttt{electoral/optimization/simplex.py}:
  \texttt{project\_simplex(v)} using the Duchi et al.\ (2008) $O(n \log n)$
  algorithm. The algorithm: sort $v$ descending, find the largest $\rho$ such
  that $v_{(\rho)} - \frac{1}{\rho}(\sum_{j \leq \rho} v_{(j)} - 1) > 0$,
  set $\theta = \frac{1}{\rho}(\sum_{j \leq \rho} v_{(j)} - 1)$, return
  $\max(v - \theta, 0)$. Write three unit tests: (i) already-valid simplex
  vector returns unchanged, (ii) negative components are zeroed and remainder
  redistributed, (iii) output always sums to 1.0 within \texttt{tol=1e-9}.}{3}

\done \task{Implement \texttt{electoral/optimization/cvx.py}:
  \texttt{solve\_rebalanced(mu\_tilde, cov\_delta, target)} using CVXPY with
  the \textbf{max P(win) objective} (not min-variance).
  Formulate the Sharpe-ratio objective:
  numerator = \texttt{mu\_tilde @ w - target} (linear in \texttt{w});
  denominator = \texttt{cp.sqrt(cp.quad\_form(w, cov\_delta))} (convex in \texttt{w});
  objective = \texttt{cp.Maximize(cp.ratio(numerator, denominator))}.
  Declare \texttt{problem.solve(qcp=True)} (Disciplined Quasiconvex Programming).
  \textbf{Assert \texttt{problem.is\_dcp(qcp=True) == True} before solving} ---
  if this assertion fails, the objective is not quasi-convex and the solver
  cannot guarantee the global optimum. This assertion must be a unit test.
  On \texttt{SolverError} or infeasibility, return \texttt{EquilibriumData}
  with \texttt{feasible=False}.}{4}

\done \task{Implement \texttt{solve\_equal\_weight\_rebalanced(blocs, delta\_w)}
  as a fallback in \texttt{cvx.py}: if the QP is infeasible after 5 retries,
  return an equal-weight distribution with \texttt{feasible=False} and
  \texttt{target\_met=False}. Log a warning. This ensures the web app
  never crashes due to an infeasible shock input.}{2}
\end{taskbox}

% ── Day 2 ────────────────────────────────────────────────────────────────────
\dayhead{2}{Tue --- Covariance Bootstrap \& Full Shock Chain}

\begin{taskbox}
\done \task{Implement \texttt{electoral/models/bootstrap.py} using
  \textbf{Ledoit-Wolf shrinkage} rather than raw bootstrap covariance.
  With only 6--8 election cycles, bootstrapping a $10 \times 10$ covariance
  matrix yields a rank-deficient result ($\leq 7$ degrees of freedom for
  55 parameters). Applying \texttt{psd\_repair} to a rank-deficient matrix
  just adds noise --- the optimizer would be minimizing against your
  artificial $\epsilon I$ injection rather than historical electoral reality.

  Use \texttt{sklearn.covariance.LedoitWolf} instead: it applies an
  analytically optimal shrinkage toward a diagonal target, is specifically
  designed for the $p > n$ regime, and produces a well-conditioned,
  full-rank covariance estimate. Usage: instantiate
  \texttt{LedoitWolf(assume\_centered=False)}, call \texttt{.fit(delta\_matrix)}
  where \texttt{delta\_matrix} has shape \texttt{(n\_cycles, 10)}, read
  \texttt{.covariance\_} --- always full-rank, no \texttt{psd\_repair} needed.
  Keep \texttt{bootstrap\_cov} as an alternative for comparison; document
  in \texttt{DECISIONS.md} that Ledoit-Wolf is the primary estimator.
  Write unit tests: (i) output is symmetric and PSD, (ii) condition number
  is finite and below $10^6$, (iii) same input produces identical output.}{3}

\done \task{Implement the sensitivity-weighted covariance variant
  \texttt{bootstrap\_cov\_weighted(panel\_df, social\_elasticities,
  n\_bootstrap=1000, seed=42)}: weight each cycle's contribution by its
  mean absolute social media elasticity across platforms, so cycles with
  stronger social media reactions contribute more to the covariance estimate.
  This fuses the social media signal directly into the stochastic model.
  Store both variants; the paper will compare them.}{2}

\done \task{Integrate the full shock response chain in
  \texttt{electoral/kernels/shock.py}: \texttt{build\_shock\_response(config,
  event\_description, intensity)}. Step 1: call \texttt{estimate\_shock}
  to get \texttt{delta\_mu} dict ($\Delta\mu_i$ per bloc). Step 2: load
  $\boldmu$ and $\mathbf{w}^{(0)}$ from the \texttt{BaselinePortfolioData}
  artifact. Step 3: compute $\tilde{\boldmu} = \boldmu + \boldsymbol{\Delta\mu}$.
  Step 4: call \texttt{bootstrap\_cov} to get $\boldSigma_\Delta$.
  Step 5: call \texttt{solve\_rebalanced(mu\_tilde, cov\_delta, target)}
  to produce \texttt{EquilibriumData} (which stores both $\tilde{\mathbf{w}}$
  and $\tilde{\boldmu}$). Step 6: construct and validate
  \texttt{ShockResponseData} wrapping \texttt{delta\_mu} and the covariance.
  Wire into \texttt{stages.py}.}{3}
\end{taskbox}

% ── Day 3 ────────────────────────────────────────────────────────────────────
\dayhead{3}{Wed --- Monte Carlo Engine}

\begin{taskbox}
\done \task{Implement \texttt{electoral/simulation/montecarlo.py}:
  \texttt{run\_simulations(w\_star, mu\_tilde, cov\_delta, target, n, seed)}.
  Use the \textbf{Logistic-Normal with ILR parameterization} --- not Dirichlet
  (which forces negative off-diagonal covariances) and not Gaussian-then-project
  (which destroys covariance structure asymmetrically at simplex boundaries).

  \textbf{Implementation steps:}
  \begin{itemize}[leftmargin=1.4em, itemsep=1pt]
    \item Construct a Helmert contrast matrix $V \in \mathbb{R}^{(K-1)\times K}$.
    \item Map $\mathbf{w}^*$ to ILR coordinates:
          $\mathbf{z}^* = V^\top \log(\mathbf{w}^*)$
    \item Propagate $\Sigma_\Delta$ to ILR space via the Jacobian $J$ of the
          ILR transformation: $\Sigma_{\text{ILR}} = J \Sigma_\Delta J^\top$
    \item Draw $N$ samples:
          $\mathbf{y}^{(n)} \sim \mathcal{N}(\mathbf{z}^*, \Sigma_{\text{ILR}})$
    \item Back-transform via softmax:
          $w_i^{(n)} = \exp(y_i^{(n)}) / \sum_j \exp(y_j^{(n)})$
    \item Compute win flags and confidence interval bounds (5th and 95th percentiles)
  \end{itemize}
  If $w_i^* = 0$ exactly for any bloc, report it as \texttt{infeasible\_bloc}
  in \texttt{EquilibriumData} and exclude it from the simulation --- do
  \emph{not} floor to 0.01 (flooring distorts log-ratio space non-linearly).

  \textbf{Required test:} \texttt{test\_montecarlo\_degenerate} --- with
  $w^*_0 = 0.99$, $w^*_{1..4} = 0.0025$, assert all $N$ samples sum to
  1.0 within tol=$10^{-9}$ and win\_probability is in $[0, 1]$.}{5}

\done \task{Implement \texttt{project\_simplex\_batch(W)} in
  \texttt{simplex.py}: takes an $(N \times d)$ matrix and returns each
  row projected onto the probability simplex. This function is a
  \textbf{fallback utility} used only when the CVXPY optimizer returns an
  infeasible solution and equal-weight distribution is substituted --- it
  is \emph{not} the primary Monte Carlo sampling path (which uses Dirichlet
  and produces simplex-valid samples directly).

  Implement Duchi's $O(n \log n)$ algorithm per row. A Python \texttt{for}
  loop over rows is acceptable and correct here: at $N = 10{,}000$,
  $d = 10$, the dominant cost is the sort, and Python overhead per row is
  negligible. Vectorizing the dynamic per-row threshold $\rho$ in pure NumPy
  is brittle without a loop --- do not attempt it. If the fallback path
  ever becomes a throughput bottleneck (it won't during the SRP), use
  \texttt{numba.jit} to compile the inner loop.

  Tests: (i) all rows sum to 1.0 within \texttt{tol=1e-9}, (ii) all values
  non-negative, (iii) already-valid row returns unchanged,
  (iv) $N=10{,}000$ completes in under 5 seconds on M5 (loop is fine).}{2}

\done \task{Write \texttt{scripts/run\_montecarlo\_laguna.sh}: SLURM batch
  script for USC Laguna requesting 8 CPU cores, 16GB RAM, 2-hour walltime.
  The script loads a shock event from \texttt{artifacts/shock\_\{id\}.json},
  runs \texttt{python -m electoral.simulation.montecarlo} with $N=50{,}000$,
  and saves \texttt{SimulationData} to \texttt{artifacts/sim\_\{id\}.json}.
  Submit for the Ayatollah event as a test job. Confirm the job
  completes and produces a valid artifact.}{1}

\done \task{Write \texttt{tests/test\_montecarlo.py}: (i) smoke run $N=100$
  completes in under 2 seconds on M5, (ii) win probability is in $[0,1]$,
  (iii) with $\boldmu$ all equal to 0.8 and \texttt{target=0.535}, win
  probability $> 0.95$, (iv) with $\boldmu$ all equal to 0.2, win
  probability $< 0.05$, (v) two runs with same seed produce identical
  \texttt{SimulationData.to\_dict()} outputs, (vi) percentiles list has
  exactly 5 values per bloc in $[0,1]$.}{1}
\end{taskbox}

% ── Day 4 ────────────────────────────────────────────────────────────────────
\dayhead{4}{Thu --- Convergence Tests, Sensitivity Analysis \& Kernel Wiring}

\begin{taskbox}
\done \task{Write \texttt{tests/test\_mc\_convergence.py}: run simulations at
  $N=1{,}000$, $N=5{,}000$, and $N=10{,}000$ with the same seed. Assert the
  $N=5{,}000$ and $N=10{,}000$ win probability estimates differ by less than
  $0.005$. Also assert the $N=1{,}000$ estimate is within $0.02$ of the
  $N=10{,}000$ estimate (looser bound). Log the three estimates side by side.}{2}

\done \task{Write \texttt{tests/test\_optimization.py}: (i) a feasible shock
  (small delta, weak covariance) returns weights summing to 1.0 with
  \texttt{feasible=True} and \texttt{target\_met=True}, (ii) an extreme
  shock that pushes the portfolio below target returns \texttt{target\_met=False}
  and logs the relaxation steps, (iii) \texttt{project\_simplex} output
  always non-negative and sums to 1.0, (iv) same inputs produce identical
  \texttt{EquilibriumData.to\_dict()} across two calls.}{2}

\done \task{Implement a shock sensitivity analysis script
  \texttt{scripts/sensitivity\_analysis.py}: for a given shock event,
  vary the intensity parameter from 0.5 to 2.0 in steps of 0.25 and
  plot how win probability and per-bloc weights change. Save the plot as
  \texttt{artifacts/sensitivity\_\{shock\_id\}.png}. Run on the Ayatollah
  event as a case study. This figure will be useful in the paper.}{2}

\done \task{Wire \texttt{electoral/kernels/optimize.py} into \texttt{stages.py}:
  \texttt{build\_optimization(config, shock\_response\_artifact)} that loads
  the \texttt{ShockResponseData}, runs \texttt{solve\_rebalanced}, and writes
  \texttt{EquilibriumData}. Run the full pipeline smoke config end-to-end
  and confirm all seven artifact files exist and validate.}{2}
\end{taskbox}

% ── Day 5 ────────────────────────────────────────────────────────────────────
\dayhead{5}{Fri --- Metrics, Reporting \& PR}

\begin{taskbox}
\done \task{Implement \texttt{electoral/metrics/performance.py}. Required
  functions: \texttt{win\_probability(sim\_data)} (reads from
  \texttt{SimulationData}), \texttt{equilibrium\_gap(weights, mu, target)}
  returning $\mathbf{w}^\top\boldmu - \veq$ (positive = above threshold),
  \texttt{bloc\_delta\_summary(baseline\_weights, rebalanced\_weights)}
  returning a dict of signed deltas per bloc sorted by absolute magnitude.
  All functions must accept and return JSON-native types.}{2}

\done \task{Implement \texttt{electoral/metrics/tables.py}: write
  \texttt{build\_shock\_summary\_table} (takes \texttt{shock\_ids} and
  \texttt{config}) that iterates over shock event IDs, loads the
  corresponding \texttt{EquilibriumData} and \texttt{SimulationData}
  artifacts, and assembles a summary table with columns: shock id,
  category, win probability, equilibrium gap, top-moving bloc, and
  top-moving delta. Return as a JSON-serializable list of dicts.}{2}

\done \task{Implement \texttt{electoral/reporting/export.py}:
  \texttt{export\_json(path, data)}, \texttt{export\_csv(path, rows, fieldnames)},
  and \texttt{export\_latex\_table(path, rows, caption, label)} that produces
  a \texttt{booktabs}-style LaTeX table fragment. Run the summary table for
  all 25+ shocks and save to \texttt{artifacts/shock\_summary\_table.csv}
  and \texttt{artifacts/shock\_summary\_table.tex}.}{2}

\done \task{Weekly PR: \texttt{feature/week6-optimization-montecarlo}. Include
  the convergence test output, the sensitivity analysis plot for the Ayatollah
  event, and the shock summary table (first 10 rows) in the PR description.}{1}

\done \task{Brief supervision check-in with Prof.\ Espinosa: share win
  probability estimates for 4--5 historical shocks. Get domain input on
  whether results align with known electoral outcomes. Juhi decides which
  covariance variant to use in the paper and documents the decision.}{1}
\end{taskbox}

\begin{taskbox}{Day 4.5 --- Synthetic Data Augmentation Pipeline}

\task{Design a four-axis event taxonomy for synthetic fine-tuning data
  generation. Axes: (i) domain (econ, foreign, social, candidate,
  institutional, scandal, exogenous), (ii) ideological valence
  (far\_left through far\_right, libertarian, populist, bipartisan),
  (iii) electoral effect (helps\_dem, helps\_rep, splits, realigns, neutral),
  (iv) party perspective (democrat, republican). Key design principle:
  ideological valence does NOT determine electoral effect --- a far-left event
  can hurt Democrats with moderates. Store as
  \texttt{configs/synthetic\_events.json} (62 balanced seeds) and
  \texttt{docs/synthetic\_event\_taxonomy.md}. Balance targets:
  helps\_dem/helps\_rep both $\geq$15\%, both party perspectives $\geq$40\%.}{2}

\task{Implement \texttt{scripts/synthetic/generate\_deepseek.py}:
  reads seed events from \texttt{configs/synthetic\_events.json},
  expands each seed into $N$ training records via the DeepSeek API
  (OpenAI-compatible, \texttt{deepseek-chat} model), validates every
  record against the 9-token bin schema and canonical bloc keys, and
  writes to \texttt{data/finetune/candidates.jsonl}. Key implementation
  details: batch size of 5 records per API call (20 causes truncation at
  4k tokens), max\_tokens=8000, retry logic with exponential backoff.
  Generated 1,237 records from 62 seeds with only 3 rejections. Effect
  balance: splits 48\%, helps\_dem 23\%, helps\_rep 16\%, neutral 8\%,
  realigns 5\%.}{2}

\task{Implement \texttt{scripts/synthetic/gemini\_review.py}:
  three-stage review pipeline using Gemini 2.5 Flash as a political
  plausibility reviewer. For each candidate record, Gemini checks:
  (i) direction sanity per bloc, (ii) delta\_eff sign coherence,
  (iii) cross-bloc logic (evangelical vs secular), (iv) magnitude
  realism. Triages into APPROVE / REVISE (with corrected labels) /
  HUMAN\_REVIEW. Revisions are audited in \texttt{data/finetune/revisions.jsonl}.
  A 5\% random spot-check of approved records is routed to the human
  queue regardless of verdict. Human queue
  (\texttt{data/finetune/human\_review\_queue.jsonl}) is adjudicated
  in a separate Opus/human pass during low-usage hours.
  Results: 579 approved, 647 queued for human review (Gemini 2.5
  Flash is conservative by design).}{2}

\task{Merge Gemini-approved synthetic records into the training set.
  Strip \texttt{\_review} and \texttt{\_seed\_meta} metadata fields
  before merging. Rebuild 80/20 train/eval split with seed=42.
  Final dataset: 1,136 synthetic records (577 original + 559
  Gemini-approved). Retrain Mistral 7B r=16 adapter on the expanded
  dataset. Document in \texttt{DECISIONS.md}: synthetic data pipeline
  rationale, DeepSeek generation settings, Gemini review thresholds,
  human review queue process, and the key insight that ideological
  valence and electoral effect are independent axes.}{2}

\task{NEP-weighted baseline coalition formula. Discovered that the
  original \texttt{build\_baseline\_portfolio} used CVXPY to derive
  coalition weights $\mathbf{w}^{(0)}$, producing degenerate results
  (asian=0, latino=0, african\_american=0.51). Replaced with the
  principled formula: $w_i = (\text{NEP}_i \times \mu_i) /
  \sum_j (\text{NEP}_j \times \mu_j)$, where NEP shares are from
  Edison Research 2024 National Exit Poll (actual voter shares, not
  census). Values: white=0.63, african\_american=0.12, latino=0.15,
  asian=0.05, other\_race=0.05. Resulting $\mathbf{w}^{(0)}$:
  white=0.528, african\_american=0.189, latino=0.180, asian=0.051,
  other\_race=0.052. Baseline $\mu_\text{eff}$ dropped from 0.716
  (equal weights) to 0.605 (NEP-weighted). Stored in
  \texttt{configs/base.json} as \texttt{bloc\_population\_shares}.}{2}

\end{taskbox}

\newpage

% ════════════════════════════════════════════════════════════════════════════
\section{Week 6 --- Web Application (40 hrs)}
% ════════════════════════════════════════════════════════════════════════════

\textbf{Goal.} By Friday, the core public web application is live end-to-end
with SSE streaming and deployed on Modal + Vercel. The analyst dashboard
moves to Week~8 Days~1--2 to give it the time it actually requires.

\textbf{End condition.}
\begin{itemize}[leftmargin=1.6em, itemsep=1pt]
  \item \done Public app renders per-bloc shifts progressively via SSE.
  \item \done Audit log records every \texttt{/estimate} call with latency.
  \item \done Modal + Vercel deployment live with public URL.
  \item \done PR merged.
\end{itemize}

% ── Day 1 ────────────────────────────────────────────────────────────────────
\dayhead{1}{Mon --- Next.js Setup \& Typed API Client}

\begin{taskbox}
\done \task{Initialize the Next.js app: \texttt{npx create-next-app webapp
  --typescript --tailwind --app --no-src-dir}. Remove all boilerplate page
  content. Set up \texttt{webapp/.env.local}: \texttt{NEXT\_PUBLIC\_API\_URL=
  http://localhost:8000}. Confirm \texttt{npm run dev} starts without errors
  and the browser shows a blank page.}{1}

\done \task{Install frontend dependencies: \texttt{npm install recharts zod
  @radix-ui/react-slider @radix-ui/react-collapsible lucide-react}.
  These cover charting, schema validation, the intensity slider, the
  ``How it works'' expandable, and icons respectively.}{1}

\done \task{Define TypeScript types in \texttt{webapp/lib/types.ts}: all existing
  payload types (\texttt{ShockResponseData}, \texttt{EquilibriumData},
  \texttt{SimulationData}, \texttt{EstimateResponse}) plus new dashboard
  types: \texttt{AuditEntry} (id, timestamp, event\_text, party, intensity,
  win\_prob, feasible, llm\_ms, optimizer\_ms, montecarlo\_ms),
  \texttt{CoverageCell} (bloc\_id, cycle, quality, vote\_share, sources),
  and \texttt{TrainingRun} (run\_id, epochs, train\_loss, val\_loss, mae).
  Add \texttt{Party = "democrat" | "republican"} as a shared literal type.
  All types must mirror their FastAPI response models exactly.}{1}

\done \task{Implement Zod schemas in \texttt{webapp/lib/schemas.ts} for each
  type. These are used to validate API responses at runtime so the frontend
  fails loudly rather than rendering NaN or undefined. Export
  \texttt{EstimateResponseSchema.parse(data)} as the single validation
  entry point.}{1}

\done \task{Implement \texttt{webapp/lib/api.ts}: \texttt{estimateShock(event:
  string, intensity: number): Promise<EstimateResponse>} that POSTs to
  \texttt{/estimate}, parses the response with the Zod schema, and throws
  a typed \texttt{ApiError} (with \texttt{message} and \texttt{status})
  on any failure. Implement \texttt{getBlocs(): Promise<string[]>} and
  \texttt{healthCheck(): Promise<boolean>}.}{2}

\done \task{Write a smoke integration test for the API client: in
  \texttt{webapp/lib/\_\_tests\_\_/api.test.ts}, mock \texttt{fetch} using
  \texttt{jest.fn()}, provide a valid \texttt{EstimateResponse} fixture,
  and assert: (i) the correct URL is called, (ii) the response parses
  without error, (iii) a malformed response throws \texttt{ApiError}.}{2}
\end{taskbox}

% ── Day 2 ────────────────────────────────────────────────────────────────────
\dayhead{2}{Tue --- ShockInput Component \& State Management}

\begin{taskbox}
\done \task{Implement the page-level state in \texttt{webapp/app/page.tsx}
  using \texttt{useState}: \texttt{party} (\texttt{"democrat" | "republican"},
  default \texttt{"democrat"}), \texttt{event} (string), \texttt{intensity}
  (number, default 1.0), \texttt{loading} (boolean), \texttt{error}
  (string | null), \texttt{deltaBins} (\texttt{Record<string,string> | null}),
  \texttt{shifted} (\texttt{Record<string,number> | null}),
  \texttt{equilibrium} (\texttt{EquilibriumData | null}),
  \texttt{simulation} (\texttt{SimulationData | null}).

  Implement \texttt{handleSubmit}: open an SSE \texttt{EventSource} to
  \texttt{/estimate/stream}; on \texttt{event: deltas} set \texttt{deltaBins}
  and \texttt{shifted} (populates \texttt{ShockNarrative} immediately, and
  renders \texttt{CoalitionChart} with opaque bars only --- rebalance layer
  is absent until the next event); on \texttt{event: equilibrium} set
  \texttt{equilibrium} (adds the translucent rebalance layer to
  \texttt{CoalitionChart} in-place --- no re-render of the opaque bars);
  on \texttt{event: simulation} set \texttt{simulation} (populates
  \texttt{WinGauge}); on \texttt{[DONE]} set \texttt{loading=false}.

  Layout: \texttt{ShockNarrative} above the chart, full width.
  \texttt{CoalitionChart} below it, full width (the chart is self-contained
  with both layers). \texttt{WinGauge} to the right on desktop, below on
  mobile. The progressive reveal --- opaque bars first, translucent overlay
  second, gauge last --- gives the user something to read at each SSE step.}{3}

\done \task{Implement \texttt{webapp/components/ShockInput.tsx}. Props:
  \texttt{party}, \texttt{setParty}, \texttt{event}, \texttt{setEvent},
  \texttt{intensity}, \texttt{setIntensity}, \texttt{onSubmit},
  \texttt{loading}. Render: (i) a two-button toggle labeled
  \textbf{Democrat} / \textbf{Republican}, styled blue/red respectively,
  that sets \texttt{party} state --- this is the first control the user sees;
  (ii) a \texttt{<textarea>} with placeholder ``Describe a hypothetical
  political event...'', minimum 3 rows, full width; (iii) a labeled
  \texttt{Slider} (Radix UI) for intensity 0.5--2.0 in steps of 0.1;
  (iv) a submit \texttt{<button>} showing a spinner when loading.
  Style with Tailwind. The party selection must be prominent --- it
  determines the entire analysis direction.}{3}

\done \task{Add five preset shock buttons below the textarea. Each button
  pre-fills the event field with a template string. Presets: (i) Security:
  ``An assassination attempt is made on the Democratic candidate''; (ii)
  Geopolitical: ``The US enters a major armed conflict in the Middle East'';
  (iii) Moral/Scandal: ``A major financial scandal involving the Republican
  candidate is revealed''; (iv) Electoral Surprise: ``A major October
  Surprise leak dominates the news cycle''; (v) Economic: ``A sudden
  recession is declared in the final month of the campaign''. Style each
  button with its shock category color.}{2}

\done \task{Implement the error display: below the submit button, show a
  styled red alert box when \texttt{error} is non-null. The message should
  be user-friendly (``The estimation service is unavailable. Please try again.'')
  rather than exposing the raw API error. Log the raw error to the browser
  console for debugging.}{1}
\end{taskbox}

% ── Day 3 ────────────────────────────────────────────────────────────────────
\dayhead{3}{Wed --- Prediction Panel: ShockNarrative \& CoalitionChart}

\begin{taskbox}
\done \task{Implement \texttt{webapp/components/ShockNarrative.tsx}.
  Props: \texttt{deltaBins: Record<string, string>}, \texttt{party: Party}.
  This component renders the plain-language \textbf{prediction} --- what the
  model expects will happen to each voter group after the shock, in words
  the user can read without knowing what a delta bin is.

  Logic: map each \texttt{delta\_bin} token to a phrase pair
  (direction + magnitude): \texttt{strong\_neg} $\to$ ``significant loss among'',
  \texttt{mild\_pos} $\to$ ``modest gain with'', \texttt{neutral} $\to$
  ``no meaningful change among'', etc. Group blocs by direction (gaining vs.\
  losing) and render two sentences:
  \begin{itemize}[leftmargin=1.6em, itemsep=1pt]
    \item ``This shock is predicted to \textbf{hurt} the [party] coalition's
          standing with: [Evangelicals (significant), Catholics (mild)].''
    \item ``And \textbf{help} with: [Latino voters (moderate), Women (mild)].''
  \end{itemize}
  If all blocs are neutral, render: ``The model predicts minimal coalition
  impact from this event.'' The narrative populates immediately when the
  \texttt{event: deltas} SSE event fires --- the user sees the prediction
  in words before the charts finish rendering.}{3}

\task{Implement \texttt{webapp/components/CoalitionChart.tsx} --- the
  single unified chart that shows both the prediction and the rebalance in
  one view. Props: \texttt{baseline: Record<string, number>} (pre-shock
  loyalty $\mu_i$), \texttt{shifted: Record<string, number>} (post-shock
  loyalty $\tilde{\mu}_i$ --- the \textbf{prediction}),
  \texttt{rebalanced: Record<string, number> | null} (optimizer coalition
  weights $\tilde{w}_i$ --- the \textbf{rebalance}; null until the
  \texttt{event: equilibrium} SSE event fires), \texttt{feasible: boolean},
  \texttt{party: Party}.

  \textbf{Visual design --- three series on one horizontal bar chart:}
  \begin{itemize}[leftmargin=1.6em, itemsep=2pt]
    \item \textbf{Opaque bars (prediction):} post-shock loyalty
          $\tilde{\mu}_i$, rendered at full opacity. Color: blue for Democrat,
          red for Republican. This is what the model says \emph{will happen}
          to each group's loyalty after the shock.
    \item \textbf{Translucent bars (rebalance):} optimizer coalition weight
          $\tilde{w}_i$, rendered at 35\% opacity (\texttt{fillOpacity=0.35})
          on top of the opaque bars using the same color. This is what the
          model says the coalition \emph{should look like} to still win.
          Absent until \texttt{event: equilibrium} fires; the opaque bars are
          visible and readable while the rebalance is loading.
    \item \textbf{Reference tick marks:} a vertical \texttt{ReferenceLine} at
          $\mu_i$ (pre-shock baseline) per bloc, rendered as a thin gray
          tick, so users can see how far the shock moved loyalty from the
          starting point.
  \end{itemize}

  \textbf{Why same axis works:} both $\tilde{\mu}_i$ and $\tilde{w}_i$ are
  in $[0, 1]$. $\tilde{\mu}_i$ is within-bloc loyalty (typically 0.15--0.90).
  $\tilde{w}_i$ is coalition weight (typically 0.03--0.30). They occupy
  different ranges and are visually distinct. A clear legend is required:
  ``\textbf{---} Prediction (loyalty shift)'' and
  ``\textbf{---} Rebalance (coalition weight, translucent)''. The tooltip
  on hover must label them explicitly with different units.

  \textbf{Sorting:} sort blocs by \texttt{Math.abs(shifted[bloc] -
  baseline[bloc])} descending so the most-moved blocs appear at the top.

  \textbf{Labels:} on each opaque bar, show the signed loyalty delta
  (\texttt{+8pp}, \texttt{-3pp}) in green/red. On each translucent bar,
  show the coalition weight as a percentage (\texttt{14\%}).

  \textbf{Feasibility:} if \texttt{feasible=false} and \texttt{rebalanced}
  is non-null, replace the translucent bars with a striped red pattern and
  add a banner: ``No feasible coalition path under this shock.''}{5}

\done \task{Add a section header above \texttt{CoalitionChart} with two
  sub-labels that map to the two visual layers:
  \textbf{``What will happen''} (opaque, bold) and
  \textbf{``Most likely rebalance''} (translucent, muted) --- styled as a
  color-coded legend inline with the header so users understand what each
  layer means before reading the chart. Below the chart, add a summary
  line: ``Equilibrium status: \textbf{MET / NOT MET}'' in green/red, plus
  the gap formatted as \texttt{"+3.2 pp above target"}.}{1}
\end{taskbox}

% ── Day 4 ────────────────────────────────────────────────────────────────────
\dayhead{4}{Thu --- WinGauge, SSE Streaming \& Audit Log}

\begin{taskbox}
\done \task{Implement \texttt{webapp/components/WinGauge.tsx}. Props:
  \texttt{winProbability: number}, \texttt{percentiles: Record<string,
  number[]>}. Render a semicircular SVG gauge (viewBox \texttt{0 0 200 110}).
  Draw the arc as three colored segments: red (0--45\%), yellow (45--55\%),
  green (55--100\%). Draw a needle rotating from 0° (left) to 180° (right)
  based on \texttt{winProbability}. Animate the needle with a CSS transition
  on value change (\texttt{transition: transform 0.8s ease-out}).}{2}

\done \task{Add a percentile distribution strip below the gauge: a compact
  horizontal box-plot showing p5, p25, p50, p75, p95 per bloc from
  \texttt{SimulationData.percentiles}. Use a small Recharts
  \texttt{ComposedChart} with a \texttt{Bar} for the IQR and
  \texttt{ErrorBar} for p5--p95 whiskers.}{1}

\done \task{Add the \textbf{market prior display} to the \texttt{WinGauge}
  component. After the model win probability renders, call
  \texttt{GET /api/market-prior?party=\&event=} (new FastAPI route that
  queries Polymarket and Manifold for the nearest relevant contract via
  \texttt{electoral/markets/aggregator.py}). Display the volume-weighted
  result as ``Market consensus: \textbf{X\%}'' in a muted gray style
  directly below the model estimate. If no active contract is found,
  show ``No active market'' in italics. The juxtaposition is the key
  epistemic display: the market aggregates all current public information;
  the model shows how the coalition structure shifts under the hypothetical.
  Add disclaimer: ``Model: $N=10{,}000$ simulations. Market: live
  Polymarket/Manifold prices.''}{2}

\done \task{Implement \textbf{SSE streaming} on the FastAPI backend.
  Add \texttt{GET /estimate/stream} using \texttt{StreamingResponse}
  with \texttt{media\_type="text/event-stream"}. The generator yields
  three events: \texttt{event: deltas} (LLM output), \texttt{event:
  equilibrium} (optimizer), \texttt{event: simulation} (Monte Carlo).
  Wrap the optimizer and Monte Carlo calls correctly for their respective
  concurrency models:
  \begin{itemize}[leftmargin=1.6em, itemsep=1pt]
    \item \textbf{CVXPY optimizer} --- use \texttt{ProcessPoolExecutor}:
          \texttt{loop.run\_in\_executor(process\_pool, solve\_rebalanced,
          ...)}.
          CVXPY's underlying C solvers (SCS/ECOS/OSQP) are not thread-safe
          and will segfault under \texttt{ThreadPoolExecutor}. Passing
          \texttt{None} defaults to a thread pool --- never pass \texttt{None}.
          Instantiate \texttt{ProcessPoolExecutor(max\_workers=1)} once at
          startup and reuse it.
    \item \textbf{Monte Carlo} (pure NumPy) --- use \texttt{ThreadPoolExecutor}:
          \texttt{loop.run\_in\_executor(thread\_pool, run\_simulations, ...)}.
          NumPy releases the GIL during array operations; a thread pool is
          safe and lower-overhead than spawning a new process.
  \end{itemize}
  Each stage starts as soon as the previous one emits its SSE event.
  Expected latency improvement: $\sim$30\% on total round-trip. Update
  \texttt{webapp/lib/api.ts}: add \texttt{estimateShockStream(event,
  intensity, onDeltas, onEquilibrium, onSimulation)} that opens an
  \texttt{EventSource}, calls each callback as its event fires, and
  closes on \texttt{[DONE]}.}{3}

\done \task{Implement \textbf{audit logging}. Create
  \texttt{electoral/api/audit.py}: \texttt{AuditLogger} class connecting
  to \texttt{data/audit.duckdb}. On first connect, run:
  \texttt{CREATE TABLE IF NOT EXISTS estimates (id INTEGER PRIMARY KEY,
  timestamp TIMESTAMP, event\_text VARCHAR, intensity FLOAT, deltas\_json
  VARCHAR, feasible BOOLEAN, target\_met BOOLEAN, win\_prob FLOAT, llm\_ms
  INTEGER, optimizer\_ms INTEGER, montecarlo\_ms INTEGER, backend VARCHAR)}.
  Implement \texttt{log\_estimate(...)} inserting one row per call. Wire
  into both \texttt{/estimate} and \texttt{/estimate/stream} endpoints.
  Add \texttt{GET /api/audit?limit=100} returning the most recent rows as
  JSON. Verify: call the endpoint 5 times, query DuckDB directly, confirm
  5 rows exist.}{2}
\end{taskbox}

% ── Day 5 ────────────────────────────────────────────────────────────────────
\dayhead{5}{Fri --- Polish, Deployment \& PR}

\begin{taskbox}
\done \task{Assemble \texttt{webapp/app/page.tsx}: two-column layout on desktop,
  single-column on mobile. Add a sticky header with project name, a
  one-sentence description, and a link to the methodology document.
  Add the ``How it works'' expandable (\texttt{@radix-ui/react-collapsible})
  and the research disclaimer. Test on iPhone viewport (375px).}{2}

\done \task{Add a ``Share this result'' button that serializes \texttt{party},
  \texttt{event}, and \texttt{intensity} as URL query params and copies to
  clipboard. Test that the URL pre-fills the full form including party
  selection on reload.}{1}

\done \task{Deploy: \texttt{modal deploy deploy/modal\_app.py}. Connect
  Vercel to the GitHub repo, set \texttt{NEXT\_PUBLIC\_API\_URL} to the
  Modal URL. Confirm the full public end-to-end flow works for both party
  modes. Scaffold \texttt{deploy/hpc\_vllm.sh} and
  \texttt{deploy/backend\_router.py} (leave deactivated).}{2}

\done \task{Weekly PR: \texttt{feature/week7-webapp}. Include: three public app
  screenshots (empty state, Democrat Ayatollah result, Republican result,
  mobile layout) and the live Modal + Vercel URLs. Tag ``demo-ready''.
  \textbf{Analyst dashboard is Week~8 Days~1--2; do not attempt it here.}}{1}

\done \task{Supervision check-in with Prof.\ Espinosa: live demo via public URL.
  Walk through the Ayatollah shock in both party modes. Collect domain feedback
  on result framing and disclaimer wording for use in Week~8 paper writing.}{2}
\end{taskbox}

\newpage

% ════════════════════════════════════════════════════════════════════════════
\section{Week 7 --- Analyst Dashboard, Hardening \& Deliverables (40 hrs)}
% ════════════════════════════════════════════════════════════════════════════

\textbf{Goal.} By Friday, the analyst dashboard is live, all integration tests
pass, the pipeline is replication-ready, and the paper and poster are complete.
The dashboard was moved here from Week~7 to give it realistic time: SSE wiring,
DuckDB, password auth, and six data-backed panels is a two-day build, not one.

\textbf{End condition.}
\begin{itemize}[leftmargin=1.6em, itemsep=1pt]
  \item \done Analyst dashboard live at \texttt{/dashboard} with all six panels.
  \item \done All six integration tests pass on a clean checkout.
  \item \done Replication package complete and documented.
  \item \done \texttt{CONTINUOUS\_UPDATE.md} written; \texttt{just continuous} runs clean.
  \item \done Research paper draft submitted to SRP committee --- must include:
        (i) formal quasiconvexity proof for the DQCP objective;
        (ii) explicit additive independence assumption statement;
        (iii) additive vs.\ raked model comparison on held-out cycles;
        (iv) MMD + PCA + PCD synthetic data validation report.
  \item \done Poster submitted to SRP.
  \item \done Final PR merged and tagged \texttt{v1.0}.
\end{itemize}

% ── Day 1 ────────────────────────────────────────────────────────────────────
\dayhead{1}{Mon --- Analyst Dashboard: Auth, Backend Routes \& Panels 1--3}

\begin{taskbox}
\done \task{Implement dashboard authentication. Create
  \texttt{webapp/app/dashboard/page.tsx}: render a login form when no valid
  \texttt{dashboard\_session} cookie exists; on submit POST to
  \texttt{/api/dashboard/auth} which compares against the
  \texttt{DASHBOARD\_PASSWORD} env var and sets an 8-hour session cookie.
  Implement \texttt{webapp/app/api/dashboard/auth/route.ts} as the Next.js
  route handler. On success, render the dashboard shell with a
  \texttt{<DashboardNav>} sidebar (links to each panel, logout button).}{2}

\done \task{Implement the three FastAPI backend routes for dashboard data and
  wire them in \texttt{shock\_endpoint.py}. \texttt{GET /api/coverage}: load
  the \texttt{VoterPanelData} artifact and return a JSON matrix of per-cell
  quality (\texttt{"multi-source"}, \texttt{"single-source"},
  \texttt{"imputed"}, \texttt{"missing"}) with raw \texttt{vote\_share} and
  sources. \texttt{GET /api/sentiment-dist?category=}: load
  \texttt{SentimentData} and return per-bloc elasticity score arrays,
  optionally filtered by shock category. \texttt{GET /api/bio-coverage}:
  return per-shock bio classifier coverage counts (keyword, SetFit,
  fallback).}{2}

\done \task{Implement \textbf{Panel 1 --- Data Coverage Matrix} in
  \texttt{CoverageMatrix.tsx}: a Recharts \texttt{ScatterChart} rendered
  as a heatmap (bloc $\times$ cycle), colored dark blue (multi-source)
  through light blue (single) through gray (imputed) to white (missing).
  Tooltip shows raw \texttt{vote\_share} and sources.}{2}

\done \task{Implement \textbf{Panel 2 --- RoBERTa Score Distributions}: a
  \texttt{ComposedChart} showing min, p25, median, p75, max per bloc across
  all shocks; shock category dropdown re-fetches with \texttt{?category=}.
  Implement \textbf{Panel 3 --- Bio Classifier Coverage}: stacked
  \texttt{BarChart}, one column per shock, three segments (keyword /
  SetFit / fallback). Both panels share a two-column grid row.}{2}
\end{taskbox}

% ── Day 2 ────────────────────────────────────────────────────────────────────
\dayhead{2}{Tue --- Analyst Dashboard: Backend Routes \& Panels 4--6}

\begin{taskbox}
\done \task{Implement the three remaining FastAPI backend routes.
  \texttt{GET /api/training-logs}: read HPC training log files from
  \texttt{rawdata/hpc\_logs/} (synced via Syncthing from HPC) and return
  per-epoch train/val loss arrays per run ID. \texttt{GET /api/convergence}:
  load the most recent \texttt{SimulationData} artifact and return win
  probability at $N = 1{,}000$, $5{,}000$, $10{,}000$ with p5/p95 bands.
  \texttt{GET /api/audit?limit=100\&search=}: query
  \texttt{data/audit.duckdb} with an optional \texttt{WHERE event\_text
  LIKE \%search\%} clause; return rows with all columns including
  \texttt{party}.}{2}

\done \task{Implement \textbf{Panel 4 --- Fine-Tuning Loss Curves}:
  Recharts \texttt{LineChart} with train (solid) and val (dashed) series,
  horizontal reference line at the 2020 MAE, run-selector dropdown.
  Implement \textbf{Panel 5 --- Monte Carlo Convergence}: \texttt{AreaChart}
  with p5/p95 shaded band, dashed final estimate line, $\pm 0.005$
  tolerance band, Pass/Fail badge.}{2}

\done \task{Implement \textbf{Panel 6 --- Audit Log Query Interface}:
  sortable table with columns timestamp, event (40-char truncated),
  party, intensity, win probability, feasibility, total latency. Free-text
  search box. Row-expand shows the full per-bloc delta vector as a mini
  \texttt{CoalitionChart} and per-stage latency breakdown.}{2}

\done \task{Assemble the full dashboard layout: two-column grid for Panels
  1--2, full-width for Panels 3--6. Add page header ``Electoral Equilibrium
  --- Analyst Dashboard'' with live estimate count from DuckDB. Test all
  six panels with real pipeline data. Verify panel filters work.}{1}
\end{taskbox}

% ── Day 3 ────────────────────────────────────────────────────────────────────
\dayhead{3}{Wed --- Feedback, Integration Tests \& Replication Package}

\begin{taskbox}
\done \task{Implement any domain feedback from Prof.\ Espinosa's Week~7
  supervision check-in. Common changes: bloc display name corrections,
  chart axis relabeling, disclaimer rewording. Re-run the manual end-to-end
  test after each change. Commit as \texttt{fix: incorporate-supervision-feedback}.}{2}

\done \task{Write \texttt{tests/test\_integration\_end\_to\_end.py}: run the
  full pipeline on the smoke config from a clean state. Assert all stage
  artifact files exist at the expected paths, parse into their typed payload
  classes without error, and pass \texttt{validate()}. Assert the artifact
  envelope contains \texttt{stage}, \texttt{run\_key}, \texttt{metadata},
  \texttt{data} keys.}{2}

\done \task{Write \texttt{tests/test\_determinism\_full\_run.py}: run the
  pipeline twice with the same config and seed writing to two separate
  output directories. For each stage, load both artifacts, parse the
  \texttt{data} payload dict, and assert field-by-field equality for all
  non-path fields. Log any non-deterministic fields found.}{2}

\done \task{Write \texttt{tests/test\_artifact\_schemas.py}: for each artifact
  produced by the smoke run, assert the envelope keys exist with correct
  types (\texttt{stage} is str, \texttt{metadata} is dict, etc.), and assert
  all required payload fields declared in the artifact classes are present
  in the \texttt{data} dict.}{2}

\done \task{Run the full test suite on a clean virtual environment from a fresh
  \texttt{git clone}. Fix any import errors, missing \texttt{\_\_init\_\_.py}
  files, hard-coded absolute paths, or undeclared dependencies until the full
  suite passes on a clean machine state.}{2}

\done \task{Run the complete pipeline on real data with
  \texttt{configs/paper\_baseline.json}. Record total runtime per stage.
  Commit all stage artifacts to \texttt{artifacts/paper\_baseline/}.
  Tag \texttt{v1.0-paper-baseline}. These artifacts must not be modified
  after this point --- they are the paper's ground truth.}{2}

\done \task{Write \texttt{README.md} ``How to Reproduce This Run'' section
  and \texttt{scripts/verify\_artifacts.py}. README must include exact Python
  version, step-by-step commands, HPC setup, expected runtimes, and how to
  verify outputs. The script loads paper-baseline and fresh artifacts
  side-by-side and prints a pass/fail report.}{2}
\end{taskbox}

% ── Day 4 ────────────────────────────────────────────────────────────────────

\begin{taskbox}
\done \task{Write the \textbf{Methodology} section (target: 1,800 words).
  Structure: (1) \textit{Data Sources} --- describe the five survey sources
  and the social media collection across four platforms; introduce the
  baseline-adjustment procedure for selection bias. (2) \textit{Voter
  Equilibrium Framework} --- state the portfolio analogy formally with the QP
  from the development plan; define $\veq = 0.535$ and explain the 3\%
  margin-of-error rationale. (3) \textit{Two-Stage NLP Pipeline} --- describe
  RoBERTa as feature engineering layer, explain the fine-tuning dataset
  construction including social media elasticity features, state the QLoRA
  hyperparameters in a table. (4) \textit{Constrained Decoding} --- one
  paragraph explaining why standard sampling is insufficient for the web app
  and how \texttt{outlines} guarantees schema validity.}{4}

\done \task{Write the \textbf{Results} section (target: 1,200 words).
  Structure: (1) \textit{Baseline Portfolio} --- report the optimized bloc
  weights and cross-validate against Prof.\ Espinosa's documented thresholds
  (Evangelical 24\%, etc.); include the LOCO validation accuracy table.
  (2) \textit{Baseline Comparison} --- present the three-way MAE table
  (news-only, social-only, unified LLM); highlight which blocs benefit most
  from the social media signal. (3) \textit{Shock Simulations} --- present
  win probability and bloc-weight shifts for 3--4 representative historical
  shocks (one per category); include the sensitivity analysis plot for the
  Ayatollah event. (4) \textit{Platform Correlations} --- one paragraph
  reporting the per-platform correlation with 14-day lagged favorability
  polls.}{4}
\end{taskbox}

% ── Day 4 ────────────────────────────────────────────────────────────────────
\dayhead{4}{Thu --- Paper Completion \& Poster}

\begin{taskbox}
\done \task{Write the \textbf{Introduction} (target: 600 words). Open with
  a concrete example: the Ayatollah shock and its measurable effect on
  evangelical favorability. State the three gaps in existing work: (i)
  polling aggregators are static; (ii) advanced systems are enterprise-only;
  (iii) religious and demographic granularity is underused. State the
  three contributions: (i) the voter equilibrium framework; (ii) the
  two-stage RoBERTa-to-LLM pipeline with social media augmentation;
  (iii) the public web application. End with a roadmap of the paper.}{2}

\done \task{Write the \textbf{Literature Review} (target: 700 words).
  Four subsections: (1) \textit{Electoral Forecasting} --- FiveThirtyEight,
  Economist model, prediction markets; their shared limitation as polling
  aggregators. (2) \textit{Religion and Race in Presidential Politics} ---
  cite Prof.\ Espinosa's series (all five books, key articles); the Evangelical
  threshold, Latino Pentecostal patterns, Catholic swing behavior.
  (3) \textit{NLP in Political Science} --- sentiment analysis for event
  detection, existing social media political forecasting work.
  (4) \textit{Portfolio Theory in Non-Financial Contexts} --- precedents
  for applying mean-variance optimization outside finance.}{2}

\done \task{Write the \textbf{Discussion and Limitations} (500 words) and
  \textbf{Conclusion} (300 words). Discussion must address: (i) platform
  selection bias and the baseline-adjustment mitigation; (ii) small
  fine-tuning dataset and why QLoRA regularizes against overfitting;
  (iii) the model should display uncertainty prominently and must not be
  used for campaign strategy. Conclusion: summarize contributions, state
  the 2026 midterm validation plan, note the open-source release.}{2}

\done \task{Design the conference poster ($48'' \times 36''$) in Canva.
  Sections: (i) title block with authors and CMC; (ii) pipeline diagram
  (the full flow from data to web app); (iii) the baseline comparison table
  (\texttt{roberta\_news\_only} vs.\ \texttt{roberta\_social\_only} vs.\
  \texttt{llm\_unified} MAE per bloc); (iv) a \texttt{CoalitionChart} screenshot for the
  Ayatollah event; (v) the win gauge for one Electoral Surprise shock;
  (vi) a dashboard screenshot (Panel 1 coverage matrix or Panel 6 audit log);
  (vii) QR code linking to the live Modal + Vercel URL. Export as PDF and PNG.
  Share with Prof.\ Espinosa for any domain input before finalizing.}{2}
\end{taskbox}

% ── Day 5 ────────────────────────────────────────────────────────────────────
\dayhead{5}{Fri --- Deployment, Final Submission \& Retrospective}

\begin{taskbox}
\done \task{Incorporate any domain feedback from the Week~8 supervision
  check-in. Finalize paper and poster as PDFs. Submit both to the SRP
  committee per their submission instructions.}{2}

\done \task{Verify the Modal + Vercel deployment is stable: run 20 back-to-back
  inference calls from the public URL, confirm all return valid JSON, confirm
  mean latency is under 5 seconds (Modal cold start included). If latency is
  too high, switch Modal to a \texttt{keep\_warm=1} container to eliminate
  cold starts. Document final latency numbers in \texttt{docs/latency\_benchmark.md}.}{1}

\done \task{Verify the CMC HPC hosting scaffold: \texttt{sbatch deploy/hpc\_vllm.sh}
  on CMC HPC and confirm the vLLM server starts and the health check passes on
  the HPC-internal URL. Do not expose externally yet --- this is a dry run to
  confirm the scaffold works. Log the internal URL and SLURM job ID in
  \texttt{docs/hpc\_deployment\_notes.md} for future activation.}{1}

\done \task{Pin all dependency versions: run \texttt{pip freeze >
  requirements.txt} in the activated venv. Update \texttt{pyproject.toml}
  to match. Commit. Confirm a fresh install from \texttt{requirements.txt}
  passes the full test suite.}{1}

\done \task{Write \texttt{CONTINUOUS\_UPDATE.md} --- the post-SRP maintenance
  specification. This document ensures the model does not become stale after
  the project ends. Required sections:

  \begin{enumerate}[leftmargin=1.6em, itemsep=2pt]
    \item \textbf{Nightly data ingestion.} A launchd plist on the
          \textbf{Intel Mac} triggers a Prefect flow at 2am using
          \texttt{pipeline\_mode="continuous"}: runs \texttt{BlueSkyCollector}
          and \texttt{ApifyCollector} for any shock events in the last 7 days;
          runs \texttt{sample\_archives.py} on any new archive data;
          runs \texttt{clean\_with\_llm.py} on the new sample;
          scores new posts via RoBERTa; appends to
          \texttt{data/finetune/new\_events.jsonl}. The Intel Mac is the right
          machine for this: it is always on, already running collectors, and
          can submit USC Laguna fine-tuning jobs when the threshold is reached without
          any human intervention. Include the launchd plist and the
          \texttt{just continuous} Justfile target.
    \item \textbf{Fine-tuning trigger.} When \texttt{new\_events.jsonl} reaches
          10+ new shock events, submit a new USC Laguna fine-tuning job. Archive the
          current adapter with a timestamp before replacing it.
    \item \textbf{Voter panel refresh.} Annual check: ARDA, GSS, Pew, and Gallup
          all release updated data. When new data is available, re-run the
          Week~2 pipeline. The win threshold $\veq^{(P)}$ and constraint bounds
          update automatically via \texttt{build\_constraint\_spec}.
    \item \textbf{Prediction market calibration.} Quarterly: query
          \texttt{audit.duckdb} for resolved shock events; compare
          \texttt{win\_probability} predictions against actual market outcomes.
          Document any systematic divergence and adjust fine-tuning example
          weights accordingly.
    \item \textbf{How to trigger a manual run.} Commands for \texttt{just
          continuous}, \texttt{just train}, and \texttt{just deploy}.
  \end{enumerate}}{2}

\done \task{Create \texttt{configs/continuous.json}: the configuration file for
  nightly incremental runs. Set \texttt{pipeline\_mode: "continuous"},
  \texttt{lookback\_days: 7} (only collect shocks from the last 7 days),
  \texttt{skip\_voter\_panel\_rebuild: true} (voter panel only rebuilt
  annually), \texttt{finetune\_trigger\_threshold: 10} (trigger retraining
  after 10 new shock events). Add \texttt{just continuous} to the
  \texttt{Justfile}: \texttt{prefect run pipeline.py --config
  configs/continuous.json}.}{1}

\done \task{Final PR: \texttt{feature/week8-hardening-deliverables}. This PR
  must include: (i) all integration and hardening tests, (ii) the frozen
  paper-baseline artifacts, (iii) \texttt{requirements.txt} with pinned
  versions, (iv) \texttt{scripts/verify\_artifacts.py}, (v) the complete
  \texttt{README.md}, (vi) the \texttt{Justfile} with all targets verified
  to work on a clean checkout including \texttt{just continuous},
  (vii) \texttt{DOMAIN.md} reviewed and complete,
  (viii) \texttt{CONTINUOUS\_UPDATE.md} complete.
  Merge to \texttt{main}. Tag as \texttt{v1.0}.}{1}

\done \task{Write \texttt{docs/retrospective.md}: a one-page honest assessment
  covering (i) what went as planned, (ii) what took longer than expected
  and why (likely: API approvals, HPC queue wait times, data cleaning),
  (iii) what you would do differently with a 10-week timeline, (iv) the
  concrete publication plan (target journal, submission timeline), (v) the
  2026 midterm validation strategy (which shocks to monitor, how to
  update the model with new data).}{2}

\done \task{Final supervision check-in with Prof.\ Espinosa: project wrap-up.
  Share the submitted paper and poster. Get domain sign-off on the framing
  and discuss the 2026 midterm validation plan. Agree on publication
  timeline and the 2--3 most important follow-up directions.}{2}
\end{taskbox}

% ════════════════════════════════════════════════════════════════════════════
\newpage
\section{Week 8 --- Debugging \& Hardening}
% ════════════════════════════════════════════════════════════════════════════

This week consolidates debugging and hardening deferred during the Week~5--6
feature sprints. Tasks are ordered by whether they block the SRP deliverable
(training must complete; model must be directionally correct) versus polish.
Infrastructure and training pipeline fixes are prerequisites; model quality
and data pipeline tasks can proceed in parallel once training is unblocked.

\begin{taskbox}{Standing Process --- Flag Capture}

\task{Maintain a running flag log for the remainder of the project. Every
  time a development session surfaces a bug, an unverified assumption, an
  ``it should work but isn't confirmed'', a ``decide with Prof.\ Espinosa'',
  or a ``this contradicts that'' --- it gets appended here as a checkbox
  item with a one-line description and a severity (blocking / correctness /
  hardening / polish). Nothing gets fixed silently and nothing gets
  forgotten by living only in a chat transcript. New flags append to the
  end; resolved flags are checked off, not deleted, so the audit trail
  survives.}{1}

\task{Flag-log entries captured so far (append below as new ones arise):
  \begin{itemize}[leftmargin=1.6em, itemsep=1pt]
    \item[$\square$] \textbf{[correctness]} Meaning of \texttt{target}
      (0.5066) undefined --- popular vote? EC threshold? Drives
      win\_prob=1.0 saturation. Resolve with Espinosa.
    \item[$\square$] \textbf{[blocking]} Retrained model must be confirmed
      to fix the Newsom all-minorities-slight\_pos directional error.
    \item[$\square$] \textbf{[correctness]} Party-neutral preset text +
      separate \texttt{party} field: confirm model reads the field and
      flips deltas by party.
    \item[$\square$] \textbf{[correctness]} ShockNarrative (pre-rebalance
      sentiment) vs CoalitionChart (post-rebalance weights) must not read
      as contradictory.
    \item[$\square$] \textbf{[blocking]} Laguna bitsandbytes hang ---
      bf16 fallback written but never confirmed working.
    \item[$\square$] \textbf{[hardening]} Hopper 50\,GB quota: HF cache,
      checkpoint bloat, 3.8\,GB \texttt{.git} with committed weights.
    \item[$\square$] \textbf{[correctness]} Confirm Ledoit-Wolf covariance
      used in production, not the degenerate diagonal fallback.
    \item[$\square$] \textbf{[correctness]} No artifact may serialize
      \texttt{Infinity}/\texttt{NaN} --- shared sanitizer + test.
    \item[$\square$] \textbf{[correctness]} Apply + verify all six Copilot
      PR fixes (binomial bootstrap, field comments, dual artifact
      filenames $\times$3, fallback method label).
    \item[$\square$] \textbf{[hardening]} Green test suite on a clean
      checkout after mid-development schema churn.
    \item[$\square$] \textbf{[hardening]} Cross-machine hardcoded paths
      audit (Mac / Laguna / Hopper).
    \item[$\square$] \textbf{[hardening]} 647-record Gemini human-review
      queue: adjudicate or explicitly defer.
    \item[$\square$] \textbf{[polish]} candidates.jsonl overwrite guard;
      DeepSeek truncation handling; malformed-JSONL defensive parsing.
    \item[$\square$] \textbf{[hardening]} Backfill or document 60--66
      missing (cycle, bloc) survey panel cells.
    \item[$\square$] \textbf{[correctness]} Data-provenance audit: every
      ``real'' number traces to a documented source.
    \item[$\square$] \textbf{[correctness]} CoalitionChart dual-axis
      overlay: loyalty and weight occupy distinct visual ranges on a shared
      [0,1] axis --- verify no per-bloc collision on real shock data before
      shipping as the SRP headline figure.
    \item[$\square$] \textbf{[correctness]} WinGauge model-vs-market
      juxtaposition: model answers a counterfactual-conditional question,
      market answers an unconditional-current question --- UI must label
      them as different, and must flag when the matched contract is the
      generic election market (not event-conditioned). Confirm framing with
      Espinosa.
    \item[$\square$] \textbf{[correctness (reporting)]} SSE streaming
      ``30\% latency improvement'' claim is wrong --- stages are sequential
      so streaming improves time-to-first-content only, not total
      round-trip. Update devplan/docs to claim perceived-latency
      improvement, not throughput speedup.
  \end{itemize}
  The SRP is not ``done'' until every box here is checked or has a written,
  Espinosa-approved deferral.}{2}

\end{taskbox}

\begin{taskbox}{HPC \& Infrastructure}
\task{Resolve USC Laguna \texttt{bitsandbytes} hang. Laguna GPU nodes (both c22
  and c23) run kernel 4.18.0, below the 5.5.0 minimum. QLoRA 4-bit
  quantization via \texttt{bitsandbytes} hangs indefinitely at training step~0
  on every Laguna node. Implement and TEST the \texttt{ELECTORAL\_NO\_4BIT=1}
  bf16 fallback path in \texttt{trainer.py} (skip \texttt{BitsAndBytesConfig},
  load in bfloat16 --- L40S 46\,GB fits 7B bf16 + LoRA easily). The fallback was
  written but never confirmed working because the env var did not propagate to
  the trainer in the last attempt. Verify the bf16 branch is actually taken
  (grep the log for absence of the bitsandbytes 4-bit loading line) and that
  training advances past step~0.}{3}

\task{Fix Hopper disk quota management. The 50\,GB home quota is repeatedly
  exhausted. Root causes: (i)~14\,GB HuggingFace Mistral-7B cache in
  \texttt{\textasciitilde{}/.cache/huggingface}, (ii)~intermediate training
  checkpoints accumulating ($\sim$13\,MB each, every 113 steps),
  (iii)~a bloated 3.8\,GB \texttt{.git} directory containing committed model
  weights. Fix all three: (a)~set \texttt{HF\_HOME} and
  \texttt{TRANSFORMERS\_CACHE} to node-local \texttt{/tmp} during runs,
  (b)~add \texttt{save\_total\_limit=1} to \texttt{TrainingArguments} in
  \texttt{trainer.py}, (c)~run \texttt{git gc} / filter the \texttt{.git}
  history to purge committed weights.}{3}

\task{Purge model weights from git history. The \texttt{.git} directory is
  3.8\,GB because model safetensors were accidentally committed before
  \texttt{models/} was gitignored. Use \texttt{git filter-repo} or BFG to
  remove \texttt{models/} and \texttt{checkpoints/} from all history, then
  force-push. Coordinate timing so Hopper and Laguna clones are re-cloned
  cleanly afterward. Verify \texttt{.git} drops below 50\,MB.}{2}

\task{Document the cross-HPC dependency matrix. The working dependency set
  differs by machine and must be pinned:
  \texttt{outlines==0.0.37}, \texttt{pyairports},
  \texttt{numpy\textless{}2.0}, \texttt{scipy\textless{}1.14},
  \texttt{cvxpy==1.3.4}, \texttt{transformers==4.46.3}.
  The \texttt{outlines}/\texttt{cvxpy}/\texttt{numpy} versions conflict
  (\texttt{cvxpy\textgreater{}=1.7.5} wants \texttt{numpy\textgreater{}=2.0}
  but \texttt{outlines} 0.0.37 wants \texttt{numpy\textless{}2.0}).
  Document the resolved pinned set in a \texttt{requirements-hpc.txt} and
  in \texttt{DECISIONS.md}, noting Hopper requires
  \texttt{module load cuda/13.2} and
  \texttt{LD\_LIBRARY\_PATH=/hopper/software/cuda/13.2/lib64} in every
  \texttt{srun}.}{2}
\end{taskbox}

\begin{taskbox}{Training Pipeline}
\task{Resolve \texttt{eval\_mae=inf} during training. Both \texttt{r=16} and
  \texttt{r=32} training runs report \texttt{eval\_mae=inf} because the 4-bit
  model's in-loop evaluation is broken. Currently MAE is measured post-hoc
  via \texttt{run\_baseline\_comparison.py}. Either fix the in-loop eval to
  produce a real number or formally remove the broken eval from the training
  loop and rely on the post-hoc script. Replace any \texttt{Infinity} written
  to \texttt{trainer\_state.json} with \texttt{null} (JSON-invalid
  \texttt{Infinity} already caused a downstream parse bug, patched once ---
  verify \texttt{trainer.py} no longer emits it).}{2}

\task{Fix \texttt{outlines} constrained-decoding latency.
  \texttt{ShockEstimator.estimate()} takes $\sim$12\,s per call (p95 well
  above the 3\,s target) because \texttt{outlines} 0.0.37 FSM-based
  constrained decoding processes $\sim$200 tokens one-at-a-time. Documented
  in \texttt{docs/latency\_benchmark.md} as deferred. Investigate \texttt{vLLM}
  \texttt{guided\_json} decoding (expected 2--4\,s) or a distilled smaller
  model. This blocks any interactive/web use of the endpoint.}{4}

\task{Fix the KV-cache / \texttt{use\_cache} workaround. \texttt{ShockEstimator}
  currently sets \texttt{model.config.use\_cache=False} to avoid a
  \texttt{NoneType.device} crash in \texttt{outlines}' \texttt{reorder\_kv\_cache}
  under PEFT models. This disables KV caching and worsens latency. Find a
  proper fix that keeps caching enabled.}{3}

\task{Audit the intensity-clipping behavior. \texttt{estimate()} clips deltas
  to $[-0.15,\, 0.15]$ \emph{after} intensity scaling to satisfy
  \texttt{ShockResponseData} validation. At intensity $= 1.5$--$2.0$ this
  clips many blocs to the boundary, flattening the sensitivity analysis.
  Decide whether clipping should happen pre- or post-scaling and document
  the modeling justification.}{2}
\end{taskbox}

\begin{taskbox}{Model Quality}
\task{Validate the retrained \texttt{r=16} model fixes the directional errors.
  The pre-augmentation model predicted \texttt{slight\_pos} for \emph{all}
  minority blocs on the Newsom 2028 nomination event (politically wrong ---
  Newsom is expected to hurt with white working-class and some Latino voters).
  After retraining on the 1,136-record augmented dataset, rerun the Newsom
  sensitivity analysis and confirm white/latino now show the expected negative
  or mixed deltas. If still wrong, the synthetic labels did not fix the
  underlying issue --- escalate to a data-quality review with Prof.\ Espinosa
  rather than generating more volume.}{3}

\task{Verify party-neutral event text + explicit party field produces
  party-appropriate deltas. The webapp uses an authoritative party toggle
  with party-neutral preset templates (e.g.\ ``the leading presidential
  candidate'', ``the sitting administration'') --- the party is passed as a
  separate field in the event dict, not named in the description. But the
  model was fine-tuned on synthetic seeds that DO name the party in the
  text (``Democratic nominee\ldots{}'', ``Republican administration\ldots{}'').
  Confirm the model reads the separate \texttt{party} field correctly: send
  the same neutral event text with \texttt{party="democrat"} and
  \texttt{party="republican"} and verify the predicted deltas flip in the
  politically expected direction. If they do NOT (i.e.\ the model ignores
  the party field when the text is neutral), the fix is to have the
  frontend inject the party name into the event text at submit time
  (template ``the \{party\} candidate'') rather than relying on the field.
  Decide and document which approach is authoritative.}{2}

\task{Resolve the \texttt{win\_probability=1.0} saturation. With NEP-weighted
  $w_0$, the Democrat baseline $\mu_{\mathrm{eff}}$ is 0.605 vs.\ target
  0.5066 --- a 0.098 gap that makes \texttt{win\_probability} saturate at 1.0
  across all shock intensities. The diagonal fallback covariance (near-zero
  variance, used when fewer than 3 panel cycles) produces degenerate CIs of
  $[1.0,1.0]$ or $[0.0,0.0]$. Confirm the Ledoit-Wolf covariance from 14 real
  cycle-deltas is actually being used in production (not the fallback), and
  verify the target value 0.5066 is the correct Electoral College threshold ---
  Democrats are structurally disadvantaged and may need $\sim$0.52--0.54
  popular-vote-equivalent to win the EC. Clarify what ``target'' represents
  with Prof.\ Espinosa.}{3}

\task{Verify the NEP-weighted baseline formula end-to-end. The $w_0$
  computation was changed from a degenerate CVXPY-derived result
  (\texttt{asian}$=0$, \texttt{latino}$=0$, \texttt{african\_american}$=0.51$)
  to $w_i = (\mathrm{NEP}_i \cdot \mu_i) / \sum_j(\mathrm{NEP}_j \cdot \mu_j)$
  using Edison Research 2024 exit-poll voter shares. Confirm
  \texttt{build\_baseline\_portfolio} uses this formula (not
  \texttt{solve\_baseline}) for $w_0$, and that the resulting weights
  (\texttt{white}$=0.528$, \texttt{african\_american}$=0.189$,
  \texttt{latino}$=0.180$, \texttt{asian}$=0.051$,
  \texttt{other\_race}$=0.052$) propagate correctly into the optimizer and
  Monte Carlo. \texttt{solve\_baseline} should be used \emph{only} for
  post-shock rebalancing.}{2}
\end{taskbox}

\begin{taskbox}{Data Pipeline}
\task{Fix the \texttt{candidates.jsonl} overwrite hazard. The DeepSeek
  generator opens the output file in write mode, so an accidental second run
  silently overwrote 1,237 good records with 0 (the run kept=0 because the
  file already existed at the seed level). Add a guard: refuse to overwrite a
  non-empty output file without \texttt{-{}-force}, or append with dedup.
  Also reduce per-attempt warning noise (warn only on final failure).}{2}

\task{Fix DeepSeek truncation at scale. Generating 20 records per call
  truncated mid-record at the 4k token limit (kept=0 for several seeds).
  Worked around with \texttt{BATCH\_SIZE=5} and \texttt{max\_tokens=8000}.
  Make batch size and token budget configurable and add a post-parse check
  that warns when a response appears truncated (ends mid-object).}{1}

\task{Handle malformed JSONL in review outputs. 7 of 647 human-queue records
  are unparseable JSON (Gemini occasionally emits invalid output). All readers
  must parse defensively (skip-and-count bad lines) rather than crashing on
  \texttt{json.loads}. Audit every place that reads these JSONL files.}{1}

\task{Backfill missing survey panel cells. \texttt{build\_voter\_panel} reports
  60--66 missing (cycle, bloc) pairs across ANES/CES/VOTER sources --- notably
  evangelical, secular, jewish, muslim, \texttt{other\_gender} absent for
  pre-2004 cycles, and most blocs missing for 2024. The religion strata impute
  to 0.50 with warnings. Source the missing data or document the imputation as
  a known limitation with its effect on $\mu_{\mathrm{religion}}$.}{3}
\end{taskbox}

\begin{taskbox}{Synthetic Review Backlog}
\task{Adjudicate the 647-record human review queue. Gemini 2.5 Flash flagged
  647 of 1,237 synthetic records as \texttt{HUMAN\_REVIEW} (conservative by
  design; 579 were auto-approved and already merged, bringing the training set
  to 1,136). Process \texttt{human\_review\_queue.jsonl} in an Opus session
  (manual upload during dormant hours): approve, correct, or reject each. The
  queue includes $\sim$15 records Gemini actually approved but were routed via
  the 5\% spot-check mechanism --- these validate that approved records are
  trustworthy. Merge keepers and retrain.}{4}

\task{Spot-check the merged synthetic set for shared-blindspot errors. Three
  models (DeepSeek generate, Gemini review, the prior Mistral) can share the
  same political blind spots --- e.g.\ all assuming a far-left event helps
  Democrats. Sample 50 high-confidence approved records plus the
  highest-disagreement tail and review with Prof.\ Espinosa for systematic
  directional bias before the labels are trusted as ground truth.}{3}

\task{Add coverage-gap detection to the generator. Review accuracy does not
  fix coverage gaps --- if DeepSeek under-produces a taxonomy cell
  (e.g.\ \texttt{helps\_rep}/\texttt{far\_right} institutional events), no
  amount of review surfaces the absence. Add per-axis-cell counting to
  \texttt{generate\_deepseek.py} and oversample underfilled cells in a second
  pass.}{2}
\end{taskbox}

\begin{taskbox}{API \& Frontend Contract}

\task{Widen the \texttt{/estimate} endpoint to return the full artifact
  bundle. The webapp's \texttt{EstimateResponse} type assumes the endpoint
  returns \texttt{shock}, \texttt{equilibrium}, and \texttt{simulation}
  together, but \texttt{shock\_endpoint.py} currently returns only
  \texttt{ShockResponseData}. Extend the POST \texttt{/estimate} route to
  run the full chain (\texttt{build\_shock\_response} →
  \texttt{solve\_rebalanced} → \texttt{run\_ilr\_montecarlo}) and return a
  composite JSON object \texttt{\{shock, equilibrium, simulation\}} matching
  the \texttt{EstimateResponse} TypeScript interface. Add per-stage timing
  fields (\texttt{llm\_ms}, \texttt{optimizer\_ms}, \texttt{montecarlo\_ms})
  so the frontend \texttt{AuditEntry} can be populated from a single call.
  Alternative considered and rejected: three separate endpoints --- adds
  round-trips and state-sync complexity. Keep it one call.}{3}

\task{Establish a single source of truth for the API contract. The
  TypeScript types in \texttt{webapp/lib/types.ts} are hand-mirrored from
  the Python dataclasses in \texttt{electoral/artifacts.py} and will drift.
  Either (a) generate the TS types from the Pydantic
  \texttt{ShockResponseSchema} via \texttt{pydantic-to-typescript} or a
  JSON Schema export, or (b) add a CI check that fails when the two
  definitions diverge. Cover \texttt{ShockResponseData},
  \texttt{EquilibriumData}, \texttt{SimulationData} (including the
  \texttt{win\_probability\_low}/\texttt{high} fields added in Week~6), and
  the \texttt{Party}/\texttt{DeltaBin} literal types. Document the chosen
  approach in DECISIONS.md.}{3}

\task{Reconcile the latency budget with the endpoint design. The full
  \texttt{/estimate} chain includes the $\sim$12\,s outlines constrained
  decoding (documented in \texttt{docs/latency\_benchmark.md}), so a
  single composite call will take $\sim$12--15\,s end-to-end. The frontend
  must handle this: add a loading/streaming state, and decide whether
  Monte Carlo (the fast part) should stream separately from the LLM
  estimate (the slow part). Cross-reference the Week~8 latency task
  (vLLM \texttt{guided\_json}) --- if that lands, this budget shrinks to
  $\sim$3--5\,s. Until then the UI must not appear hung during the 12\,s
  wait.}{2}

\task{Define the \texttt{TrainingRun.val\_loss} and \texttt{mae} null
  contract. These webapp fields are nullable because in-loop evaluation
  currently produces \texttt{eval\_mae=inf} (the 4-bit eval bug). The
  frontend must render ``n/a'' gracefully when these are null, and the
  \texttt{TrainingRun} feed must serialize Python \texttt{inf}/\texttt{NaN}
  as JSON \texttt{null} (never the invalid \texttt{Infinity} literal that
  already caused one parse bug). Add a serialization guard.}{2}

\task{Populate \texttt{CoverageCell} quality flags from the real panel
  data. The webapp coverage grid expects per-(bloc, cycle) quality of
  ``real'' / ``imputed'' / ``missing'', but the panel builder currently only
  logs missing pairs as warnings (60--66 missing cells, mostly religion
  strata pre-2004 and most blocs in 2024). Emit a structured coverage
  artifact from \texttt{build\_voter\_panel} that the endpoint can serve,
  rather than the frontend guessing.}{2}

\task{Implement the \texttt{/estimate/stream} SSE endpoint in
  \texttt{shock\_endpoint.py}. The frontend \texttt{page.tsx} opens an
  \texttt{EventSource} expecting named events \texttt{deltas},
  \texttt{equilibrium}, \texttt{simulation}, and \texttt{done}, emitted in
  that order as each pipeline stage completes. Use FastAPI
  \texttt{StreamingResponse} with \texttt{media\_type="text/event-stream"}.
  Yield the \texttt{deltas} event as soon as
  \texttt{ShockEstimator.estimate()} returns ($\sim$12\,s), then
  \texttt{equilibrium} after \texttt{solve\_rebalanced} ($\sim$ms), then
  \texttt{simulation} after \texttt{run\_ilr\_montecarlo}, then
  \texttt{done}. This is what makes the progressive-render UX work ---
  without it the user stares at a blank dashboard for the full 12\,s. Note:
  \texttt{EventSource} only supports GET, so this endpoint takes query
  params (\texttt{event}, \texttt{intensity}, \texttt{party}), unlike the
  POST \texttt{/estimate}.}{3}

\task{Reconcile the two endpoint shapes. The webapp now depends on BOTH
  POST \texttt{/estimate} (returns the full bundle, validated by
  \texttt{EstimateResponseSchema}) and GET \texttt{/estimate/stream} (SSE,
  incremental). Decide whether both are maintained or the non-streaming
  one is dropped. If both stay, the Zod \texttt{EstimateResponseSchema}
  should also validate the assembled-from-stream object so the streamed
  path gets the same runtime safety as the single-shot path.}{2}

\task{Verify ShockNarrative and CoalitionChart tell the same story. The
  narrative renders plain-language ``hurt/help'' sentences from raw
  \texttt{delta\_bins} (the model's pre-rebalance sentiment prediction),
  while CoalitionChart shows both the baseline bars AND the
  post-optimization rebalanced overlay. These are two different layers ---
  sentiment movement vs.\ campaign resource reallocation --- and must not
  be read as contradicting each other. Confirm: (i) the narrative
  explicitly describes direction-of-movement (``standing with''), never
  final level (``the party is winning X''); (ii) a code comment marks
  ShockNarrative as intentionally pre-rebalance so no one later rewires it
  to equilibrium weights; (iii) for a test shock, the blocs the narrative
  says are ``hurt'' are not simultaneously shown gaining weight in the
  chart in a way a layperson would read as contradictory --- and if they
  can be (a bloc whose sentiment drops but whose optimal campaign weight
  rises), add a one-line chart caption explaining the distinction. The
  plain-language layer is the one a reader takes at face value, so it must
  be faithful to exactly what the number means.}{2}

\task{Pressure-test the CoalitionChart dual-axis overlay on real result
  data. The chart overlays within-bloc loyalty ($\mu_i$, $\sim$0.15--0.90)
  and coalition weight ($w_i$, $\sim$0.03--0.30) on one shared [0,1] axis,
  relying on them occupying distinct ranges to stay visually unambiguous.
  This breaks when a bloc's loyalty and weight land near the same value
  (e.g.\ White loyalty 0.30 vs weight 0.28) --- the two bars overlap and a
  layperson reads the translucent bar as a loyalty drop, not a different
  quantity. Render the chart for the 3--4 most consequential real shocks
  and check for this collision. If it occurs, either (a) split into two
  stacked sub-charts sharing a Y axis (loyalty panel + weight panel), or
  (b) move weight to a secondary top X axis with its own scale and label.
  Do NOT ship the shared axis as the SRP headline figure until this is
  verified on real numbers. Severity: correctness --- the figure must not
  mislead.}{2}

\task{Make the model-vs-market juxtaposition epistemically honest. The
  WinGauge shows model $P(\text{win} \mid \text{hypothetical shock})$ next
  to market $P(\text{win} \mid \text{everything currently true})$. These
  answer DIFFERENT questions --- the model is counterfactual-conditional,
  the market is unconditional-current --- and for most hypothetical events
  no contract is actually conditioned on the event, so ``nearest relevant
  contract'' returns the generic election market (the unconditional prior).
  Presenting ``Model 62\% / Market 55\%'' as if they were the same question
  is misleading. Required: (i) UI copy must label them as different
  questions, not competing estimates --- e.g.\ ``Model: if this happens
  $\to$'' vs ``Market (current, unconditional): $\to$''; (ii) when the
  matched contract is NOT conditioned on the event (the common case), say
  so explicitly (``nearest market is the general election winner, not this
  scenario''); (iii) the \texttt{aggregator.py} match must return whether
  the contract is event-conditional or generic, and the UI must surface
  that. Severity: correctness --- this is the headline epistemic display
  and must not imply a false comparison. Confirm framing with
  Prof.\ Espinosa.}{2}

\end{taskbox}

\begin{taskbox}{Final Sweep --- Ambiguities \& Loose Ends}

\task{Resolve the meaning of \texttt{target}. The win threshold is set to
  0.5066 but it is unclear whether this represents the national two-party
  popular-vote share needed to win the Electoral College, a weighted
  effective-loyalty threshold, or something else. Democrats are
  structurally disadvantaged in the EC (need $\sim$52--54\% popular vote),
  so 0.5066 may be too low and is the likely cause of the
  win\_probability=1.0 saturation. Pin down the exact definition with
  Prof.\ Espinosa, document it in DECISIONS.md, and adjust the value if
  the current one is wrong.}{2}

\task{Apply all six Copilot PR review fixes that were accepted but may not
  all be committed: (i) Bernoulli-binomial bootstrap in montecarlo.py
  (avoid the 50k$\times$500 index matrix), (ii) corrected
  win\_probability\_low/high field comments in artifacts.py (bootstrap of
  the estimator, not ``across ILR draws''), (iii)+(iv)+(v) shock.py must
  write BOTH legacy \texttt{shock\_response.json}/\texttt{equilibrium.json}
  AND per-shock \texttt{shock\_\{id\}.json}/\texttt{equilibrium\_\{id\}.json},
  restoring the ``event'' metadata key (fixes test\_stages.py and
  test\_shock\_kernel.py), (vi) the equal-weight fallback in cvx.py must
  use \texttt{method="equal\_weight\_fallback"}, not ``cvxpy\_dqcp''. Verify
  each is committed and its test passes.}{3}

\task{Run the full test suite on a clean checkout and fix every failure.
  Many tests were written against schemas and artifact filenames that
  changed mid-development (SimulationData gained win\_probability fields;
  shock.py changed artifact names; cvx fallback method label changed).
  Enumerate failures, fix the code or the test as appropriate, and
  confirm a green suite. Do not leave any test xfail/skip without a
  documented reason.}{3}

\task{Reconcile every cross-machine path and config assumption. The repo
  has hardcoded paths for three environments (M5 Mac \texttt{.venv} Py3.9,
  Laguna scratch, Hopper home). Audit for hardcoded \texttt{/Volumes/},
  \texttt{/Users/juhidamley/}, \texttt{/scratch/\ldots},
  \texttt{/hopper/home/\ldots} that will break when run elsewhere. Move them
  to config or env vars. Confirm \texttt{configs/base.json}
  vs \texttt{smoke.json} data\_path behavior is documented.}{2}

\task{Confirm the production covariance path. Verify
  \texttt{run\_ilr\_montecarlo} uses the Ledoit-Wolf $\Sigma_\Delta$ from
  14 real cycle-deltas in production runs, not the near-zero diagonal
  fallback (which produces degenerate [1.0,1.0]/[0.0,0.0] CIs). Add an
  explicit log line stating which covariance source was used, and a test
  asserting the fallback is NOT silently used when real panel data is
  present.}{2}

\task{Final data-provenance audit. Confirm every ``real'' number traces to
  a documented source: NEP 2024 voter shares (Edison Research), ANES/CES
  panel mu values, the 1,136-record synthetic training set (577 original
  + 559 Gemini-approved, with the 647-record human-review queue either
  adjudicated or explicitly marked as deferred). No silently-imputed value
  should be presentable as empirical. Produce a one-page provenance table
  for the SRP appendix.}{3}

\task{Verify no JSON artifact can contain \texttt{Infinity}/\texttt{NaN}.
  The \texttt{eval\_mae=inf} bug already produced one invalid-JSON parse
  failure. Audit every \texttt{write\_artifact} / \texttt{to\_dict} path
  and the \texttt{TrainingRun} serializer to guarantee non-finite floats
  become \texttt{null}. Add a single shared sanitizer and a test feeding
  it \texttt{inf}/\texttt{-inf}/\texttt{nan}.}{2}

\end{taskbox}

% ════════════════════════════════════════════════════════════════════════════
\newpage
\section*{AI Research Assistant --- Quick Reference}
% ════════════════════════════════════════════════════════════════════════════

\textbf{Setup:} OpenClaw on Raspberry Pi 5 $\to$ Telegram/WhatsApp interface $\to$
Gemini Pro (default) or Claude Sonnet (\texttt{@claude} prefix).
Syncthing shared folder across the three local machines (M5, Intel Mac, Pi).

\subsection*{Machine Allocation at a Glance}

\begin{center}
\begin{tabular}{p{0.28\linewidth}p{0.32\linewidth}p{0.32\linewidth}}
\toprule
\textbf{Machine} & \textbf{Always running} & \textbf{When needed} \\
\midrule
M5 MacBook Pro & Development, Claude sessions & MLX fine-tune experiments \\
Intel MacBook & Bluesky + Apify collectors (24/7); nightly pipeline (2am launchd) & FastAPI endpoint in Week~7 \\
Intel Mac & News scraper (was Windows role), collectors & launchd 2am nightly pipeline \\
Pi 5 + AI HAT & Bio classifier NPU, OpenClaw & --- \\
USC Laguna HPC & --- & Fine-tuning, scoring, Monte Carlo (compute) \\
CMC HPC & --- & vLLM hosting \& deployment (hosting only) \\
\bottomrule
\end{tabular}
\end{center}

\subsection*{When to use Gemini vs.\ Claude}

\begin{center}
\begin{tabular}{p{0.35\linewidth}p{0.12\linewidth}p{0.45\linewidth}}
\toprule
\textbf{Task} & \textbf{Model} & \textbf{Example message} \\
\midrule
Summarize a paper & Gemini & ``Summarize arxiv.org/abs/2401.00000'' \\
Morning project brief & Gemini & Automated 8am cron \\
Quick code question & Gemini & ``What does \texttt{scale\_pos\_weight} do in XGBoost?'' \\
Find references & Gemini & ``Find 3 papers on NLP and electoral forecasting'' \\
Draft a paragraph & Gemini & ``Draft an intro paragraph on portfolio theory'' \\
\midrule
Review pipeline logic & Claude & ``@claude review kernels/shock.py'' \\
Debug test failure & Claude & ``@claude debug this pytest output: [paste]'' \\
Write or fix LaTeX & Claude & ``@claude fix the overflows in section 4'' \\
Architecture decision & Claude & ``@claude should I use Ridge or Lasso for the elasticity regression?'' \\
Full session work & Claude & Open claude.ai directly \\
\bottomrule
\end{tabular}
\end{center}

\subsection*{Useful commands}

\begin{center}
\begin{tabularx}{\linewidth}{p{0.42\linewidth}X}
\toprule
\textbf{Command} & \textbf{Action} \\
\midrule
\texttt{openclaw doctor --fix} & Validate config and fix common issues \\
\texttt{openclaw skills install <name>} & Install a skill from ClawHub \\
\texttt{openclaw logs} & View recent agent activity \\
\texttt{openclaw config edit} & Open config in editor \\
\texttt{@claude <message>} & Route this message to Claude Sonnet \\
\texttt{@gemini <message>} & Explicitly route to Gemini Pro \\
\bottomrule
\end{tabularx}
\end{center}

\subsection*{Cost guardrails --- spend only where it matters}

\begin{center}
\begin{tabularx}{\linewidth}{p{0.22\linewidth}p{0.16\linewidth}p{0.22\linewidth}X}
\toprule
\textbf{Service} & \textbf{Cost} & \textbf{Programme} & \textbf{Action} \\
\midrule
USC Laguna HPC & Free & USC student allocation & Use for training \& compute \\
CMC HPC & Free & CMC student allocation & Use for hosting (vLLM) only \\
Modal GPU inference & $\sim$\$0.0003/call & \$30/mo free credit & $\sim$100k calls free \\
Vercel frontend & Free & Hobby tier & No action needed \\
GitHub & Free & Student Developer Pack & Apply at education.github.com \\
Namecheap domain & Free & via GitHub Student Pack & Claim via Student Pack \\
Gemini Pro & Free & Google AI Studio API & 60 RPM, 1k req/day \\
Bluesky AT Protocol & Free & Open source & Primary social; full firehose \\
Apify X Scraper & Free & Free tier & Secondary live; 500 results/run \\
Meta Content Library & Free & Academic research & Apply Week 0 \\
Claude Pro Max & Already paid & Existing subscription & Reserve for hard problems \\
\bottomrule
\multicolumn{4}{l}{\small Expected total cost during development: \textbf{\$0--\$5/month.}} \\
\end{tabularx}
\end{center}

Set a \$10/month hard cap on the Modal API at \url{modal.com/settings/billing}
to prevent unexpected charges if demo traffic spikes.
Do not use \texttt{keep\_warm} on the Modal container during development ---
cold starts cost nothing and the \$30 free credit is sufficient for all
testing. Reserve paid Modal time for the live demo day only.

\end{document}